% Citizenship and human rights — Anglais Tle Tech — Cours signé OnBuch+
% Programme officiel MINESEC. Compiler avec : tectonic N04-citizenship-and-human-rights.tex
\newcommand{\DOCMATIERE}{Anglais}
\newcommand{\DOCNIVEAU}{Tle Tech}
\newcommand{\DOCMODULE}{Anglais – classe de Terminale technique}
\newcommand{\DOCLECON}{N04}
\newcommand{\DOCTITRE}{Citizenship and human rights}
% OnBuch+ — préambule « cours signés » ENRICHI (compilé avec tectonic / XeLaTeX)
% Système visuel « Pop » d'OnBuch+ : fond crème (jamais blanc), bordures encre
% épaisses, ombres « dures » décalées, titres Archivo Black en capitales.
% Version MATHÉMATIQUES : graphes (pgfplots/tikz), boîtes pédagogiques
% (définition, propriété, méthode, attention, le savais-tu, exercice résolu,
% exercice, TP, bilan), page de garde et signature OnBuch+ ancrée en bas.
%
% Macros attendues dans main.tex AVANT \input{preamble} :
%   \DOCMATIERE  \DOCNIVEAU  \DOCMODULE  \DOCLECON (numéro)  \DOCTITRE
\documentclass[11pt]{article}

\usepackage{fontspec}
\usepackage[french]{babel}
\usepackage[a4paper,margin=2cm,top=3.4cm,bottom=2.8cm,headheight=1.4cm,footskip=1.3cm]{geometry}
\usepackage{amsmath,amssymb,amsfonts}
\usepackage{mathtools} % \xmapsto, \coloneqq... souvent utilisés par les modèles
\usepackage{cancel} % simplifications barrées (bilans de forces)
% \abs{z} (module, valeur absolue) et \norm{u} (norme), utilisés par les modèles
\providecommand{\abs}{}\let\abs\relax\DeclarePairedDelimiter{\abs}{\lvert}{\rvert}
\providecommand{\norm}{}\let\norm\relax\DeclarePairedDelimiter{\norm}{\lVert}{\rVert}
\usepackage[table]{xcolor} % [table] : fournit aussi \rowcolors (alternance de lignes)
\usepackage{pagecolor}
\usepackage{tikz}
\usetikzlibrary{arrows.meta,positioning,calc,shapes.geometric,shapes.misc,shapes.arrows,decorations.pathmorphing,decorations.pathreplacing,decorations.markings,patterns,shadows,mindmap,trees,fit,backgrounds,matrix,intersections,angles,quotes}
\usepackage{pgfplots}
\usepackage[version=4]{mhchem} % équations nucléaires \ce{^{235}_{92}U}
\usepackage[european,siunitx]{circuitikz} % schémas électriques
\pgfplotsset{compat=1.18}
\usepgfplotslibrary{fillbetween,groupplots} % aires entre courbes (calcul intégral)
\usepackage{enumitem}
\usepackage{fancyhdr}
% [most] charge toutes les bibliothèques tcolorbox (listings, minted, raster,
% documentation...) et gonfle inutilement la mémoire TeX sur un document long
% (~300 boîtes) : on ne charge que ce qui est réellement utilisé.
\usepackage[skins,breakable]{tcolorbox}
\usepackage{graphicx}
\usepackage{colortbl}
\usepackage{booktabs}
\usepackage{tabularx}
\usepackage{multirow}
\usepackage{diagbox} % case d'en-tête diagonale des tableaux à double entrée
% Colonnes extensibles centrée / gauche / droite (C, L, R), souvent utilisées par les modèles
\newcolumntype{C}{>{\centering\arraybackslash}X}
\newcolumntype{L}{>{\raggedright\arraybackslash}X}
\newcolumntype{R}{>{\raggedleft\arraybackslash}X}
\usepackage{array}
\usepackage{multicol}
\usepackage{siunitx}
% parse-numbers=false : accepte aussi \SI{2,5 \times 10^{-3}}{...}
\sisetup{locale=FR,output-decimal-marker={,},group-minimum-digits=4,parse-numbers=false}
\DeclareSIUnit\ohm{\ensuremath{\symup{\Omega}}} % U+2126 absent de la police
\DeclareSIUnit\division{div} % réglages d'oscilloscope (ms/div, V/div)
\DeclareSIUnit\tr{tr} % tours (vitesse de rotation, ex. \SI{1500}{\tr\per\minute})
\DeclareSIUnit\molar{M} % concentration molaire (ex. \SI{3}{\molar}, \SI{1}{\micro\molar})
\DeclareSIUnit\an{an} % année (le modèle confond parfois avec \year, un compteur TeX) — ex. \SI{5}{\kilo\meter\per\an}
\DeclareSIUnit\FCFA{FCFA} % franc CFA (ex. \SI{500}{\FCFA\per\kilo\gram})
\DeclareSIUnit\degreeBrix{\textdegree Brix} % teneur en sucre d'un sirop/jus (réfractométrie)
\DeclareSIUnit\month{mois} % le modèle utilise parfois \month, un compteur TeX, comme unité de durée
\DeclareSIUnit\microliter{\micro\liter} % le modèle écrit parfois \microliter en un seul token
\DeclareSIUnit\million{millions} % dénombrement (ex. \SI{15}{\million\per\mL} spermatozoïdes)
\usepackage{qrcode}
% \degree : commande fréquemment utilisée par les modèles pour un angle ou une
% température en dehors de \SI{}{} (non fournie par les paquets chargés).
\providecommand{\degree}{^{\circ}}
% \qed (fin de démonstration), utilisé par les modèles sans amsthm
\providecommand{\qed}{\hfill$\square$}
% \barre : double barre des valeurs interdites dans un tableau de variations (tkz-tab)
\providecommand{\barre}{\ensuremath{\|}}
% environnement proof (amsthm non chargé) utilisé par les modèles
\ifdefined\proof\else\newenvironment{proof}[1][Démonstration]{\par\noindent\textit{#1.}\ }{\qed\par}\fi
\usepackage{lastpage}
\usepackage{hyperref}
\hypersetup{hidelinks}

% ============================================================
% Palette « Pop » OnBuch+ — mêmes hex que la classe P (plus_ui.dart)
% ============================================================
\definecolor{poporange}{HTML}{F59321}
\definecolor{popdark}{HTML}{E07A0C}
\definecolor{popcream}{HTML}{FFF6E8}
\definecolor{popink}{HTML}{1C1714}
\definecolor{poppurple}{HTML}{7A5AE0}
\definecolor{popgreen}{HTML}{2E9E6A}
\definecolor{poppink}{HTML}{E0457B}
\definecolor{popblue}{HTML}{2F7DE1}
\definecolor{popgold}{HTML}{C98A12}
\definecolor{popmuted}{HTML}{8A7F76}
\definecolor{popline}{HTML}{EADFD2}
\definecolor{popgray}{HTML}{8A7F76}
\colorlet{popgrey}{popgray}
% Alias tolérants (le modèle nomme parfois les couleurs différemment)
\colorlet{popwhite}{white}
\colorlet{popblack}{popink}
\colorlet{popred}{poppink}
\colorlet{popyellow}{popgold}
% Teintes claires (fonds de tableaux/encadrés) — le modèle nomme parfois
% les couleurs avec un suffixe L (ex. popblueL) sans les avoir définies.
\colorlet{poporangeL}{poporange!20}
\colorlet{popdarkL}{popdark!20}
\colorlet{poppurpleL}{poppurple!20}
\colorlet{popgreenL}{popgreen!20}
\colorlet{poppinkL}{poppink!20}
\colorlet{popblueL}{popblue!20}
\colorlet{popgoldL}{popgold!20}
\colorlet{popredL}{popred!20}
\colorlet{popgrayL}{popgray!20}
% Teintes claires pour fonds de tableaux / zones
\colorlet{popblueL}{popblue!10!white}
\colorlet{poporangeL}{poporange!14!white}
\colorlet{popgreenL}{popgreen!12!white}
\colorlet{poppurpleL}{poppurple!12!white}
\colorlet{poppinkL}{poppink!12!white}
\colorlet{popgoldL}{popgold!14!white}

\pagecolor{popcream}

% ============================================================
% Typographie
% ============================================================
\defaultfontfeatures{Ligatures=TeX}
\setmainfont{PlusJakartaSans-Medium.ttf}[Path=../../../fonts/,BoldFont=PlusJakartaSans-Bold.ttf]
% Maths en sans-serif (Fira Math) avec chiffres et lettres droites de Plus
% Jakarta, pour que les formules se fondent dans le texte.
\usepackage[math-style=ISO,bold-style=ISO]{unicode-math}
\setmathfont{FiraMath-Regular.otf}[Path=../../../fonts/]
\setmathfont{PlusJakartaSans-Medium.ttf}[Path=../../../fonts/,range={up/num,up/{Latin,latin},\mathup}]
\usepackage{icomma} % virgule décimale sans espace en mode math
% Caractères mathématiques Unicode absents des polices de texte (Plus Jakarta,
% Archivo Black) : rendus par leur équivalent mathématique, en texte comme en
% formule — sinon ils s'affichent en carré vide (ex. « Arithmétique dans ℤ »).
\usepackage{newunicodechar}
\newunicodechar{ℝ}{\ensuremath{\mathbb{R}}}
\newunicodechar{ℤ}{\ensuremath{\mathbb{Z}}}
\newunicodechar{ℂ}{\ensuremath{\mathbb{C}}}
\newunicodechar{ℕ}{\ensuremath{\mathbb{N}}}
\newunicodechar{ℚ}{\ensuremath{\mathbb{Q}}}
\newunicodechar{𝔻}{\ensuremath{\mathbb{D}}}
\newunicodechar{α}{\ensuremath{\alpha}}
\newunicodechar{β}{\ensuremath{\beta}}
\newunicodechar{γ}{\ensuremath{\gamma}}
\newunicodechar{δ}{\ensuremath{\delta}}
\newunicodechar{ε}{\ensuremath{\varepsilon}}
\newunicodechar{θ}{\ensuremath{\theta}}
\newunicodechar{λ}{\ensuremath{\lambda}}
\newunicodechar{μ}{\ensuremath{\mu}}
\newunicodechar{σ}{\ensuremath{\sigma}}
\newunicodechar{τ}{\ensuremath{\tau}}
\newunicodechar{φ}{\ensuremath{\varphi}}
\newunicodechar{ψ}{\ensuremath{\psi}}
\newunicodechar{ω}{\ensuremath{\omega}}
\newunicodechar{Σ}{\ensuremath{\Sigma}}
% majuscules produites par \MakeUppercase dans les titres (α-aminés -> Α)
\newunicodechar{Α}{\ensuremath{\alpha}}
\newunicodechar{Β}{\ensuremath{\beta}}
\newunicodechar{Γ}{\ensuremath{\Gamma}}
\newunicodechar{Δ}{\ensuremath{\Delta}}
\newunicodechar{Ε}{\ensuremath{\varepsilon}}
\newunicodechar{Θ}{\ensuremath{\Theta}}
\newunicodechar{Λ}{\ensuremath{\Lambda}}
\newunicodechar{Μ}{\ensuremath{\mu}}
\newunicodechar{Τ}{\ensuremath{\tau}}
\newunicodechar{Φ}{\ensuremath{\Phi}}
\newunicodechar{Ψ}{\ensuremath{\Psi}}
\newunicodechar{Ω}{\ensuremath{\Omega}}
\newunicodechar{↦}{\ensuremath{\mapsto}}
\newunicodechar{∈}{\ensuremath{\in}}
\newunicodechar{∉}{\ensuremath{\notin}}
\newunicodechar{∧}{\ensuremath{\wedge}}
\newunicodechar{⇔}{\ensuremath{\Leftrightarrow}}
\newunicodechar{⟺}{\ensuremath{\Longleftrightarrow}}
\newunicodechar{⇒}{\ensuremath{\Rightarrow}}
\newunicodechar{⟹}{\ensuremath{\Longrightarrow}}
\newunicodechar{∘}{\ensuremath{\circ}}
\newunicodechar{∪}{\ensuremath{\cup}}
\newunicodechar{∩}{\ensuremath{\cap}}
\newunicodechar{≡}{\ensuremath{\equiv}}
\newunicodechar{⊂}{\ensuremath{\subset}}
\newunicodechar{⊥}{\ensuremath{\perp}}
\newunicodechar{∃}{\ensuremath{\exists}}
\newunicodechar{∀}{\ensuremath{\forall}}
\newunicodechar{∛}{\ensuremath{\sqrt[3]{\,}}}
\newunicodechar{∜}{\ensuremath{\sqrt[4]{\,}}}
\newunicodechar{⃗}{\ensuremath{{}^{\to}}}
\newunicodechar{ⁿ}{\ifmmode^{n}\else\textsuperscript{\ensuremath{n}}\fi}
\newunicodechar{ˣ}{\ifmmode^{x}\else\textsuperscript{\ensuremath{x}}\fi}
\newunicodechar{ᵇ}{\ifmmode^{b}\else\textsuperscript{\ensuremath{b}}\fi}
\newunicodechar{ʳ}{\ifmmode^{r}\else\textsuperscript{\ensuremath{r}}\fi}
\newunicodechar{ᵘ}{\ifmmode^{u}\else\textsuperscript{\ensuremath{u}}\fi}
\newunicodechar{ʸ}{\ifmmode^{y}\else\textsuperscript{\ensuremath{y}}\fi}
\newunicodechar{ᵃ}{\ifmmode^{a}\else\textsuperscript{\ensuremath{a}}\fi}
\newunicodechar{ᵉ}{\ifmmode^{e}\else\textsuperscript{\ensuremath{e}}\fi}
\newunicodechar{ᵏ}{\ifmmode^{k}\else\textsuperscript{\ensuremath{k}}\fi}
\newunicodechar{ᵖ}{\ifmmode^{p}\else\textsuperscript{\ensuremath{p}}\fi}
\newunicodechar{ᴺ}{\ifmmode^{N}\else\textsuperscript{\ensuremath{N}}\fi}
\newunicodechar{ᶜ}{\ifmmode^{c}\else\textsuperscript{\ensuremath{c}}\fi}
\newunicodechar{ʰ}{\ifmmode^{h}\else\textsuperscript{\ensuremath{h}}\fi}
\newunicodechar{ᵠ}{\ifmmode^{q}\else\textsuperscript{\ensuremath{q}}\fi}
\newunicodechar{ₙ}{\ifmmode_{n}\else\textsubscript{\ensuremath{n}}\fi}
\newunicodechar{ₐ}{\ifmmode_{a}\else\textsubscript{\ensuremath{a}}\fi}
\newunicodechar{ᵢ}{\ifmmode_{i}\else\textsubscript{\ensuremath{i}}\fi}
\newunicodechar{ᵣ}{\ifmmode_{r}\else\textsubscript{\ensuremath{r}}\fi}
\newunicodechar{ₖ}{\ifmmode_{k}\else\textsubscript{\ensuremath{k}}\fi}
\newunicodechar{ᵦ}{\ifmmode_{\beta}\else\textsubscript{\ensuremath{\beta}}\fi}
\newunicodechar{ₓ}{\ifmmode_{x}\else\textsubscript{\ensuremath{x}}\fi}
\newfontfamily\popdisplay{ArchivoBlack-Regular.ttf}[Path=../../../fonts/]
\newfontfamily\poplogo{SpaceGrotesk-Bold.ttf}[Path=../../../fonts/]
\newfontfamily\popsemi{PlusJakartaSans-SemiBold.ttf}[Path=../../../fonts/]
\newfontfamily\popxbold{PlusJakartaSans-ExtraBold.ttf}[Path=../../../fonts/]
\newfontfamily\popxboldls{PlusJakartaSans-ExtraBold.ttf}[Path=../../../fonts/,LetterSpace=3.0]

\setlist[enumerate]{leftmargin=1.7em,itemsep=5pt,topsep=4pt}
\setlist[itemize]{leftmargin=1.7em,itemsep=4pt,topsep=4pt,label={\color{poporange}\textbullet}}
\setlength{\parindent}{0pt}
\setlength{\parskip}{5pt}
\renewcommand{\arraystretch}{1.25}

% pgfplots : style Pop par défaut
\pgfplotsset{
  every axis/.append style={axis line style={popink,line width=1pt},tick style={popink},
    grid style={popline,line width=0.5pt},label style={font=\small},tick label style={font=\footnotesize}},
  popcurve/.style={poporange,line width=1.8pt,no markers,smooth},
}
% Même style disponible dans un \draw TikZ simple (hors pgfplots)
\tikzset{popcurve/.style={poporange,line width=1.8pt}}
\tikzset{popgrid/.style={popline,very thin}}
\colorlet{popcurve}{poporange} % le nom du style est parfois utilisé comme couleur

% ============================================================
% Pill / squiggle
% ============================================================
\newcommand{\poppill}[3]{%
  \tikz[baseline=-0.55ex]\node[draw=popink,line width=1.2pt,rounded corners=2.6pt,%
    fill=#2,inner xsep=6.5pt,inner ysep=3pt]{%
    \popxboldls\fontsize{7}{7}\selectfont\color{#3}\MakeUppercase{#1}};%
}
\newcommand{\popsquiggle}[1]{%
  \tikz[baseline=-0.4ex]{
    \draw[#1,line width=2.6pt,line cap=round]
      plot [smooth] coordinates {(0,0.09) (0.34,0.01) (0.68,0.21) (1.02,0.01) (1.36,0.10)};
  }%
}
% Mot-clé surligné (marqueur orange sous le texte)
\newcommand{\cle}[1]{{\popxbold\color{popdark}#1}}

% ============================================================
% En-tête / pied
% ============================================================
\pagestyle{fancy}
\fancyhf{}
\renewcommand{\headrulewidth}{0pt}
\renewcommand{\footrulewidth}{0pt}
\lhead{\raisebox{-0.15ex}{\poplogo\fontsize{13}{13}\selectfont\color{popink}OnBuch\color{poporange}+}}
\rhead{\poppill{\DOCMATIERE\ \textbullet\ \DOCNIVEAU\ \textbullet\ Leçon \DOCLECON}{white}{popink}}
\lfoot{\popxbold\fontsize{7.6}{7.6}\selectfont\color{popmuted}\MakeUppercase{Cours signé OnBuch+}\ \textbullet\ {\popsemi\DOCTITRE}}
\rfoot{\tikz[baseline=-0.6ex]\node[rounded corners=8.5pt,fill=popink,minimum height=17pt,minimum width=17pt,inner xsep=5pt,inner sep=0]{\popxbold\fontsize{8}{8}\selectfont\color{popcream}\thepage\,/\,\pageref*{LastPage}};}

% ============================================================
% Titres
% ============================================================
\newcommand{\chaptitre}[2]{%
  \begin{tcolorbox}[enhanced,colback=poporange,colframe=popink,boxrule=2.6pt,arc=11pt,
      drop shadow=popink,left=15pt,right=15pt,top=12pt,bottom=11pt,before skip=3pt,after skip=18pt]
    \poppill{#2}{popink}{popcream}\par\vspace{9pt}
    {\raggedright\hyphenpenalty=10000\exhyphenpenalty=10000\popdisplay\fontsize{22}{24}\selectfont\color{popcream}\MakeUppercase{#1}\par}
  \end{tcolorbox}
}
% Section numérotée : pastille orange avec le numéro + titre Archivo
\newcounter{popsec}
\newcommand{\coursec}[1]{%
  \stepcounter{popsec}\setcounter{popsub}{0}%
  \par\vspace{16pt}\noindent
  \begin{minipage}{\linewidth}
  \tikz[baseline=-0.7ex]\node[draw=popink,line width=1.6pt,fill=poporange,rounded corners=4pt,minimum size=22pt,inner sep=0]{\popdisplay\fontsize{12}{12}\selectfont\color{popcream}\thepopsec};%
  \hspace{8pt}{\popdisplay\fontsize{15}{17}\selectfont\color{popink}\MakeUppercase{#1}}\par
  \vspace{4pt}\noindent{\color{poporange}\rule{36pt}{3.6pt}}
  \end{minipage}\par\nopagebreak\vspace{8pt}
}
\newcounter{popsub}
\newcommand{\courssub}[1]{%
  \stepcounter{popsub}%
  \par\vspace{9pt}\noindent{\popxbold\fontsize{11.5}{13}\selectfont\color{popink}{\color{poporange}\thepopsec.\thepopsub}\hspace{6pt}#1}\par\nopagebreak\vspace{4pt}
}

% ============================================================
% Boîtes pédagogiques — même recette : fond blanc, bordure épaisse colorée,
% ombre dure encre, badge plein. Titre optionnel [#1].
% ============================================================
\tcbset{popbase/.style={breakable,enhanced,colback=white,boxrule=2.3pt,arc=9pt,
  shadow={3pt}{-3pt}{0pt}{popink},left=14pt,right=14pt,top=11pt,bottom=11pt,before skip=11pt,after skip=15pt,
  fontupper=\popsemi\fontsize{10.6}{14.5}\selectfont\color{popink},
  fontlower=\popsemi\fontsize{10.6}{14.5}\selectfont\color{popink}}}
% Coche dessinée (le glyphe ✓ n'existe pas dans les polices du document)
\newcommand{\popcheck}[1]{\tikz[baseline=-0.5ex]\draw[#1,line width=1.6pt,line cap=round,line join=round](-0.13,0.0)--(-0.04,-0.09)--(0.13,0.1);}
\providecommand{\coche}{\popcheck{popgreen}}
\providecommand{\resultat}[1]{\par\smallskip\noindent\fcolorbox{popgreen}{popgreenL}{\parbox{\dimexpr\linewidth-2\fboxsep-2\fboxrule\relax}{\textbf{Résultat.} #1}}\par\smallskip}
\newcommand{\popboxhead}[3]{\poppill{#1}{#2}{white}\ifx\relax#3\relax\else\hspace{6pt}{\popxbold\fontsize{10.6}{12}\selectfont\color{popink}#3}\fi\par\vspace{7pt}}

\newtcolorbox{aretenir}[1][]{popbase,colframe=popgreen,before upper={\popboxhead{À retenir}{popgreen}{#1}}}
\newtcolorbox{exemplebox}[1][]{popbase,colframe=poporange,before upper={\popboxhead{Exemple}{poporange}{#1}}}
\newtcolorbox{definition}[1][]{popbase,colframe=popblue,before upper={\popboxhead{Définition}{popblue}{#1}}}
\newtcolorbox{propriete}[1][]{popbase,colframe=poppurple,before upper={\popboxhead{Propriété}{poppurple}{#1}}}
\newtcolorbox{methode}[1][]{popbase,colframe=popgold,before upper={\popboxhead{Méthode}{popgold}{#1}}}
\newtcolorbox{attention}[1][]{popbase,colframe=poppink,before upper={\popboxhead{Attention}{poppink}{#1}}}
\newtcolorbox{experience}[1][]{popbase,colframe=popblue,colback=popblueL,before upper={\popboxhead{Expérience}{popblue}{#1}}}
% « Le savais-tu ? » : carte violette pleine (contexte, culture, Cameroun)
\newtcolorbox{savaistu}[1][]{popbase,colback=poppurple,colframe=popink,
  fontupper=\popsemi\fontsize{10.4}{14.2}\selectfont\color{white},
  before upper={\poppill{Le savais-tu\,?}{white}{poppurple}\ifx\relax#1\relax\else\hspace{6pt}{\popxbold\color{white}#1}\fi\par\vspace{7pt}}}
% Exercice résolu : énoncé en haut, \tcblower puis la solution sur fond crème
\newcounter{popexo}\newcounter{popexr}
\newtcolorbox{exoresolu}[1][]{popbase,colframe=poporange,colbacklower=popcream,
  segmentation style={popink,line width=1.2pt,dashed},
  before upper={\stepcounter{popexr}\popboxhead{Exercice résolu \thepopexr}{poporange}{#1}},
  before lower={\poppill{Solution}{popink}{popcream}\par\vspace{6pt}}}
% Exercice (non résolu en place) — la correction est dans la partie Corrigés
\newtcolorbox{exercice}[1][]{popbase,colframe=popink,
  before upper={\stepcounter{popexo}\popboxhead{Exercice \thepopexo}{popink}{#1}}}
% Corrigé
\newtcolorbox{corrige}[1][]{popbase,colframe=popgreen,colback=popgreenL,
  before upper={\popboxhead{Corrigé}{popgreen}{#1}}}
% Objectifs / prérequis / bilan
\newtcolorbox{objectifs}{popbase,colframe=popink,colback=poporangeL,
  before upper={\popboxhead{Objectifs de la leçon}{popink}{}}}
\newtcolorbox{prerequis}{popbase,colframe=popink,colback=popblueL,
  before upper={\popboxhead{Prérequis}{popblue}{}}}
\newtcolorbox{bilan}{popbase,colframe=popink,colback=white,boxrule=2.8pt,
  before upper={\popboxhead{Fiche bilan}{poporange}{L'essentiel à connaître pour l'examen}}}
% Figure encadrée avec légende
\newtcolorbox{popfigure}[1][]{enhanced,breakable=false,colback=white,colframe=popline,boxrule=1.4pt,arc=8pt,
  left=10pt,right=10pt,top=10pt,bottom=8pt,before skip=10pt,after skip=12pt,halign=center,
  lower separated=false,halign lower=center,
  fontlower=\popsemi\fontsize{9}{11.5}\selectfont\color{popmuted},
  #1}
\newcommand{\legende}[1]{\par\vspace{6pt}{\popsemi\fontsize{9}{11.5}\selectfont\color{popmuted}#1}}

% ============================================================
% Signature OnBuch+
% ============================================================
\newcommand{\poplogoinline}[1]{{\poplogo\fontsize{#1}{#1}\selectfont\color{popink}OnBuch\color{poporange}+}}
\newcommand{\signatureonbuch}{%
  \begin{tcolorbox}[enhanced,colback=popink,colframe=popink,boxrule=2.2pt,arc=10pt,
      drop shadow=poporange,left=14pt,right=14pt,top=10pt,bottom=10pt,before skip=0pt,after skip=0pt]
    \tikz[baseline=-0.7ex]\node[circle,fill=popgreen,draw=popcream,line width=1.3pt,minimum size=22pt,inner sep=0]{\popcheck{white}};%
    \hspace{9pt}\begin{minipage}[c]{0.62\linewidth}
      {\popxbold\fontsize{12}{13}\selectfont\color{popcream}Signé\ \textbullet\ {\poplogo OnBuch\color{poporange}+}}\par\vspace{2pt}
      {\raggedright\popsemi\fontsize{8}{10.5}\selectfont\color{popline}\MakeUppercase{Équipe pédagogique OnBuch+}\\\MakeUppercase{Conforme au programme officiel MINESEC}\par}
    \end{minipage}\hfill
    {\poplogo\fontsize{22}{22}\selectfont\color{popcream}OnBuch\color{poporange}+}
  \end{tcolorbox}%
}

% ============================================================
% Page de garde
% #1 titre · #2 matière · #3 niveau · #4 module · #5 n° leçon · #6 durée ·
% #7 description / objectifs (texte ou itemize)
% ============================================================
\newcommand{\pagegarde}[7]{%
  \thispagestyle{empty}
  \newgeometry{a4paper,margin=1.7cm,top=1.6cm,bottom=1.4cm}%
  \begin{tikzpicture}[remember picture,overlay]
    \fill[poporange!18] ([xshift=-3.2cm,yshift=-3.6cm]current page.north east) circle (4.6cm);
    \fill[poppurple!14] ([xshift=2.4cm,yshift=7.6cm]current page.south west) circle (3.8cm);
  \end{tikzpicture}%
  {\poplogo\fontsize{30}{30}\selectfont\color{popink}OnBuch\color{poporange}+}\hfill
  \raisebox{0.6ex}{\poppill{Cours signé \textbullet\ Premium}{popink}{popcream}}\par
  \vspace{1pt}\popsquiggle{poppurple}\par
  \vspace{8pt}
  \begin{tcolorbox}[enhanced,colback=poporange,colframe=popink,boxrule=2.8pt,arc=13pt,
      drop shadow=popink,left=18pt,right=18pt,top=12pt,bottom=13pt,after skip=10pt]
    \poppill{#2 \textbullet\ #3}{popink}{popcream}\hspace{5pt}\poppill{#4}{white}{popink}\par\vspace{8pt}
    {\popdisplay\fontsize{14}{14}\selectfont\color{popink}LEÇON #5}\par\vspace{4pt}
    {\raggedright\hyphenpenalty=10000\exhyphenpenalty=10000\popdisplay\fontsize{25}{27}\selectfont\color{popcream}\MakeUppercase{#1}\par}
    \vspace{8pt}
    {\popxbold\fontsize{9.5}{9.5}\selectfont\color{popink}\MakeUppercase{Durée conseillée : #6}}
  \end{tcolorbox}
  \begin{tcolorbox}[enhanced,colback=white,colframe=popink,boxrule=2.2pt,arc=10pt,
      drop shadow=popink,left=16pt,right=16pt,top=10pt,bottom=10pt,after skip=8pt]
    \poppill{Dans cette leçon}{poppurple}{white}\par\vspace{5pt}
    \popsemi\fontsize{9.3}{12.6}\selectfont\color{popink}#7
  \end{tcolorbox}
  \noindent\begin{minipage}[t]{0.487\linewidth}
    \begin{tcolorbox}[enhanced,colback=popgreen,colframe=popink,boxrule=2.2pt,arc=10pt,
        drop shadow=popink,left=11pt,right=11pt,top=10pt,bottom=10pt,height=3.1cm,valign=center]
      \tcbox[colback=white,colframe=popink,boxrule=1.3pt,arc=5pt,boxsep=0pt,nobeforeafter,
        left=3pt,right=3pt,top=3pt,bottom=3pt,valign=center]{\qrcode[height=1.9cm]{https://play.google.com/store/apps/details?id=com.luvvix.onbuch&hl=fr}}%
      \hspace{8pt}\begin{minipage}[c]{0.5\linewidth}
        {\popdisplay\fontsize{11}{12.5}\selectfont\color{white}\MakeUppercase{Télécharge OnBuch}}\par\vspace{4pt}
        {\raggedright\popsemi\fontsize{8.4}{11}\selectfont\color{white}Gratuit \textbullet\ Google Play\\Tuteur IA, annales, concours, résultats\par}
      \end{minipage}
    \end{tcolorbox}
  \end{minipage}\hfill
  \begin{minipage}[t]{0.487\linewidth}
    \begin{tcolorbox}[enhanced,colback=poppurple,colframe=popink,boxrule=2.2pt,arc=10pt,
        drop shadow=popink,left=11pt,right=11pt,top=10pt,bottom=10pt,height=3.1cm,valign=center]
      \tikz[baseline=-0.7ex]\node[circle,fill=white,draw=popink,line width=1.3pt,minimum size=1.6cm,inner sep=0]{\popdisplay\fontsize{24}{24}\selectfont\color{poppurple}+};%
      \hspace{8pt}\begin{minipage}[c]{0.6\linewidth}
        {\popdisplay\fontsize{11}{12.5}\selectfont\color{white}\MakeUppercase{Passe à OnBuch+}}\par\vspace{4pt}
        {\raggedright\popsemi\fontsize{8.4}{11}\selectfont\color{white}Tous les cours signés, Tuteur IA illimité, fiches PDF — onglet OnBuch+ de l'app\par}
      \end{minipage}
    \end{tcolorbox}
  \end{minipage}
  \vspace{8pt}
  \popcontenu
  \vfill
  \signatureonbuch
  \restoregeometry
  \newpage
}

% Rangée « Ce cours contient » (page de garde)
\newcommand{\poptile}[3]{% #1 couleur #2 gros texte #3 légende
  \begin{tcolorbox}[enhanced,colback=#1,colframe=popink,boxrule=2pt,arc=9pt,shadow={2.5pt}{-2.5pt}{0pt}{popink},
      left=6pt,right=6pt,top=9pt,bottom=9pt,halign=center,valign=center,height=1.75cm,width=\linewidth,nobeforeafter]
    {\popdisplay\fontsize{12.5}{13}\selectfont\color{white}\MakeUppercase{#2}}\par\vspace{3pt}
    {\centering\popsemi\fontsize{7.8}{9.5}\selectfont\color{white}#3\par}
  \end{tcolorbox}}
\newcommand{\popcontenu}{%
  \par\noindent{\popxboldls\fontsize{7.5}{7.5}\selectfont\color{popmuted}CE COURS SIGNÉ CONTIENT}\par\vspace{7pt}
  \noindent\begin{minipage}[t]{0.235\linewidth}\poptile{popblue}{Cours}{complet, illustré, pas à pas}\end{minipage}\hfill
  \begin{minipage}[t]{0.235\linewidth}\poptile{poporange}{Méthodes}{et exercices résolus}\end{minipage}\hfill
  \begin{minipage}[t]{0.235\linewidth}\poptile{poppink}{Exercices}{type examen + corrigés}\end{minipage}\hfill
  \begin{minipage}[t]{0.235\linewidth}\poptile{popgreen}{Fiche bilan}{l'essentiel à retenir}\end{minipage}\par}

% Sommaire
\newtcolorbox{sommaire}{enhanced,colback=white,colframe=popink,boxrule=2.2pt,arc=10pt,
  drop shadow=popink,left=18pt,right=18pt,top=13pt,bottom=13pt,before skip=6pt,after skip=20pt,
  before upper={\poppill{Sommaire}{poppurple}{white}\par\vspace{9pt}\popsemi\fontsize{10.6}{14.5}\selectfont\color{popink}}}

% Signature de fin de cours — poussée tout en bas de la dernière page
\newcommand{\findecours}{%
  \par\vfill\nopagebreak
  \begin{center}\popsquiggle{poporange}\end{center}\vspace{4pt}
  \signatureonbuch
}
\AtBeginDocument{\def\llbracket{\mathopen{[\mkern-3.5mu[}}\def\rrbracket{\mathclose{]\mkern-3.5mu]}}} % intervalles d'entiers (glyphe ⟦ absent de la police)
% Glyphes absents de Fira Math : redéfinis à partir de glyphes disponibles
\AtBeginDocument{%
  \def\setminus{\mathbin{\backslash}}\let\smallsetminus\setminus
  \def\vdots{\mathord{\vbox{\baselineskip3.2pt\lineskiplimit0pt\kern4pt\hbox{.}\hbox{.}\hbox{.}}}}%
  \def\ddots{\mathinner{\mkern1mu\raise6.5pt\hbox{.}\mkern2mu\raise3.6pt\hbox{.}\mkern2mu\raise0.7pt\hbox{.}\mkern1mu}}%
  \def\triangle{\mathord{\scalebox{0.9}{$\symup{\Delta}$}}}%
  \def\nabla{\mathord{\raisebox{\depth}{\scalebox{1}[-1]{$\symup{\Delta}$}}}}%
  \def\checkmark{\popcheck{popdark}}%
  \def\star{\mathbin{\ast}}\def\bigstar{\mathord{\ast}}%
  \def\diamond{\mathbin{\scalebox{0.6}{$\Diamond$}}}%
  \def\bigcup{\mathop{\mathchoice{\raisebox{-0.3em}{\scalebox{1.7}{$\cup$}}}{\scalebox{1.25}{$\cup$}}{\cup}{\cup}}\displaylimits}%
  \def\bigcap{\mathop{\mathchoice{\raisebox{-0.3em}{\scalebox{1.7}{$\cap$}}}{\scalebox{1.25}{$\cap$}}{\cap}{\cap}}\displaylimits}%
  \def\lesseqqgtr{\mathrel{\lessgtr}}\def\bigoplus{\mathop{\oplus}\displaylimits}%
}
\providecommand{\ssi}{si et seulement si} % abréviation fréquente
\AtBeginDocument{\providecommand{\wideparen}{\overparen}} % arc de cercle

%% FIGFIX-BEGIN v4 — correctifs globaux des figures (docs/onbuch-generation/tools/figfix)
\usepackage{adjustbox}\usepackage{etoolbox}
% toute figure TikZ/pgfplots plus large que la ligne est réduite à la largeur disponible
\BeforeBeginEnvironment{tikzpicture}{\begin{adjustbox}{max width=\linewidth}}
\AfterEndEnvironment{tikzpicture}{\end{adjustbox}}
% légendes pgfplots lisibles par-dessus les courbes
\pgfplotsset{every axis/.append style={legend style={fill=white,fill opacity=0.92,text opacity=1,draw=black!15,inner sep=3pt}}}
% étiquettes d'arêtes (syntaxe edge["..."]) avec fond blanc
\tikzset{every edge quotes/.append style={fill=white,inner sep=1.2pt}}
%% FIGFIX-END
\begin{document}

\pagegarde{Citizenship and human rights}{Anglais}{Tle Tech}{Anglais – classe de Terminale technique}{N04}{13 séances (fiche de progression harmonisée)}{Cette leçon permet aux élèves de Terminale technique d'échanger des informations sur les campagnes électorales scolaires, les sondages, les formations à la démocratie et le volontariat pour l'égalité des sexes. Elle mobilise les questions rapportées, l'ordre des mots, les phrasal verbs, la concordance des temps et les passifs impersonnels, aboutissant à la rédaction d'une composition sur le volontariat citoyen au Cameroun.
\vspace{4pt}
\begin{multicols}{2}\small
\begin{itemize}
\item À la fin de cette leçon, je sais utiliser le vocabulaire lié aux campagnes électorales et au vote des leaders scolaires.
\item À la fin de cette leçon, je sais rapporter correctement une question (reported questions) dans un dialogue.
\item À la fin de cette leçon, je sais comprendre un texte audio sur les formations à la démocratie et en dégager les idées essentielles.
\item À la fin de cette leçon, je sais respecter l'ordre des mots dans une phrase anglaise déclarative, interrogative et exclamative.
\item À la fin de cette leçon, je sais lire et exploiter un texte sur le volontariat pour la promotion de l'égalité des sexes.
\item À la fin de cette leçon, je sais employer les phrasal verbs dans des expressions simples de but.
\end{itemize}\end{multicols}}

\begin{sommaire}
\begin{multicols}{2}
\begin{enumerate}
\item Speaking \& Vocabulary: Campaigning and voting for school leaders
\item Grammar focus 1: Reported questions
\item Listening \& Vocabulary: Training courses on democracy
\item Grammar focus 2: Word order
\item Reading \& Vocabulary: Volunteering in gender equality promotion
\item Grammar focus 3 \& Writing: Phrasal verbs, sequence of tenses, impersonal passives
\item Activité d'intégration
\item Méthodes et exercices résolus
\item Exercices
\item Corrigés des exercices
\item Fiche bilan
\end{enumerate}
\end{multicols}
\end{sommaire}

\begin{objectifs}
\begin{itemize}
\item À la fin de cette leçon, je sais utiliser le vocabulaire lié aux campagnes électorales et au vote des leaders scolaires.
\item À la fin de cette leçon, je sais rapporter correctement une question (reported questions) dans un dialogue.
\item À la fin de cette leçon, je sais comprendre un texte audio sur les formations à la démocratie et en dégager les idées essentielles.
\item À la fin de cette leçon, je sais respecter l'ordre des mots dans une phrase anglaise déclarative, interrogative et exclamative.
\item À la fin de cette leçon, je sais lire et exploiter un texte sur le volontariat pour la promotion de l'égalité des sexes.
\item À la fin de cette leçon, je sais employer les phrasal verbs dans des expressions simples de but.
\item À la fin de cette leçon, je sais appliquer la concordance des temps et construire des passifs impersonnels (It is said that...).
\item À la fin de cette leçon, je sais rédiger une composition structurée sur le volontariat en faveur de l'égalité des sexes au Cameroun.
\item À la fin de cette leçon, je sais justifier ma démarche et mes réponses dans une situation concrète de vie citoyenne.
\end{itemize}
\end{objectifs}

\begin{prerequis}
\begin{itemize}
\item Le discours direct et le discours rapporté (reported statements) vus en classe de Première.
\item Les temps de base de l'anglais : present simple, past simple, present perfect, future.
\item Le vocabulaire élémentaire de la vie scolaire et de la citoyenneté (vote, leader, rights, duties).
\item La structure de base de la phrase anglaise : sujet + verbe + complément.
\item Les notions de droits humains vues en Éducation à la citoyenneté (Human Rights Education).
\end{itemize}
\end{prerequis}

\newpage

\coursec{Speaking \& Vocabulary: Campaigning and voting for school leaders}

Pendant les élections des délégués de classe au lycée technique de Nkongsamba, Amina, candidate au poste de \cle{class prefect}, doit présenter son programme en anglais devant toute la classe. Un camarade lui demande : \emph{``What did you promise the voters yesterday?''} et elle hésite à rapporter la question correctement. Plus tard, elle découvre une affiche annonçant un \cle{training course on democracy} organisé par \cle{Elections Cameroon (ELECAM)} et une campagne de volontariat pour la promotion de l'égalité des sexes dans son quartier. Comment parler, comprendre et écrire sur la citoyenneté et les droits humains en anglais correct ? C'est tout l'enjeu de cette leçon.

\courssub{Setting the Scene: School Elections in Cameroon}

Dans les lycées techniques du Cameroun, l'élection des délégués (\cle{class prefects}), du \cle{head boy} et de la \cle{head girl} est un moment fort de la vie citoyenne. C'est une simulation grandeur nature du processus démocratique : les élèves apprennent à choisir leurs représentants, à examiner des programmes (\cle{manifestos}), à voter à bulletin secret et à accepter les résultats. Maîtriser le vocabulaire et les expressions de ce domaine vous permet non seulement de réussir vos examens (Probatoire, Baccalauréat technique), mais aussi de participer activement à la vie de votre établissement et, plus tard, aux élections nationales organisées par \cle{ELECAM}.

\courssub{Key Vocabulary: Campaigning and Voting}

Voici les mots et expressions essentiels, classés par étape du processus électoral. Apprenez-les avec leur équivalent français et essayez de les utiliser dans des phrases complètes.

\begin{center}
\begin{tabularx}{\linewidth}{|l|X|X|}
\hline
\rowcolor{popblueL}
\textbf{English Term} & \textbf{Definition / Context} & \textbf{Glossaire français} \\
\hline
\cle{to campaign} & To work actively to get votes. & faire campagne \\
\cle{candidate} & A person standing for election. & candidat(e) \\
\cle{manifesto} & A written declaration of policies and promises. & programme électoral, profession de foi \\
\cle{slogan} & A short, striking phrase used in campaign. & slogan, devise \\
\cle{poster} & A large printed notice for publicity. & affiche \\
\cle{to canvass (for votes)} & To ask people directly for their support. & solliciter des voix, faire du porte-à-porte \\
\cle{electorate} & All the people entitled to vote. & corps électoral, électorat \\
\cle{polling day} & The day when voting takes place. & jour du scrutin, jour de vote \\
\cle{ballot paper} & The paper on which voters mark their choice. & bulletin de vote \\
\cle{ballot box} & The sealed box where votes are placed. & urne \\
\cle{to cast a vote} & To formally deposit one's ballot. & voter, déposer son bulletin \\
\cle{polling station} & The place where voting occurs. & bureau de vote \\
\cle{results} & The official outcome of the election. & résultats \\
\cle{to win / lose an election} & To be elected / not elected. & gagner / perdre une élection \\
\hline
\end{tabularx}
\end{center}

\begin{definition}[Vocabulaire clé : Poll, Survey, Manifesto]
\begin{itemize}
\item \textbf{A poll / A survey} = une enquête d'opinion. On interroge un échantillon (\cle{sample}) de personnes (\cle{respondents}) pour connaître leurs intentions de vote ou leurs opinions. Les résultats sont souvent donnés en pourcentages (\cle{results in percentage}).
\item \textbf{A manifesto} = le programme électoral d'un candidat. C'est un document écrit où il/elle explique ce qu'il/elle fera s'il/elle est élu(e) : \emph{``My manifesto promises clean water in the dormitories and better library hours.''}
\end{itemize}
\end{definition}

\begin{exemplebox}[Exemple chiffré : Sondage au GTHS Yaoundé]
Lors d'un sondage (\cle{opinion poll}) mené auprès de 200 élèves (\cle{sample size = 200}) au GTHS Yaoundé, 120 élèves (\cle{respondents}) ont déclaré qu'ils voteraient pour une fille au poste de \cle{head girl}.
\begin{itemize}
\item \textbf{Calcul :} $\frac{120}{200} \times 100 = 60\%$.
\item \textbf{Phrase type :} \emph{``According to the survey, 60\% of the respondents said they would vote for a female head girl.''}
\end{itemize}
\end{exemplebox}

\begin{savaistu}[ELECAM : Le garant des élections au Cameroun]
Au Cameroun, \textbf{ELECAM} (\emph{Elections Cameroon}) est l'organisme indépendant chargé d'organiser, de gérer et de superviser toutes les opérations électorales (présidentielles, législatives, municipales, mais aussi la formation à la citoyenneté). Lors des \cle{training courses on democracy}, des agents d'ELECAM expliquent le rôle de la carte d'électeur, le déroulement du scrutin et l'importance du vote secret. Participer à ces formations est un excellent atout pour votre dossier d'élève citoyen !
\end{savaistu}

\courssub{Vocabulary for Polls and Surveys (Sondages)}

Au-delà du vote lui-même, les élèves mènent souvent des \cle{surveys} ou \cle{polls} dans la classe pour tester la popularité des candidats avant le jour J.

\begin{center}
\begin{tabularx}{\linewidth}{|l|X|X|}
\hline
\rowcolor{popgreenL}
\textbf{English Term} & \textbf{Usage} & \textbf{Glossaire français} \\
\hline
\cle{poll} / \cle{survey} & General terms for asking opinions. & sondage, enquête \\
\cle{questionnaire} & The set of written questions. & questionnaire \\
\cle{respondent} & The person answering. & personne interrogée, répondant \\
\cle{sample} & The group surveyed (representative). & échantillon \\
\cle{opinion poll} & Survey on voting intentions. & sondage d'opinion \\
\cle{to take part in a survey} & To answer the questions. & participer à une enquête \\
\cle{results in percentage} & Data shown as \%. & résultats en pourcentage \\
\hline
\end{tabularx}
\end{center}

\courssub{Useful Expressions for Exchanging Information}

Pour \textbf{échanger des informations} (fonction \emph{Speaking} du programme), utilisez ces modèles de phrases. Notez la structure : \emph{Question word + Auxiliary + Subject + Verb}.

\begin{methode}[Présenter les résultats d'un sondage]
Pour annoncer des résultats de manière formelle, suivez cette structure en 3 étapes :
\begin{enumerate}
\item \textbf{Source :} \emph{According to the survey / poll / questionnaire...} (D'après l'enquête...)
\item \textbf{Chiffre :} \emph{...X\% of respondents / students / the sample...} (...X\% des répondants / élèves / de l'échantillon...)
\item \textbf{Contenu :} \emph{...said that / declared that / think that + clause.} (...ont dit que / pensent que + proposition.)
\end{enumerate}
\textbf{Exemple :} \emph{``According to our class poll, 45\% of students think that Candidate A has the best manifesto.''}
\end{methode}

\begin{center}
\begin{tabularx}{\linewidth}{|l|X|}
\hline
\rowcolor{poporangeL}
\textbf{Expression} & \textbf{Contexte d'utilisation} \\
\hline
\emph{Who are you voting for?} & Question directe à un camarade (closed question $\to$ rising intonation). \\
\emph{What does your manifesto promise?} & Question à un candidat (open question $\to$ falling intonation). \\
\emph{I'm voting for [Name] because...} & Justifier son choix. \\
\emph{According to the survey, 60\% of students...} & Annoncer un résultat de sondage (voir Méthode). \\
\emph{The polling station opens at 8 a.m.} & Donner une info pratique (jour J). \\
\emph{Please cast your vote secretly.} & Instruction pour le vote. \\
\emph{The results will be proclaimed at 2 p.m.} & Annoncer la proclamation. \\
\hline
\end{tabularx}
\end{center}

\courssub{Pronunciation \& Intonation: Rising vs Falling}

L'intonation change le sens communicatif de la question. C'est crucial pour l'oral (Probatoire/Baccalauréat technique).

\begin{propriete}[Règles d'intonation des questions]
\begin{itemize}
\item \textbf{Questions fermées (Yes/No Questions) :} Intonation \textbf{montante} (\nearrow) à la fin. \\ \emph{``Are you a candidate?''} $\nearrow$ \emph{``Will you vote for him?''} $\nearrow$
\item \textbf{Questions ouvertes (Wh- Questions) :} Intonation \textbf{descendante} (\searrow) à la fin. \\ \emph{``What is your manifesto?''} $\searrow$ \emph{``Who did you vote for?''} $\searrow$ \emph{``When is polling day?''} $\searrow$
\end{itemize}
\end{propriete}

\begin{attention}[Piège classique : \emph{Vote FOR}]
En anglais, le verbe \textbf{vote} est intransitif quand il est suivi d'une personne. On dit toujours : \textbf{to vote FOR somebody}.
\begin{itemize}
\item \checkmark \emph{``I will vote \textbf{for} Amina.''}
\item $\times$ \emph{``I will vote Amina.''} (Incorrect ! Cela signifie « je vote Amina » comme si c'était une loi).
\end{itemize}
De même : \emph{to campaign FOR someone}, \emph{to canvass FOR votes}.
\end{attention}

\courssub{Guided Role-Plays (Jeux de rôles guidés)}

Entraînez-vous à l'oral en binôme ou en groupe de trois. Changez les rôles à chaque fois.

\begin{exemplebox}[Role-play 1 : Interviewing a Candidate]
\textbf{Contexte :} Un journaliste du journal scolaire interviewe un(e) candidat(e) au poste de \cle{head boy/head girl}.
\begin{itemize}
\item \textbf{Journalist :} \emph{``Good morning. What is your name and what position are you running for?''}
\item \textbf{Candidate :} \emph{``I am Paul, and I am running for Head Boy.''}
\item \textbf{Journalist :} \emph{``What are the three main points of your manifesto?''}
\item \textbf{Candidate :} \emph{``First, improving the canteen menu. Second, creating a study room. Third, organizing sports tournaments.''}
\item \textbf{Journalist :} \emph{``How will you canvass for votes?''}
\item \textbf{Candidate :} \emph{``I will put up posters and talk to students during breaks.''}
\item \textbf{Journalist :} \emph{``Thank you. Good luck on polling day!''}
\end{itemize}
\end{exemplebox}

\begin{exoresolu}[Role-play 2 : Presenting a Manifesto (Simulation orale)]
\textbf{Consigne :} Vous êtes candidat(e). Présentez votre manifesto en 1 minute devant la classe. Utilisez au moins 5 mots du vocabulaire ci-dessus.

\tcblower

\textbf{Proposition de réponse (Script) :}
\emph{``Dear friends, dear electorate, I am Grace, your candidate for Class Prefect. My \textbf{manifesto} is built on three pillars. First, \textbf{campaign} for a cleaner workshop environment. Second, set up a peer-tutoring \textbf{campaign} for difficult subjects like Mathematics. Third, ensure every student has a voice: I will organize a monthly \textbf{survey} to collect your ideas. On \textbf{polling day}, I ask you to \textbf{cast your vote} for experience and dedication. \textbf{Vote for} Grace! Thank you.''}

\textbf{Analyse :} Le texte utilise \cle{manifesto}, \cle{campaign} (nom et verbe), \cle{electorate}, \cle{polling day}, \cle{cast a vote}, \cle{vote for}, \cle{survey}. L'intonation est descendante sur les \emph{Wh-questions} implicites (\emph{``What is my manifesto?''}) et montante sur l'appel final (\emph{``Vote for Grace?''} $\nearrow$).
\end{exoresolu}

\begin{exemplebox}[Role-play 3 : Announcing Class Poll Results]
\textbf{Contexte :} Le délégué sortant annonce les résultats d'un sondage d'opinion (\cle{opinion poll}) réalisé la veille auprès de 40 élèves (\cle{sample}).
\begin{itemize}
\item \textbf{Speaker :} \emph{``Good morning everyone. \textbf{According to the survey} we conducted yesterday, \textbf{45\% of respondents} said they will vote for Candidate A. \textbf{35\%} support Candidate B, and \textbf{20\%} are still undecided. The \textbf{polling station} will open tomorrow at 8 a.m. Please bring your student ID.''}
\end{itemize}
\end{exemplebox}

\courssub{Visualizing the Electoral Process}

Comprendre la chronologie aide à structurer votre discours. Voici le déroulement standard d'un scrutin scolaire.

\begin{popfigure}
\begin{tikzpicture}[
    node distance=0.5cm and 2.5cm,
    stage/.style={rectangle, rounded corners, draw=popblue, thick, fill=popblueL, minimum width=2.8cm, minimum height=1.2cm, align=center, font=\small},
    arrow/.style={->, >=Stealth, thick, popdark},
    label/.style={font=\tiny, text=popmuted, align=center}
]
% Nodes
\node[stage, align=center] (c1){\textbf{1. CAMPAIGNING}\\\emph{(2 weeks before)}};
\node[stage, right=of c1, align=center] (c2){\textbf{2. POLLING DAY}\\\emph{(D-Day)}};
\node[stage, right=of c2, align=center] (c3){\textbf{3. COUNTING VOTES}\\\emph{(Immediately after)}};
\node[stage, right=of c3, align=center] (c4){\textbf{4. PROCLAMATION}\\\emph{(Results announced)}};
% Arrows
\draw[arrow] (c1.east) -- (c2.west) node[midway, above, label, align=center]{Posters, Speeches,\\ Canvassing, Manifestos};
\draw[arrow] (c2.east) -- (c3.west) node[midway, above, label, align=center]{Secret Ballot,\\ Ballot Box, ID Check};
\draw[arrow] (c3.east) -- (c4.west) node[midway, above, label, align=center]{Tally Sheets,\\ Observers, Transparency};
% Decorations
\node[above=0.3cm of c1, font=\bfseries\small, popdark] {PREPARATION};
\node[above=0.3cm of c2, font=\bfseries\small, popdark] {ACTION};
\node[above=0.3cm of c3, font=\bfseries\small, popdark] {VERIFICATION};
\node[above=0.3cm of c4, font=\bfseries\small, popdark] {OUTCOME};
\end{tikzpicture}
\legende{Figure 1 : Frise chronologique d'une élection scolaire type (Campaigning $\to$ Polling Day $\to$ Counting $\to$ Proclamation).}
\end{popfigure}

Maintenant, observons les résultats fictifs du sondage mentionné dans le Role-play 3.

\begin{popfigure}
\begin{tikzpicture}
\begin{axis}[
    ybar,
    bar width=1.2cm,
    width=\linewidth,
    height=8cm,
    enlargelimits=0.25,
    ylabel={Percentage (\%)},
    xlabel={Candidates},
    symbolic x coords={Candidate A, Candidate B, Candidate C, Undecided},
    xtick=data,
    nodes near coords,
    ymin=0, ymax=55,
    ymajorgrids=true,
    grid style={popline},
    tick label style={font=\small},
    label style={font=\small},
    axis lines=left,
    axis line style={popdark},
    tick style={popdark},
    bar shift=0pt,
]
\addplot[fill=popblue, draw=popblue!80!black] coordinates {(Candidate A,45) (Candidate B,35) (Candidate C,15) (Undecided,5)};
\legend{Intention de vote}
\end{axis}
\end{tikzpicture}
\legende{Figure 2 : Résultats d'un sondage de classe (échantillon de 40 élèves). Candidate A mène avec 45\%, mais les 5\% d'indécis peuvent changer l'issue.}
\end{popfigure}

\courssub{Practice Exercises (Exercices d'application)}

\begin{exercice}[Vocabulary \& Speaking Practice]
\textbf{Exercise 1 : Vocabulary Matching.} Match the English terms (1-8) with their French definitions (a-h).
\begin{enumerate}
\item Manifesto \hfill a. Jour du scrutin
\item Ballot box \hfill b. Programme électoral
\item Polling day \hfill c. Urne
\item To canvass \hfill d. Corps électoral
\item Electorate \hfill e. Solliciter des voix
\item Slogan \hfill f. Déposer son bulletin
\item To cast a vote \hfill g. Affiche
\item Poster \hfill h. Slogan, devise
\end{enumerate}

\textbf{Exercise 2 : Gap Fill.} Complete the sentences with the correct word from the box.
\emph{manifesto \quad polling station \quad ballot paper \quad campaign \quad respondents \quad percentage}
\begin{enumerate}
\item The \_\_\_\_\_\_\_\_\_\_ opens at 7 a.m. on Saturday.
\item Please fold your \_\_\_\_\_\_\_\_\_\_ before putting it in the box.
\item According to the survey, 60\% of \_\_\_\_\_\_\_\_\_\_ support the new library project.
\item She launched her \_\_\_\_\_\_\_\_\_\_ with a big meeting in the hall.
\item What is the \_\_\_\_\_\_\_\_\_\_ of students who voted?
\item His \_\_\_\_\_\_\_\_\_\_ promises free Wi-Fi for all students.
\end{enumerate}

\textbf{Exercise 3 : Intonation Practice.} Mark the intonation arrow ($\nearrow$ or $\searrow$) for each question. Then practise saying them aloud.
\begin{enumerate}
\item Are you ready to vote?
\item Who is your favourite candidate?
\item Will the results be announced today?
\item What time does the polling station close?
\item Did you read the manifesto?
\end{enumerate}

\textbf{Exercise 4 : Role-play Preparation (Pair Work).}
Student A: You are a candidate for \cle{Head Girl}. Prepare a 1-minute speech using \textbf{at least 6 vocabulary words} from the tables above.
Student B: You are a journalist. Prepare \textbf{4 questions} (mix of open and closed) to interview Student A.
Perform the interview in front of the class.
\end{exercice}

\begin{corrige}[Exercice 1 à 3]
\textbf{Exercise 1 :} 1-b, 2-c, 3-a, 4-e, 5-d, 6-h, 7-f, 8-g.

\textbf{Exercise 2 :}
\begin{enumerate}
\item polling station
\item ballot paper
\item respondents
\item campaign
\item percentage
\item manifesto
\end{enumerate}

\textbf{Exercise 3 :}
\begin{enumerate}
\item Are you ready to vote? $\nearrow$ (Closed/Yes-No)
\item Who is your favourite candidate? $\searrow$ (Open/Wh-)
\item Will the results be announced today? $\nearrow$ (Closed/Yes-No)
\item What time does the polling station close? $\searrow$ (Open/Wh-)
\item Did you read the manifesto? $\nearrow$ (Closed/Yes-No)
\end{enumerate}

\textbf{Exercise 4 :} (Production libre. Vérifiez l'utilisation correcte de \emph{vote FOR}, \emph{manifesto}, \emph{polling day}, \emph{ballot box}, \emph{canvass}, \emph{slogan}, \emph{electorate} et l'intonation des questions).
\end{corrige}

\begin{aretenir}[Résumé de la Section 1]
\begin{itemize}
\item \textbf{Vocabulaire clé :} \cle{campaign}, \cle{candidate}, \cle{manifesto}, \cle{slogan}, \cle{canvass}, \cle{electorate}, \cle{polling day}, \cle{ballot paper}, \cle{ballot box}, \cle{cast a vote}, \cle{polling station}, \cle{poll}, \cle{survey}, \cle{respondent}, \cle{sample}, \cle{percentage}.
\item \textbf{Structure pour résultats :} \emph{``According to the survey, [\%] of respondents said that...''}
\item \textbf{Grammaire orale :} \emph{Vote FOR someone} (pas \emph{vote someone}).
\item \textbf{Phonologie :} Questions fermées $\nearrow$ (montante) ; Questions ouvertes (Wh-) $\searrow$ (descendante).
\item \textbf{Contexte Cameroun :} ELECAM organise les élections et les formations à la démocratie.
\end{itemize}
\end{aretenir}

\begin{attention}[Transition vers la Section 2 : Reported Questions]
Remarquez la question du camarade d'Amina dans l'introduction : \emph{``What did you promise the voters yesterday?''}. Pour la rapporter plus tard, Amina devra utiliser le \textbf{discours rapporté} (\emph{reported speech}) : \emph{``He asked me what I \textbf{had promised} the voters \textbf{the day before}.''} La Section 2 (Grammar Focus 1) traitera en détail les règles de transformation des questions directes en questions indirectes (changements de temps, pronoms, adverbes de temps, ordre des mots).
\end{attention}

\coursec{Grammar focus 1: Reported questions}

In the previous section, you practised speaking about \cle{campaigning} (\emph{faire campagne}) and \cle{voting} (\emph{voter}) for school leaders. During a campaign, journalists interview candidates, students ask questions, and later everybody \emph{reports} what was said: ``The candidate asked me if I would vote for her.'' This is exactly the grammar point of this section: \cle{reported questions} (\emph{les questions rapportées / le discours indirect interrogatif}). Mastering it will allow you to report interviews, polls and surveys accurately --- a classic task at the Probatoire and Baccalauréat technique.

\courssub{What is reported speech? A quick reminder}

\begin{definition}[Reported speech]
\cle{Reported speech} (also called \emph{indirect speech}) is used to tell someone what another person said, without quoting the exact words. When we report, we usually change three things:
\begin{itemize}
\item the \cle{tenses} (they move one step back: this is the \cle{backshift});
\item the \cle{pronouns} and possessives (\emph{I} $\rightarrow$ \emph{he/she}, \emph{my} $\rightarrow$ \emph{his/her});
\item the \cle{time and place markers} (\emph{now} $\rightarrow$ \emph{then}, \emph{today} $\rightarrow$ \emph{that day}, \emph{here} $\rightarrow$ \emph{there}, \emph{tomorrow} $\rightarrow$ \emph{the next day}).
\end{itemize}
\end{definition}

\begin{center}
\begin{tabularx}{\linewidth}{|l|X|X|}
\hline
\rowcolor{popblueL} \textbf{Direct speech} & \textbf{Reported speech} & \textbf{Change} \\
\hline
present simple: \emph{I vote} & past simple: \emph{she voted} & backshift \\
\hline
present continuous: \emph{I am campaigning} & past continuous: \emph{she was campaigning} & backshift \\
\hline
past simple / present perfect: \emph{I voted / I have voted} & past perfect: \emph{she had voted} & backshift \\
\hline
will: \emph{I will win} & would: \emph{she would win} & modal backshift \\
\hline
can: \emph{I can lead} & could: \emph{she could lead} & modal backshift \\
\hline
\end{tabularx}
\end{center}

\begin{exemplebox}[From the campaign]
Direct: ``I am ready for the election,'' says Manka.\par
Reported: Manka says (that) she \textbf{is} ready for the election. (no backshift if the reporting verb is in the present)\par
Reported: Manka said (that) she \textbf{was} ready for the election. (backshift with a past reporting verb)
\end{exemplebox}

\courssub{Reporting yes/no questions: if and whether}

Imagine the scene: during the school election campaign, a journalist from the school radio asks a candidate: \emph{``Are you a candidate?''} Later, a student reports this question to a friend. He cannot keep the question form; he must transform it into a statement.

\begin{propriete}[Yes/no reported questions]
A direct question that expects the answer \emph{yes} or \emph{no} is reported with \cle{if} or \cle{whether} (\emph{si}), followed by the normal order \textbf{subject + verb}:
\begin{center}
\emph{``Are you a candidate?''} $\rightarrow$ He asked me \textbf{if/whether I was} a candidate.
\end{center}
There is \textbf{no inversion}, \textbf{no auxiliary do/does/did}, and \textbf{no question mark} in the reported question.
\end{propriete}

\begin{exemplebox}[Practice with campaign questions]
\begin{itemize}
\item ``Do you support the new candidate?'' $\rightarrow$ She asked me \textbf{if I supported} the new candidate.
\item ``Did you take part in the poll?'' \emph{(participer au sondage)} $\rightarrow$ The journalist asked \textbf{whether I had taken part} in the poll.
\item ``Will you vote tomorrow?'' $\rightarrow$ He asked me \textbf{if I would vote} the next day.
\end{itemize}
\end{exemplebox}

\begin{attention}[The disappearing question mark]
A very frequent mistake is to keep the question mark or the inversion in the reported question. Remember: a reported question is \emph{not} a question any more, it is a \textbf{declarative sentence} (\emph{une phrase déclarative}).\par
Wrong: \emph{He asked me if was I a candidate?}\par
Correct: \emph{He asked me if I was a candidate.}
\end{attention}

\courssub{Reporting wh-questions}

Wh-questions begin with an interrogative word: \cle{who}, \cle{what}, \cle{where}, \cle{when}, \cle{why}, \cle{how}, \cle{which}, \cle{whose}. In reported speech, we \textbf{keep the wh-word} as the link word, and we apply the same rules: subject + verb order, backshift, no question mark.

\begin{propriete}[Wh- reported questions]
\emph{``Who will you vote for?''} $\rightarrow$ She asked me \textbf{who I would vote for}.\par
Structure: reporting verb + (person) + wh-word + subject + verb (backshifted).
\end{propriete}

\begin{exemplebox}[More examples from the survey]
\begin{itemize}
\item ``Where do you \textbf{mark} your vote?'' \emph{(cocher/voter)} $\rightarrow$ He asked me \textbf{where I marked} my vote.
\item ``Why did you join the campaign team?'' $\rightarrow$ They asked her \textbf{why she had joined} the campaign team.
\item ``How many students will take part in the survey?'' $\rightarrow$ She wanted to know \textbf{how many students would take part} in the survey.
\end{itemize}
\end{exemplebox}

\begin{propriete}[No inversion, no do/does/did]
In a reported question, the auxiliary \emph{do/does/did} disappears completely and the subject comes before the verb: \emph{He asked where I lived} (and not \emph{where did I live}). The verb then carries the tense alone: \emph{did you live} $\rightarrow$ \emph{I had lived / I lived} depending on the backshift.
\end{propriete}

\courssub{Reporting verbs}

The most common reporting verb for questions is \cle{to ask}. But English offers several useful alternatives:

\begin{center}
\begin{tabularx}{\linewidth}{|l|X|}
\hline
\rowcolor{popblueL} \textbf{Reporting verb} & \textbf{Use and example} \\
\hline
to ask (someone) & general verb: \emph{She asked me if I was registered.} \\
\hline
to wonder & the speaker asks himself/herself: \emph{I wondered why he had refused to vote.} \\
\hline
to want to know & expresses curiosity: \emph{He wanted to know who the candidates were.} \\
\hline
to enquire & formal (\emph{se renseigner}): \emph{She enquired whether the poll was open.} \\
\hline
\end{tabularx}
\end{center}

Note: with \emph{to wonder}, we do not mention the person asked; with \emph{to ask}, we usually do (\emph{asked me / asked her}).

\courssub{Bonus for the Bac: reported orders and requests}

Questions are not the only sentences we report. Orders, commands and requests (\emph{les ordres et les demandes}) follow a different pattern: \textbf{tell / order / warn / advise / ask + object + to-infinitive} (or \emph{not to} + infinitive for the negative).

\begin{propriete}[Reported orders and requests]
\emph{``Vote for me!''} $\rightarrow$ He told me \textbf{to vote for him}.\par
\emph{``Don't cheat during the election.''} $\rightarrow$ She warned me \textbf{not to cheat} during the election.\par
Structure: reporting verb + object + (\textbf{not}) \textbf{to} + base verb.
\end{propriete}

\begin{exemplebox}[Campaign instructions]
\begin{itemize}
\item ``Please count the \textbf{ballot papers} (\emph{bulletins de vote}) carefully.'' $\rightarrow$ The president asked the \textbf{scrutineers} (\emph{scrutateurs}) \textbf{to count} the ballot papers carefully.
\item ``Don't shout in the \textbf{polling station} (\emph{bureau de vote}).'' $\rightarrow$ The supervisor ordered the students \textbf{not to shout} in the polling station.
\end{itemize}
\end{exemplebox}

\begin{popfigure}
\begin{center}
\begin{tikzpicture}[font=\small, >=Stealth]
% Example 1
\node[draw=popblue, fill=popblueL, rounded corners, text width=3.6cm, align=center] (d1) at (0,4) {``Are you a candidate?''};
\node[draw=poporange, fill=poporangeL, rounded corners, text width=2.6cm, align=center] (v1) at (5,4) {He asked me};
\node[draw=popgreen, fill=popgreenL, rounded corners, text width=4.2cm, align=center] (r1) at (10,4) {if I \textbf{was} a candidate.};
\draw[->, thick, popdark] (d1) -- (v1);
\draw[->, thick, popdark] (v1) -- (r1);
% Example 2
\node[draw=popblue, fill=popblueL, rounded corners, text width=3.6cm, align=center] (d2) at (0,2) {``Who will you vote for?''};
\node[draw=poporange, fill=poporangeL, rounded corners, text width=2.6cm, align=center] (v2) at (5,2) {She asked me};
\node[draw=popgreen, fill=popgreenL, rounded corners, text width=4.2cm, align=center] (r2) at (10,2) {who I \textbf{would vote} for.};
\draw[->, thick, popdark] (d2) -- (v2);
\draw[->, thick, popdark] (v2) -- (r2);
% Example 3
\node[draw=popblue, fill=popblueL, rounded corners, text width=3.6cm, align=center] (d3) at (0,0) {``Did you take part in the poll?''};
\node[draw=poporange, fill=poporangeL, rounded corners, text width=2.6cm, align=center] (v3) at (5,0) {The journalist asked};
\node[draw=popgreen, fill=popgreenL, rounded corners, text width=4.2cm, align=center] (r3) at (10,0) {whether I \textbf{had taken part} in the poll.};
\draw[->, thick, popdark] (d3) -- (v3);
\draw[->, thick, popdark] (v3) -- (r3);
% Column labels
\node[popblue, font=\bfseries] at (0,5.2) {Direct question};
\node[poporange, font=\bfseries] at (5,5.2) {Reporting verb};
\node[popgreen, font=\bfseries] at (10,5.2) {Reported question};
\end{tikzpicture}
\end{center}
\legende{Transformation table: a direct question (blue) is introduced by a reporting verb (orange) and becomes a reported question (green) with subject + verb order, backshift, and no question mark.}
\end{popfigure}

\courssub{Method: how to report a question in four steps}

\begin{methode}[Reporting a question]
\begin{enumerate}
\item \textbf{Identify the type of question}: does it expect \emph{yes/no}, or does it begin with a wh-word?
\item \textbf{Choose the link word}: use \emph{if/whether} for yes/no questions; \emph{keep the wh-word} for wh-questions.
\item \textbf{Apply the backshift}: move the tense one step back (present $\rightarrow$ past, past/perfect $\rightarrow$ past perfect, \emph{will} $\rightarrow$ \emph{would}), and adapt pronouns and time markers.
\item \textbf{Restore the statement order}: subject + verb, no \emph{do/does/did}, no inversion, and replace the question mark with a full stop.
\end{enumerate}
\end{methode}

\begin{exoresolu}[Reporting an interview with a school candidate]
Énoncé. During the school election campaign, a journalist interviewed the candidate Amina. Report her questions using the verbs in brackets.\par
1. ``Are you ready for the debate?'' (she asked Amina)\par
2. ``Why do you want to be head girl?'' \emph{(déléguée chef / présidente des élèves)} (she asked her)\par
3. ``How many students support your programme?'' (she wanted to know)\par
4. ``Will you organise a training course on democracy?'' (she enquired)\par
5. ``Don't forget to \textbf{register} (\emph{inscrire sur les listes électorales}) the voters.'' (she advised Amina)
\tcblower
Solution détaillée.\par
1. Yes/no question $\rightarrow$ \emph{if/whether}; present \emph{are} $\rightarrow$ past \emph{was}; order subject + verb: \textbf{She asked Amina if/whether she was ready for the debate.}\par
2. Wh-question $\rightarrow$ keep \emph{why}; present \emph{do you want} $\rightarrow$ past \emph{she wanted}; no auxiliary \emph{do}: \textbf{She asked her why she wanted to be head girl.}\par
3. Wh-question $\rightarrow$ keep \emph{how many}; present \emph{support} $\rightarrow$ past \emph{supported}: \textbf{She wanted to know how many students supported her programme.}\par
4. Yes/no question with \emph{will} $\rightarrow$ \emph{would}: \textbf{She enquired whether Amina would organise a training course on democracy.}\par
5. Negative order $\rightarrow$ \emph{advised + object + not to + verb}: \textbf{She advised Amina not to forget to register the voters.}
\end{exoresolu}

\begin{attention}[Frequent errors to avoid]
\begin{itemize}
\item Keeping the inversion: \emph{He asked me where did I vote.} $\rightarrow$ Correct: \emph{He asked me where I voted.}
\item Keeping \emph{do/does/did}: \emph{She asked if did I take part.} $\rightarrow$ Correct: \emph{She asked if I took part.}
\item Keeping the question mark: a reported question ends with a \textbf{full stop}.
\item Forgetting the backshift after a past reporting verb: \emph{He asked if I will vote} $\rightarrow$ \emph{He asked if I would vote}.
\end{itemize}
\end{attention}

\begin{savaistu}[Reported speech at the Bac technique]
At the Probatoire and Baccalauréat technique in Cameroon, the transformation from direct to indirect speech is a recurrent exercise in the \textbf{Grammar section}. You are often given three to five direct questions about civic life --- elections, polls, surveys, human rights campaigns --- and asked to report them. Examiners check four points: the correct link word (\emph{if/whether} or the wh-word), the backshift, the subject--verb order, and the final full stop. Each point earns marks, so apply the four-step method systematically!
\end{savaistu}

\begin{aretenir}[Key points]
\begin{itemize}
\item Yes/no questions are reported with \cle{if} or \cle{whether}; wh-questions keep their interrogative word.
\item A reported question has the order \textbf{subject + verb}: no inversion, no \emph{do/does/did}, no question mark.
\item Apply the \cle{backshift} after a past reporting verb, and adapt pronouns and time markers.
\item Common reporting verbs: \emph{ask, wonder, want to know, enquire}.
\item Orders and requests: \emph{tell/warn/advise/ask + object + (not) to + verb}.
\end{itemize}
\end{aretenir}

\begin{exercice}[Practice: report the poll]
Report these sentences from a school survey. Use the reporting verbs given.\par
1. ``Have you ever voted in a school election?'' (the interviewer asked me)\par
2. ``When will the results be published?'' (the students wondered)\par
3. ``Who is your favourite candidate?'' (she asked him)\par
4. ``Do you believe in gender equality?'' (the journalist enquired)\par
5. ``Don't \textbf{spoil} (\emph{annuler / invalider}) your \textbf{ballot paper} (\emph{bulletin de vote}).'' (the official warned the voters)
\end{exercice}

\begin{corrige}[Exercice 1]
1. The interviewer asked me if/whether I had ever voted in a school election. (present perfect $\rightarrow$ past perfect)\par
2. The students wondered when the results would be published. (\emph{will} $\rightarrow$ \emph{would}; statement order)\par
3. She asked him who his favourite candidate was. (\emph{your} $\rightarrow$ \emph{his}; no inversion)\par
4. The journalist enquired whether I believed in gender equality. (\emph{do you believe} $\rightarrow$ \emph{I believed}; \emph{do} disappears)\par
5. The official warned the voters not to spoil their ballot papers. (negative order: \emph{not to} + verb)
\end{corrige}

\coursec{Listening \& Vocabulary: Training courses on democracy}

In the previous section, we revised \cle{reported questions}: we learnt how to report what someone asked, for example \emph{``The facilitator asked us whether we had ever voted.''} Keep this tool in mind, because at the end of this section you will use reported questions and reported statements to reformulate what you hear. Now we move from speaking about elections to \cle{listening}: you will listen to a text about a training workshop on democracy organised in Douala, learn the vocabulary of training activities, and practise professional listening techniques.

\courssub{Why listen to texts about democracy training?}

In Cameroon, many organisations --- NGOs, ELECAM (Elections Cameroon), youth associations, churches --- organise \cle{training courses} on \cle{civic education}. As future technicians and citizens, you may attend such courses, or even facilitate them. Understanding spoken English in this context is a real-life skill tested at the Probatoire and the Baccalauréat technique.

\begin{definition}[Key definitions]
\cle{Civic education} is the training of citizens about their rights and duties (formation aux droits et devoirs du citoyen). \cle{Rule of law} (l'État de droit) means that everyone, including leaders, must obey the law, and that laws are applied fairly by independent courts.
\end{definition}

\courssub{Vocabulary: words and expressions for training activities}

Read the following table aloud with your teacher. The French glosses will help you memorise the words. Note the use of articles for countable nouns.

\begin{center}
\begin{tabularx}{\linewidth}{|l|X|}\hline
\rowcolor{popblueL} \textbf{English} & \textbf{Meaning / French gloss} \\ \hline
a training course & a series of lessons to learn a skill (une formation) \\ \hline
a workshop & a practical training session where participants work in groups (un atelier) \\ \hline
a seminar & a meeting for discussion on a special topic (un séminaire) \\ \hline
the facilitator & the person who leads the training (l'animateur / le formateur) \\ \hline
a participant & a person who attends the training (un participant) \\ \hline
a module & one unit or part of a course (un module) \\ \hline
a certificate & an official document given at the end of a course (une attestation, un certificat) \\ \hline
civic education & training about citizens' rights and duties (éducation civique) \\ \hline
democratic values & ideas like freedom, equality and justice (valeurs démocratiques) \\ \hline
good governance & managing public affairs honestly and transparently (bonne gouvernance) \\ \hline
tolerance & accepting people who are different (la tolérance) \\ \hline
rule of law & the principle that the law applies to everybody (l'État de droit) \\ \hline
to attend / to take part in & to be present at an activity (assister à / participer à) \\ \hline
to raise awareness & to inform people about a problem (sensibiliser) \\ \hline
\end{tabularx}
\end{center}

\begin{attention}[Common mistakes]
Do not confuse \cle{training} (formation pratique) and \cle{teaching} (enseignement scolaire). A technician follows \emph{a training course}; a pupil receives \emph{teaching}. Also, \emph{a training course} is followed by the preposition \cle{on} (or \cle{in}): say \emph{a course \textbf{on} democracy} or \emph{a course \textbf{in} democracy}, never \emph{a course of democracy}. In formal writing, prefer \emph{on} or \emph{in} over \emph{about}. Finally, \emph{training} is uncountable: no plural \emph{trainings}.
\end{attention}

\courssub{Grammar focus: Word order (SVOMPT)}

The programme requires revising \cle{word order} in English sentences. Unlike French, English has a very strict order. Mastering it is essential for both writing and speaking clearly during the exam.

\begin{propriete}[The SVOMPT rule]
A standard English declarative sentence follows the order: \textbf{S}ubject + \textbf{V}erb + \textbf{O}bject + \textbf{M}anner + \textbf{P}lace + \textbf{T}ime.
\begin{center}
\begin{tabularx}{\linewidth}{|l|l|X|}\hline
\rowcolor{popblueL} \textbf{Element} & \textbf{Question} & \textbf{Example} \\ \hline
S (Subject) & Who? & \emph{The facilitator} \\ \hline
V (Verb) & What action? & \emph{explained} \\ \hline
O (Object) & What/Whom? & \emph{the module} \\ \hline
M (Manner) & How? & \emph{clearly} \\ \hline
P (Place) & Where? & \emph{in the hall} \\ \hline
T (Time) & When? & \emph{yesterday} \\ \hline
\end{tabularx}
\end{center}
Full sentence: \emph{The facilitator explained the module clearly in the hall yesterday.}
\end{propriete}

\begin{attention}[Word order in reported questions]
Never keep the question form in a reported question. Say \emph{She asked us what we \textbf{had learnt}}, not \emph{She asked us what had we learnt}. After \emph{asked}, use the normal SVOMPT order: subject + verb.
\end{attention}

\courssub{The pillars of democracy}

Before listening, observe the diagram below. It presents the five \cle{pillars of democracy} (les piliers de la démocratie) that the facilitator in our listening text will mention. A democracy is like a building: if one pillar is weak, the whole building can fall.

\begin{popfigure}
\begin{center}
\begin{tikzpicture}[
  pillar/.style={rectangle, draw=popblue, fill=popblueL, rounded corners=2pt,
                 minimum width=2.8cm, minimum height=2.4cm, align=center, font=\small},
  roof/.style={rectangle, draw=popdark, fill=poporangeL, rounded corners=2pt,
               minimum width=16cm, minimum height=1cm, align=center, font=\bfseries\small},
  base/.style={rectangle, draw=popdark, fill=popgreenL, rounded corners=2pt,
               minimum width=16cm, minimum height=0.9cm, align=center, font=\small\itshape},
  pillarlabel/.style={font=\tiny, text=popdark, anchor=north, yshift=-0.1cm}
]
\node[roof] (roof) at (0,2.8) {DEMOCRACY};
\node[pillar, align=center] (p1) at (-6.4,0){Free and fair\\elections};
\node[pillar, align=center] (p2) at (-3.2,0){Rule of\\law};
\node[pillar, align=center] (p3) at (0,0){Human\\rights};
\node[pillar, align=center] (p4) at (3.2,0){Separation\\of powers\\(Executive, Legislative,\\Judiciary)};
\node[pillar, align=center] (p5) at (6.4,0){Free\\press};
\node[base] (base) at (0,-2.3) {Foundation: active and informed citizens};
\foreach \p in {p1,p2,p3,p4,p5} {
  \draw[thick, popdark] (\p.north) -- (\p.north |- roof.south);
  \draw[thick, popdark] (\p.south) -- (\p.south |- base.north);
}
\end{tikzpicture}
\end{center}
\legende{The five pillars of democracy: (1) free and fair elections, (2) the rule of law, (3) human rights, (4) the separation of powers between the executive, the legislative and the judiciary, (5) a free press. All of them rest on a foundation of active citizens --- that is why training courses on democracy matter.}
\end{popfigure}

\begin{savaistu}[Did you know?]
The Cameroonian Constitution of 1996 affirms in its preamble the attachment of the Cameroonian people to the fundamental rights of the human person, as expressed in the Universal Declaration of Human Rights. ELECAM (Elections Cameroon) is the body in charge of organising and supervising elections, and it regularly runs \emph{voter education} campaigns for young people.
\end{savaistu}

\courssub{Listening techniques: how to understand a spoken text}

Listening to English is not about understanding every single word. Even native speakers do not catch everything. Good listeners use strategies.

\begin{methode}[Before, during and after listening]
\begin{enumerate}
\item \textbf{Before listening:} read the questions or the task carefully and \cle{anticipate} the vocabulary. If the text is about a workshop, expect words like \emph{facilitator, module, certificate, participants}.
\item \textbf{First listening (global listening / gist):} listen for the general meaning only. Answer: Who is speaking? Where? What is the topic?
\item \textbf{Second listening (selective listening):} listen for precise details: dates, places, numbers of participants, names of organisations.
\item \textbf{Note-taking:} write \cle{key words}, not full sentences. Example: \emph{Douala --- 3 days --- 80 participants --- ELECAM --- certificates}.
\item \textbf{After listening:} use your notes to reformulate the content orally, with reported speech: \emph{The facilitator explained that... He asked the participants whether...}
\end{enumerate}
\end{methode}

\courssub{Guided listening: a workshop on citizenship in Douala}

Your teacher will read the following text twice (or play the recording). Close your books during the first listening; open them only for the note-taking grid during the second listening.

\begin{exemplebox}[Listening script (teacher's version)]
``Good morning, dear listeners. This is the report of our civic education workshop which took place at the Youth Centre in Bonanjo, Douala, from the 12th to the 14th of March. The training course was organised by the NGO \emph{Youth for Democracy}, in partnership with ELECAM. Eighty participants --- students, apprentices and young workers --- attended the three-day seminar. The facilitator, Mrs Etoundi, divided the course into four modules: democratic values, the rule of law, good governance, and tolerance in a multicultural society. On the second day, participants simulated an election campaign and voted for class delegates. At the end of the workshop, every participant received a certificate. Mrs Etoundi encouraged the young people to register on the electoral lists and to vote peacefully. She explained that democracy needs active and informed citizens. She asked the participants whether they were ready to share what they had learnt in their schools and neighbourhoods.''
\end{exemplebox}

\textbf{Note-taking grid (fill it in during the second listening):}

\begin{center}
\begin{tabularx}{\linewidth}{|l|X|}\hline
\rowcolor{popblueL} \textbf{Question} & \textbf{Your notes (key words only)} \\ \hline
Where did the workshop take place? & \\ \hline
When (dates)? & \\ \hline
Who organised it? & \\ \hline
How many participants? & \\ \hline
How many modules? Which ones? & \\ \hline
What did participants receive at the end? & \\ \hline
\end{tabularx}
\end{center}

\begin{exoresolu}[Comprehension questions on the listening text]
\textbf{Answer these questions in full sentences.}\par
1. Where and when did the workshop take place?\par
2. How many young people attended the seminar?\par
3. What were the four modules of the course?\par
4. What did Mrs Etoundi encourage the participants to do?
\tcblower
\textbf{Suggested answers:}\par
1. The workshop took place at the Youth Centre in Bonanjo, Douala, from the 12th to the 14th of March.\par
2. Eighty participants attended the seminar: students, apprentices and young workers.\par
3. The four modules were democratic values, the rule of law, good governance, and tolerance in a multicultural society.\par
4. She encouraged them to register on the electoral lists, to vote peacefully, and to share what they had learnt in their schools and neighbourhoods.
\end{exoresolu}

\begin{exemplebox}[A concrete figure]
In 2023, a workshop in Bafoussam trained 150 youth leaders on electoral participation. After the training, most of them organised awareness-raising sessions in their own communities. This shows that one training course can multiply its impact: \emph{trained people train other people}.
\end{exemplebox}

\courssub{After listening: reformulating with reported speech}

Now use your notes to report what was said. This is where the previous section on reported questions becomes useful. Remember the rules: the tense usually moves one step back (\emph{present simple $\rightarrow$ past simple}), and reported questions use statement word order (SVOMPT) with \emph{if/whether} for yes-no questions.

\begin{center}
\begin{tabularx}{\linewidth}{|X|X|}\hline
\rowcolor{popblueL} \textbf{Direct speech (what was said)} & \textbf{Reported speech (what you say)} \\ \hline
``Democracy needs active citizens.'' & The facilitator explained that democracy needed active citizens. \\ \hline
``Are you ready to share what you have learnt?'' & She asked the participants whether they were ready to share what they had learnt. \\ \hline
``Register on the electoral lists!'' & She encouraged them to register on the electoral lists. \\ \hline
``What did you learn during the workshop?'' & She asked us what we had learnt during the workshop. \\ \hline
\end{tabularx}
\end{center}

\courssub{Short debate: Why should young Cameroonians attend training courses on democracy?}

Work in two groups. Group A defends the idea, Group B asks critical questions. Use the expressions below to speak politely and build your arguments.

\begin{itemize}
\item Giving an opinion: \emph{In my opinion... / I believe that... / As far as I am concerned...}
\item Agreeing: \emph{I totally agree with you. / That is exactly what I think.}
\item Disagreeing politely: \emph{I see your point, but... / I am afraid I disagree because...}
\item Concluding: \emph{To sum up, training courses on democracy help young people to...}
\end{itemize}

Possible arguments: training courses teach us our rights and duties; they prepare us to vote responsibly; they promote tolerance and peace in our communities; they give us certificates that can help our careers; they train future leaders who respect the rule of law and good governance.

\begin{exercice}[Practice]
\textbf{Exercise 1 --- Vocabulary.} Complete the sentences with: \emph{facilitator, certificate, module, tolerance, rule of law, workshop}.\par
1. At the end of the course, each participant received a \ldots\ signed by ELECAM.\par
2. The \ldots\ asked the participants to work in groups of five.\par
3. The first \ldots\ of the training was about democratic values.\par
4. In a democracy, the \ldots\ means that even ministers must respect the law.\par
5. Living together in peace requires \ldots\ towards other cultures and religions.\par
6. Our school organised a two-day \ldots\ on civic education.\par
\textbf{Exercise 2 --- Reported speech.} Report these sentences from the listening text.\par
1. ``The workshop took place in Douala,'' the journalist said.\par
2. ``Did you receive your certificates?'' the facilitator asked the participants.\par
3. ``Why is the rule of law important?'' she asked the class.\par
\textbf{Exercise 3 --- Word order.} Reorder the words to make correct sentences (SVOMPT).\par
1. explained / the facilitator / clearly / the rules / yesterday\par
2. attended / eighty participants / the seminar / in Douala / last week\par
3. received / a certificate / every participant / at the end\par
\textbf{Exercise 4 --- Debate preparation.} Write three sentences giving your opinion on this question: \emph{Should civic education be compulsory in all technical schools in Cameroon?} Use \emph{In my opinion}, \emph{because} and \emph{to sum up}.
\end{exercice}

\begin{corrige}[Exercises 1, 2 and 3]
\textbf{Exercise 1:} 1. certificate; 2. facilitator; 3. module; 4. rule of law; 5. tolerance; 6. workshop.\par
\textbf{Exercise 2:} 1. The journalist said that the workshop had taken place in Douala. 2. The facilitator asked the participants whether (if) they had received their certificates. 3. She asked the class why the rule of law was important.\par
\textbf{Exercise 3:} 1. The facilitator explained the rules clearly yesterday. 2. Eighty participants attended the seminar in Douala last week. 3. Every participant received a certificate at the end.
\end{corrige}

\begin{aretenir}[What to remember]
A training course \textbf{on} democracy includes modules, a facilitator, participants and a certificate at the end. English word order is strict: Subject + Verb + Object + Manner + Place + Time (SVOMPT). To understand a listening text: read the questions first, listen for the gist, then listen selectively for dates, places and numbers, and take notes with key words only. To reformulate, use reported speech: \emph{The facilitator explained that... / She asked us whether...} The five pillars of democracy are free and fair elections, the rule of law, human rights, the separation of powers and a free press --- and they all rest on active, informed citizens like you.
\end{aretenir}

\coursec{Grammar focus 2: Word order}

In the previous section, you listened to texts about \cle{training courses on democracy} (fr: formations sur la démocratie) and enriched your vocabulary with words like \emph{workshop} (fr: atelier), \emph{facilitator} (fr: formateur/animateur), \emph{civic education} (fr: éducation civique) and \emph{training session} (fr: session de formation). You probably noticed something while listening: English sentences follow a very strict architecture. A facilitator does not say ``\emph{Explained the trainer the rules clearly}''. He says: ``\emph{The trainer explained the rules clearly.}'' Word order is not a detail in English --- it \textbf{is} the grammar. This section gives you the complete toolkit to build correct, natural English sentences, an essential skill for the written paper of the Probatoire and Baccalauréat technique.

\courssub{Why word order matters: English is a fixed-order language}

Read these two groups of words. Only one is an English sentence:

\begin{enumerate}
\item \emph{The students elected the president yesterday.}
\item \emph{Elected the president the students yesterday.}
\end{enumerate}

Sentence 2 is impossible in English, even though all the words are correct. Why? Because English words carry almost no endings (no conjugation marks, no agreement) that could tell you who does what. The \textbf{position} of each word does that job. In French, you can sometimes move elements around for style (\emph{Hier, les élèves ont élu le président} / \emph{Les élèves ont élu hier le président}). In English, moving a word usually changes the meaning or destroys the sentence.

\begin{propriete}[English: a fixed-order language]
English is a \cle{fixed word order} language (fr: langue à ordre fixe): the grammatical function of a word (subject, object, etc.) is shown by its \textbf{position}, not by endings. Two direct consequences:
\begin{itemize}
\item The \textbf{subject} almost always comes \textbf{before} the verb, and the \textbf{object} comes \textbf{after} it: \emph{The voters (S) chose (V) the candidate (O).}
\item Contrary to French, the \textbf{adjective always precedes the noun} and never agrees with it: \emph{a free election}, \emph{free elections} (never \emph{elections frees}).
\end{itemize}
\end{propriete}

\begin{attention}[The adjective before the noun]
Francophone students often write \emph{a box wooden} or \emph{a candidate Cameroonian} under the influence of French (\emph{une boîte en bois}, \emph{un candidat camerounais}). In English, the adjective always stands \textbf{before} the noun: \emph{a wooden box}, \emph{a Cameroonian candidate}. There is no exception at your level.
\end{attention}

\courssub{The canonical sentence pattern: S-V-O-Manner-Place-Time}

\courssub{Building the rule from observation}

Look at these sentences taken from our unit on democracy training. Read them aloud and observe the order of the underlined blocks:

\begin{itemize}
\item \emph{The \textbf{trainees} (stagiaires) \textbf{practised} \textbf{a mock election} (élection simulée) \textbf{actively} \textbf{in the hall} \textbf{on Monday}.}
\item \emph{\textbf{Our class} \textbf{organised} \textbf{a debate} \textbf{calmly} \textbf{in the library} \textbf{last week}.}
\item \emph{\textbf{The facilitator} \textbf{explained} \textbf{the voting procedure} \textbf{clearly} \textbf{at the training centre} \textbf{this morning}.}
\end{itemize}

In every sentence, the blocks appear in exactly the same order: \textbf{who} does it, then the \textbf{action}, then \textbf{what}, then \textbf{how}, then \textbf{where}, then \textbf{when}. This is not a coincidence: it is the canonical (standard) order of the English sentence.

\begin{definition}[The canonical order: S-V-O-M-P-T]
A standard affirmative English sentence follows the pattern:
\begin{center}
\textbf{Subject + Verb + Object + Manner + Place + Time}
\end{center}
\begin{itemize}
\item \textbf{Manner} (comment ?) answers \emph{how?} --- \emph{carefully, actively, calmly}.
\item \textbf{Place} (où ?) answers \emph{where?} --- \emph{in the hall, at the training centre}.
\item \textbf{Time} (quand ?) answers \emph{when?} --- \emph{yesterday, on Monday, last week}.
\end{itemize}
French often prefers the order Place--Manner--Time; English prefers \textbf{Manner--Place--Time}. Remember the sequence: \emph{How? Where? When?}
\end{definition}

\begin{exemplebox}[Analysing a full sentence]
\emph{The participants listened carefully to the facilitator in the hall yesterday.}
\begin{center}
\begin{tabularx}{\linewidth}{|l|X|}\hline
\rowcolor{popblueL} \textbf{Block} & \textbf{Words} \\ \hline
Subject & The participants \\ \hline
Verb & listened \\ \hline
Manner & carefully \\ \hline
Place (prepositional phrases) & to the facilitator, in the hall \\ \hline
Time & yesterday \\ \hline
\end{tabularx}
\end{center}
The sentence follows the S-V-Manner-Place-Time pattern perfectly. Note that the object of the verb \emph{listen} is a prepositional phrase (\emph{to the facilitator}), which is placed before the place adjunct (\emph{in the hall}).
\end{exemplebox}

\begin{popfigure}
\begin{tikzpicture}[
block/.style={draw, rounded corners=3pt, minimum height=1.1cm, align=center, font=\small\bfseries, text=white},
lbl/.style={font=\scriptsize\itshape, text=popmuted, align=center},
ex/.style={font=\scriptsize, align=center, text=popdark}
]
\node[block, fill=popblue, minimum width=2.2cm] (s) at (0,0) {SUBJECT};
\node[block, fill=poporange, minimum width=1.8cm] (v) at (2.6,0) {VERB};
\node[block, fill=popgreen, minimum width=2.0cm] (o) at (5.1,0) {OBJECT};
\node[block, fill=poppurple, minimum width=2.0cm] (m) at (7.7,0) {MANNER};
\node[block, fill=poppink, minimum width=1.9cm] (p) at (10.2,0) {PLACE};
\node[block, fill=popgold, minimum width=1.8cm] (t) at (12.6,0) {TIME};
\node[lbl] at (0,-0.85) {who? / what?};
\node[lbl] at (2.6,-0.85) {action};
\node[lbl] at (5.1,-0.85) {who? / what?};
\node[lbl] at (7.7,-0.85) {how?};
\node[lbl] at (10.2,-0.85) {where?};
\node[lbl] at (12.6,-0.85) {when?};
\node[ex] at (0,0.95) {The participants};
\node[ex] at (2.6,0.95) {listened};
\node[ex] at (5.1,0.95) {to the lesson};
\node[ex] at (7.7,0.95) {carefully};
\node[ex] at (10.2,0.95) {in the hall};
\node[ex] at (12.6,0.95) {yesterday};
\draw[-{Stealth}, thick, popmuted] (-1.5,-1.5) -- (14,-1.5);
\end{tikzpicture}
\legende{The canonical English sentence pattern S-V-O-Manner-Place-Time, illustrated with the example sentence ``The participants listened to the lesson carefully in the hall yesterday.'' Each coloured block answers one question: who? -- action -- what? -- how? -- where? -- when?}
\end{popfigure}

\begin{attention}[Time and place at the beginning]
You may put the \textbf{time} (or place) expression at the \textbf{beginning} of the sentence for emphasis, followed by a comma: \emph{Yesterday, the participants listened carefully to the facilitator.} But you may \textbf{never} place time or place between the verb and its object: \emph{The participants listened yesterday to the facilitator} is wrong or very awkward. Rule: \textbf{never separate the verb from its object}.
\end{attention}

\courssub{The position of frequency adverbs}

Adverbs of frequency (\emph{always, usually, often, sometimes, rarely, never}) do not go to the end of the sentence. They have their own reserved seats.

\begin{propriete}[Position of frequency adverbs]
\begin{enumerate}
\item \textbf{Before the main verb} (ordinary verbs): \emph{She \textbf{always attends} the workshops.} / \emph{Voters \textbf{often ask} questions.}
\item \textbf{After the verb \emph{be}}: \emph{She is \textbf{always} punctual.} / \emph{The elections were \textbf{never} violent.}
\item \textbf{After the first auxiliary} (when there is one): \emph{He has \textbf{never} missed a training session.} / \emph{They will \textbf{always} respect the results.}
\end{enumerate}
Note: \emph{sometimes} and \emph{usually} can also open the sentence: \emph{Sometimes, the trainer organises a quiz.}
\end{propriete}

\begin{attention}[The classic Francophone mistake]
Do not calque the French order \emph{Je vais souvent à l'école}. The sentence \emph{I go often to school} is wrong. Say: \textbf{\emph{I often go to school.}} The adverb of frequency sticks to the verb, just before it (or just after \emph{be}): \emph{We \textbf{never cheat} during elections}, \emph{Our facilitator \textbf{is always} patient}.
\end{attention}

\courssub{The order of adjectives before a noun: OSASCOM}

When several adjectives describe the same noun, English does not place them randomly. Native speakers follow a hidden but very stable order. Compare:

\begin{itemize}
\item \emph{a small old wooden ballot box} (natural)
\item \emph{a wooden old small ballot box} (impossible)
\end{itemize}

\begin{methode}[Memorise OSASCOM]
To order adjectives before a noun, use the acronym \cle{OSASCOM}:
\begin{enumerate}
\item \textbf{O} -- \textbf{Opinion} (ton avis): \emph{beautiful, fair, useful, honest}
\item \textbf{S} -- \textbf{Size} (taille): \emph{big, small, long}
\item \textbf{A} -- \textbf{Age} (âge): \emph{old, new, young}
\item \textbf{S} -- \textbf{Shape} (forme): \emph{round, square, rectangular}
\item \textbf{C} -- \textbf{Colour} (couleur): \emph{red, white, green}
\item \textbf{O} -- \textbf{Origin} (origine): \emph{Cameroonian, African, French}
\item \textbf{M} -- \textbf{Material} (matière): \emph{wooden, plastic, metal}
\end{enumerate}
Then comes the \textbf{noun}. Example: \emph{a \textbf{small} (size) \textbf{old} (age) \textbf{Cameroonian} (origin) \textbf{wooden} (material) ballot box}. In practice, avoid using more than two or three adjectives at once.
\end{methode}

\begin{exemplebox}[Applying OSASCOM]
\begin{itemize}
\item \emph{a \textbf{beautiful} (opinion) \textbf{new} (age) \textbf{white} (colour) ballot box}
\item \emph{an \textbf{honest} (opinion) \textbf{young} (age) \textbf{Cameroonian} (origin) candidate}
\item \emph{a \textbf{large} (size) \textbf{rectangular} (shape) \textbf{metal} (material) suggestion box}
\end{itemize}
\end{exemplebox}

\courssub{Word order in questions}

\courssub{The basic pattern: auxiliary + subject + verb}

In a question, English moves the \textbf{auxiliary} before the subject. If the sentence has no auxiliary, English creates one: \emph{do / does / did}.

\begin{propriete}[Question word order]
\begin{enumerate}
\item \textbf{Yes/No questions}: \emph{Auxiliary + Subject + Verb + ...?}\\
\emph{\textbf{Did} the students \textbf{elect} their class leader?} / \emph{\textbf{Have} you \textbf{ever attended} a democracy workshop?}
\item \textbf{Wh- questions}: \emph{Wh-word + Auxiliary + Subject + Verb + ...?}\\
\emph{\textbf{When did} the training course \textbf{start}?} / \emph{\textbf{What does} ``citizenship'' \textbf{mean}?}
\item \textbf{Subject questions}: when the Wh-word \emph{is} the subject, there is \textbf{no auxiliary} and no inversion: \emph{\textbf{Who won} the election?} (not \emph{Who did win}), \emph{\textbf{How many students attended} the course?}
\end{enumerate}
\end{propriete}

\begin{attention}[No inversion in indirect-style word order]
After the auxiliary, the verb returns to its base form: \emph{Did you \textbf{vote}} (not \emph{Did you voted}). And never invert when the question word is the subject: \emph{Who \textbf{called} the meeting?} is correct; \emph{Who did call the meeting?} is not.
\end{attention}

\courssub{Questions ending with a preposition}

In English, the preposition of a phrasal or prepositional verb stays at the \textbf{end} of the question; it cannot move to the front in everyday English.

\begin{exemplebox}[Final prepositions]
\begin{itemize}
\item \emph{\textbf{Who} did you vote \textbf{for}?} (Pour qui as-tu voté ?)
\item \emph{\textbf{What} are you talking \textbf{about}?}
\item \emph{\textbf{Which course} are you interested \textbf{in}?}
\item \emph{\textbf{Who} did she campaign \textbf{with}?}
\end{itemize}
Formal written English allows \emph{For whom did you vote?}, but at your level, prefer the natural form with the final preposition.
\end{exemplebox}

\courssub{Inversion after restrictive adverbs (Bac complement)}

For elegant, emphatic writing --- a real bonus at the Baccalauréat --- English inverts the subject and the auxiliary when the sentence \textbf{begins} with a restrictive or negative adverbial expression: \emph{never, rarely, seldom, hardly, no sooner, not only, only then, under no circumstances}.

\begin{propriete}[Inversion after initial restrictive adverbs]
Structure: \emph{Restrictive adverb + auxiliary + subject + verb}.
\begin{itemize}
\item \emph{\textbf{Never have I seen} such a fair election.}
\item \emph{\textbf{Rarely do} young people \textbf{take part} in polls with so much enthusiasm.}
\item \emph{\textbf{Not only did} the facilitator \textbf{explain} the rules, but he also organised a mock vote.}
\item \emph{\textbf{Under no circumstances should} voters \textbf{be} intimidated.}
\end{itemize}
If there is no auxiliary, add \emph{do/does/did}: \emph{Seldom \textbf{does} he \textbf{miss} a training session.}
\end{propriete}

\begin{exoresolu}[Reordering jumbled sentences]
Put the words in the correct order to make sentences about our unit.
\begin{enumerate}
\item \emph{attended / the students / a workshop / actively / yesterday / at the youth centre}
\item \emph{is / our facilitator / punctual / always}
\item \emph{ballot / a / bought / wooden / they / old / box / small}
\item \emph{for / did / who / vote / you / ?}
\item \emph{seen / never / I / have / transparent / such / a / poll}
\end{enumerate}
\tcblower
\textbf{Solution.}
\begin{enumerate}
\item \emph{The students attended a workshop actively at the youth centre yesterday.} (S-V-O-Manner-Place-Time)
\item \emph{Our facilitator is always punctual.} (frequency adverb after \emph{be})
\item \emph{They bought a small old wooden ballot box.} (OSASCOM: size -- age -- material)
\item \emph{Who did you vote for?} (auxiliary + subject + verb, final preposition)
\item \emph{Never have I seen such a transparent poll.} (inversion after initial \emph{never})
\end{enumerate}
\end{exoresolu}

\begin{savaistu}[No verb, no sentence!]
In English, a group of words without a \textbf{conjugated verb} is not a sentence: it is a \cle{fragment} (fr: fragment de phrase), and it is considered a mistake in writing. \emph{The students in the hall yesterday.} is a fragment --- nothing happens in it! Add a verb and it becomes a sentence: \emph{The students \textbf{were} in the hall yesterday.} This is why English learners are taught: \emph{``No verb, no sentence.''} Always check that every sentence you write in your composition contains at least one conjugated verb.
\end{savaistu}

\begin{aretenir}[Word order: the five golden rules]
\begin{enumerate}
\item Affirmative sentence: \textbf{Subject + Verb + Object + Manner + Place + Time} --- never separate the verb from its object.
\item Frequency adverbs: \textbf{before} the main verb, \textbf{after} \emph{be} or the first auxiliary.
\item Several adjectives: follow \textbf{OSASCOM} (Opinion, Size, Age, Shape, Colour, Origin, Material) + noun.
\item Questions: \textbf{auxiliary + subject + verb}; prepositions stay at the end (\emph{Who did you vote for?}); no auxiliary when the Wh-word is the subject (\emph{Who won?}).
\item Emphasis: after an initial restrictive adverb (\emph{never, rarely, not only...}), invert auxiliary and subject: \emph{Never have I seen...}
\end{enumerate}
\end{aretenir}

\begin{exercice}[Practice: campaigns, training and surveys]
\textbf{A. Put the words in order.}
\begin{enumerate}
\item \emph{the candidates / spoke / convincingly / during the campaign / yesterday}
\item \emph{often / we / part / take / in surveys}
\item \emph{a / organised / they / new / interesting / training course}
\item \emph{about / what / talking / is / the facilitator / ?}
\item \emph{have / rarely / I / met / a / dedicated / so / volunteer}
\end{enumerate}
\textbf{B. Correct the mistakes.} One sentence is correct.
\begin{enumerate}
\item \emph{She goes always to the polling station early.}
\item \emph{They elected yesterday their new class president.}
\item \emph{The trainer explained clearly the rules in the hall this morning.}
\item \emph{Who did vote for the motion?}
\end{enumerate}
\textbf{C. Writing.} Using S-V-O-M-P-T, write three sentences about a democracy training course you attended (or imagine one). Underline each block with a different colour or label it.
\end{exercice}

\begin{corrige}[Exercice: Practice]
\textbf{A.}
\begin{enumerate}
\item \emph{The candidates spoke convincingly during the campaign yesterday.}
\item \emph{We often take part in surveys.}
\item \emph{They organised an interesting new training course.} (opinion before age)
\item \emph{What is the facilitator talking about?}
\item \emph{Rarely have I met such a dedicated volunteer.}
\end{enumerate}
\textbf{B.}
\begin{enumerate}
\item \emph{She always goes to the polling station early.} (adverb before the main verb)
\item \emph{They elected their new class president yesterday.} (verb and object stay together)
\item Correct.
\item \emph{Who voted for the motion?} (subject question: no auxiliary, no inversion)
\end{enumerate}
\textbf{C.} Example answer: \emph{Our class (S) attended (V) a civic education workshop (O) actively (M) at the community centre (P) last Saturday (T).} Check that every sentence has one conjugated verb and respects the block order.
\end{corrige}

You now master the architecture of the English sentence. In the next section, \emph{Reading \& Vocabulary: Volunteering in gender equality promotion}, you will meet these structures inside authentic texts --- watch how often the S-V-O-M-P-T pattern and the frequency adverbs appear!

\coursec{Reading \& Vocabulary: Volunteering in gender equality promotion}

In the previous section, we revised \cle{word order} in sentences, a key skill for writing clear and correct English. Now, we move from sentence structure to \cle{text comprehension} and \cle{thematic vocabulary}. We will read about a powerful form of citizenship: \cle{volunteering} to promote \cle{gender equality}. This topic connects directly to \cle{Human Rights} (Article 3 of the Universal Declaration of Human Rights) and to the reality of many young Cameroonians engaged in community service.

\courssub{Key Vocabulary: Volunteering and Gender Equality}

Before reading the text, let us master the essential words and expressions. Notice the French glosses for key terms.

\begin{definition}[Core Vocabulary -- Volunteering \& Gender]
\begin{center}
\begin{tabularx}{\linewidth}{|l|X|l|}\hline
\textbf{English Term} & \textbf{Definition / Context} & \textbf{Glossaire (Français)} \\ \hline
\cle{Volunteer} (n.) & A person who works without pay. & Bénévole, volontaire \\ \hline
\cle{To volunteer} (v.) & To offer to do something freely. & Faire du bénévolat, se porter volontaire \\ \hline
\cle{NGO} (n.) & Non-Governmental Organisation. & ONG (Organisation Non Gouvernementale) \\ \hline
\cle{Community service} (n.) & Unpaid work to help a community. & Service communautaire, travail d'intérêt général \\ \hline
\cle{Gender equality} (n.) & Equal rights, chances, duties for women and men. & Égalité des sexes / genres \\ \hline
\cle{Empowerment} (n.) & Giving power/confidence to someone. & Autonomisation, renforcement des capacités \\ \hline
\cle{Discrimination} (n.) & Unfair treatment based on gender, race, etc. & Discrimination \\ \hline
\cle{Awareness campaign} (n.) & Actions to inform people about an issue. & Campagne de sensibilisation \\ \hline
\cle{To advocate} (v.) & To publicly support a cause. & Plaider pour, défendre une cause \\ \hline
\cle{Women's rights} (n.) & Rights promoting equality for women. & Droits des femmes \\ \hline
\cle{Girls' education} (n.) & Access to schooling for girls. & Scolarisation des filles \\ \hline
\cle{Role model} (n.) & A person others imitate/admire. & Modèle, exemple à suivre \\ \hline
\end{tabularx}
\end{center}
\end{definition}

\begin{attention}[Grammar Trap: Noun vs Verb / Adjective vs Adverb]
Do not confuse these forms:
\begin{itemize}
    \item \textbf{A volunteer} (Noun, \emph{un bénévole}) vs \textbf{To volunteer} (Verb, \emph{faire du bénévolat}).
    \item \textbf{Voluntary} (Adjective, \emph{bénévole, volontaire}) : \textit{Voluntary work}.
    \item \textbf{Voluntarily} (Adverb, \emph{volontairement}) : \textit{He left voluntarily}.
\end{itemize}
\textbf{Test yourself:} "She is a \_\_\_\_\_ at the Red Cross." (volunteer) / "They \_\_\_\_\_ every summer." (volunteer) / "It was a \_\_\_\_\_ decision." (voluntary).
\end{attention}

\courssub{Reading Strategies: Skimming, Scanning, and Connectors}

When you face a text in an exam (Probatoire/Bac), do not read every word slowly three times. Use strategies:

\begin{methode}[Three-Step Reading Strategy]
\textbf{Step 1: Skimming (General Idea)} -- Read the \textbf{title}, the \textbf{first sentence} of each paragraph, and the \textbf{last sentence}. Ask: "What is the \cle{main topic}? Who? Where? Why?"
\textbf{Step 2: Scanning (Specific Info)} -- Look for \textbf{keywords} from the questions (numbers, names, dates, specific verbs). Move your eyes fast; do not read every word.
\textbf{Step 3: Detailed Reading (Logic \& Connectors)} -- Read fully to answer \textit{True/False Justified} or \textit{Open Questions}. Spot \cle{logical connectors} (linking words) to understand the argument:
\begin{itemize}
    \item \textbf{Addition:} \cle{Furthermore}, \cle{Moreover}, \cle{In addition}, \cle{Also}.
    \item \textbf{Contrast:} \cle{However}, \cle{Although}, \cle{On the other hand}, \cle{Despite}.
    \item \textbf{Cause/Effect:} \cle{Therefore}, \cle{Consequently}, \cle{As a result}, \cle{Because of}.
    \item \textbf{Example:} \cle{For instance}, \cle{such as}, \cle{including}.
\end{itemize}
\end{methode}

\courssub{Reading Text: "Voices of Change in the Far North"}

Read the following text. Apply the method: Skim first (30 seconds), then scan for details.

\begin{exemplebox}[Reading Text: Voices of Change in the Far North]
\textbf{Title: Voices of Change in the Far North}

[1] In the dusty streets of Maroua, capital of the Far North Region, a group of young volunteers walks from compound to compound under a scorching sun. They belong to "Youth for Equality", a local \cle{NGO} supported by the \cle{MINPROFF} (Ministry of Women's Empowerment and the Family). Their mission is simple but difficult: \cle{advocate} for \cle{Girls' education} in communities where early marriage and poverty keep daughters at home.

[2] \textbf{Furthermore}, cultural norms often dictate that a girl's place is in the kitchen, not the classroom. \textbf{However}, the volunteers -- both young men and women -- use a powerful tool: dialogue. They sit with traditional chiefs and parents, explaining that an educated girl becomes an empowered woman who contributes to the family income and community development. "We do not fight tradition," says Aïcha, 22, a volunteer and university student. "We show that \cle{gender equality} benefits everyone."

[3] \textbf{For instance}, during the last \cle{awareness campaign} in June 2024, 300 volunteers reached over 5,000 households in the Diamaré division. They organised focus groups, radio talks in Fulfulde and French, and school visits. \textbf{As a result}, local authorities reported a 15\% increase in girls' enrolment in targeted primary schools for the 2024-2025 academic year.

[4] \textbf{Despite} these gains, \cle{discrimination} persists. Some families still hide their daughters when volunteers arrive. \textbf{Therefore}, the NGO now trains "community relays" -- respected local men and women -- to continue the work after the volunteers leave. This ensures \cle{empowerment} is sustainable. These relays become \cle{role models}, proving that \cle{volunteering} is not just \cle{community service}; it is a fight for fundamental \cle{human rights}.
\end{exemplebox}

\courssub{Comprehension Activities}

\begin{exercice}[Comprehension Questions]
\textbf{A. True or False? Justify by quoting the exact line(s) from the text.}
\begin{enumerate}
    \item The volunteers in Maroua are paid employees of the government.
    \item Cultural norms in the Far North fully support girls' education.
    \item Aïcha believes gender equality is only good for women.
    \item The June 2024 campaign involved 300 volunteers.
    \item The campaign used only French to communicate.
    \item The NGO trains "community relays" to make the action sustainable.
\end{enumerate}

\textbf{B. Open Questions (Answer in your own words, using complete sentences).}
\begin{enumerate}
    \item What is the name of the NGO mentioned in the text and which ministry supports it?
    \item List three specific activities organised during the awareness campaign.
    \item According to the text, what was the concrete result of the campaign on school enrolment?
    \item Why does the NGO train "community relays"?
    \item Find in the text a sentence showing that volunteering is linked to human rights.
\end{enumerate}

\textbf{C. Vocabulary Equivalences (Match the word from the text to its definition).}
\begin{enumerate}
    \item \cle{Advocate} (l. 4)
    \item \cle{Empowerment} (l. 14)
    \item \cle{Scorching} (l. 2)
    \item \cle{Sustainable} (l. 15)
    \item \cle{Compound} (l. 2)
    \end{enumerate}
    a) Very hot. \quad b) To speak in favour of. \quad c) A family enclosure/household. \quad d) Giving power/autonomy. \quad e) Lasting, able to continue.
\end{exercice}

\begin{corrige}[Exercice Comprehension]
\textbf{A. True / False Justified}
\begin{enumerate}
    \item \textbf{False.} Line 3: "They belong to 'Youth for Equality', a local NGO..." / Line 4: "Their mission is simple...". They are volunteers, not paid government employees.
    \item \textbf{False.} Line 7: "cultural norms often dictate that a girl's place is in the kitchen, not the classroom."
    \item \textbf{False.} Line 9: "We show that gender equality benefits \textbf{everyone}."
    \item \textbf{True.} Line 11: "during the last awareness campaign in June 2024, \textbf{300 volunteers} reached..."
    \item \textbf{False.} Line 12: "radio talks in \textbf{Fulfulde and French}". 
    \item \textbf{True.} Line 15: "the NGO now trains 'community relays'... to continue the work after the volunteers leave. This ensures empowerment is \textbf{sustainable}."
\end{enumerate}

\textbf{B. Open Questions (Sample Answers)}
\begin{enumerate}
    \item The NGO is "Youth for Equality". It is supported by the MINPROFF (Ministry of Women's Empowerment and the Family).
    \item Focus groups, radio talks (in Fulfulde and French), and school visits.
    \item A 15\% increase in girls' enrolment in targeted primary schools for the 2024-2025 academic year.
    \item To continue the work after the volunteers leave and ensure empowerment is sustainable.
    \item Line 17: "...proving that volunteering is not just community service; it is a fight for fundamental human rights."
\end{enumerate}

\textbf{C. Vocabulary Equivalences}
1 -- b (\cle{Advocate} = To speak in favour of). \quad 2 -- d (\cle{Empowerment} = Giving power/autonomy). \quad 3 -- a (\cle{Scorching} = Very hot). \quad 4 -- e (\cle{Sustainable} = Lasting). \quad 5 -- c (\cle{Compound} = Family enclosure/household).
\end{corrige}

\courssub{Focus on Logical Connectors (Text Analysis)}

The text uses connectors to build a coherent argument. Identifying them helps you understand the \cle{logic} of the writer.

\begin{exoresolu}[Analyzing Connectors]
\textbf{Task:} Identify the connector in each sentence below and explain its function (Contrast, Addition, Consequence, Example).

\begin{enumerate}
    \item "\textbf{Furthermore}, cultural norms often dictate..." (Paragraph 2)
    \item "\textbf{However}, the volunteers... use a powerful tool: dialogue." (Paragraph 2)
    \item "\textbf{For instance}, during the last awareness campaign..." (Paragraph 3)
    \item "\textbf{As a result}, local authorities reported a 15\% increase..." (Paragraph 3)
    \item "\textbf{Despite} these gains, discrimination persists." (Paragraph 4)
    \item "\textbf{Therefore}, the NGO now trains community relays..." (Paragraph 4)
\end{enumerate}

\textbf{Solution:}
\begin{enumerate}
    \item \textbf{Furthermore} -- \cle{Addition}. Adds a second problem (cultural norms) to the first (poverty/early marriage mentioned in Para 1).
    \item \textbf{However} -- \cle{Contrast}. Opposes the negative context (norms) to the positive action (volunteers using dialogue).
    \item \textbf{For instance} -- \cle{Example}. Introduces a concrete illustration of the work described generally before.
    \item \textbf{As a result} -- \cle{Consequence}. Shows the effect (enrolment increase) of the cause (campaign activities).
    \item \textbf{Despite} -- \cle{Concession/Contrast}. Introduces a persisting problem even though there were gains.
    \item \textbf{Therefore} -- \cle{Consequence/Conclusion}. Introduces the strategic response (training relays) to the persisting problem.
\end{enumerate}
\end{exoresolu}

\courssub{Illustration 1: Comparative Enrollment Data (pgfplots)}

Data visualization helps understand the \cle{context} of the text. The chart below shows \cle{fictional but realistic data} inspired by INS (Institut National de la Statistique) trends for the 2023-2024 school year. It highlights the \cle{gender gap} the volunteers are fighting.

\begin{popfigure}
\begin{tikzpicture}
\begin{axis}[
    ybar,
    enlargelimits=0.15,
    legend style={at={(0.5,-0.2)}, anchor=north, legend columns=-1},
    ylabel={Gross Enrollment Rate (\%)},
    symbolic x coords={Far North, North, Adamawa, East, Centre, Littoral, West, South West, North West, South},
    xtick=data,
    x tick label style={rotate=45, anchor=east, font=\small},
    ymin=0, ymax=110,
    bar width=6pt,
    width=\linewidth,
    height=8cm,
    major grid style={draw=popmuted},
    ymajorgrids=true,
    xmajorgrids=false,
]
\addplot[popblue, fill=popblue!60] coordinates {
    (Far North, 62) (North, 70) (Adamawa, 75) (East, 72) (Centre, 95)
    (Littoral, 98) (West, 96) (South West, 85) (North West, 82) (South, 90)
};
\addplot[poppink, fill=poppink!60] coordinates {
    (Far North, 48) (North, 58) (Adamawa, 65) (East, 60) (Centre, 92)
    (Littoral, 97) (West, 94) (South West, 80) (North West, 78) (South, 88)
};
\legend{Boys, Girls}
\end{axis}
\end{tikzpicture}
\legende{Figure 1: Gross Primary Enrollment Rates by Region and Gender (Cameroon, Fictional Data 2023-2024). The gap is widest in the Far North, justifying targeted volunteer action.}
\end{popfigure}

\begin{aretenir}[Reading the Chart]
\begin{itemize}
    \item \textbf{Far North:} Boys 62\% vs Girls 48\% (\textbf{Gap: 14 points}). This is the zone of the text.
    \item \textbf{North / Adamawa / East:} Significant gaps (10-12 points).
    \item \textbf{Centre / Littoral / West:} Near parity (Gap < 3 points). Urban advantage.
    \item \textbf{Analysis:} Volunteering is most critical where the gap is widest (Far North, North).
\end{itemize}
\end{aretenir}

\courssub{Illustration 2: Volunteer Action Zones Map (TikZ)}

The text mentions Maroua (Far North). Volunteering happens nationwide. This simplified map locates three key intervention zones.

\begin{popfigure}
\begin{tikzpicture}[scale=0.9]
% Simplified Cameroon Outline (Polygon approximation)
\draw[popdark, thick, fill=popcream] 
    (0,0) -- (1,0.2) -- (2,0.5) -- (3,1) -- (4,1.5) -- (5,2) -- (5.5,2.5) -- (5.8,3) -- (5.5,3.5) -- (5,4) -- (4.5,4.5) -- (3.5,5) -- (2.5,5.2) -- (1.5,5) -- (0.5,4.5) -- (0,4) -- (-0.5,3) -- (-0.5,2) -- (0,1.5) -- (0.5,1) -- (1,0.5) -- cycle;
% Regions Labels (Approximate)
\node[font=\small\bfseries, popdark] at (2.5, 4.2) {Far North};
\node[font=\small\bfseries, popdark] at (2.5, 3.0) {North};
\node[font=\small\bfseries, popdark] at (3.5, 2.0) {Adamawa};
\node[font=\small\bfseries, popdark] at (4.5, 3.0) {East};
\node[font=\small\bfseries, popdark] at (3.0, 1.5) {Centre};
\node[font=\small\bfseries, popdark] at (2.0, 1.0) {Littoral};
\node[font=\small\bfseries, popdark] at (1.5, 2.0) {West};
\node[font=\small\bfseries, popdark] at (1.0, 2.8) {South West};
\node[font=\small\bfseries, popdark] at (0.5, 3.5) {North West};
\node[font=\small\bfseries, popdark] at (1.5, 3.8) {South};
% Capital Yaoundé
\node[circle, fill=popdark, inner sep=1.5pt] (Yde) at (3.0, 1.5) {};
\node[anchor=west, font=\small] at (3.1, 1.5) {Yaoundé};
% Zone 1: Far North (Maroua) - RED
\node[circle, fill=poporange, inner sep=2.5pt] (Mar) at (2.2, 4.0) {};
\node[anchor=west, font=\small\bfseries, poporange] at (2.4, 4.1) {Maroua (Far North)};
\draw[poporange, thick, ->, shorten >=2pt] (2.4, 3.8) -- (2.2, 3.5) node[below, font=\tiny, align=center] {Zone 1: Girl Education \\ Awareness Campaigns};
% Zone 2: Littoral (Douala) - BLUE
\node[circle, fill=popblue, inner sep=2.5pt] (Dou) at (1.8, 1.0) {};
\node[anchor=east, font=\small\bfseries, popblue] at (1.7, 1.1) {Douala (Littoral)};
\draw[popblue, thick, ->, shorten >=2pt] (1.6, 1.2) -- (1.2, 1.5) node[left, font=\tiny, align=center] {Zone 2: Urban \\ Gender-Based Violence \\ Prevention};
% Zone 3: Centre (Yaoundé) - GREEN
\node[circle, fill=popgreen, inner sep=2.5pt] (Yde2) at (3.0, 1.5) {}; % Same as capital
\node[anchor=south, font=\small\bfseries, popgreen] at (3.0, 1.7) {Yaoundé (Centre)};
\draw[popgreen, thick, ->, shorten >=2pt] (3.2, 1.8) -- (3.8, 2.0) node[right, font=\tiny, align=center] {Zone 3: School Clubs \\ \& Leadership Training};
% North Arrow
\node[font=\small] at (5.5, 4.5) {\textbf{N}};
\draw[->, thick] (5.5, 4.3) -- (5.5, 4.6);
% Scale (Approximate)
\draw[popdark] (0.5, 0.2) -- (1.5, 0.2) node[midway, below, font=\tiny] {~200 km};
\end{tikzpicture}
\legende{Figure 2: Simplified Map of Cameroon showing Three Volunteer Intervention Zones. Zone 1 (Far North) matches the reading text. Zones 2 \& 3 show other typical NGO activities.}
\end{popfigure}

\courssub{Discussion: Youth as Agents of Change}

The text shows young people (Aïcha, 22) leading change. In your context (Technical High School), you are also agents of change.

\begin{exercice}[Oral / Written Discussion]
\textbf{Work in pairs or small groups. Discuss and take notes.}
\begin{enumerate}
    \item \textbf{At School:} How can boys and girls in Terminale promote gender equality in your technical workshops (atelier), classrooms, or student government (bureau des élèves)?
    \item \textbf{In the Community:} What \cle{volunteering} activities could your class organise during holidays to support \cle{Girls' education} or fight \cle{discrimination}?
    \item \textbf{Role Models:} Who is a \cle{role model} for you (famous or local) regarding gender equality? Why?
\end{enumerate}
\textbf{Report:} Choose a spokesperson to present 3 concrete ideas to the class. Use vocabulary: \textit{advocate, raise awareness, empower, volunteer, NGO, community service}.
\end{exercice}

\courssub{Link with Human Rights: Article 3 UDHR}

Volunteering for gender equality is not just charity; it is defending a legal right.

\begin{savaistu}[Cultural \& Legal Context: Cameroon \& Human Rights]
\textbf{Article 3 of the Universal Declaration of Human Rights (1948):}
\begin{quote}
\textit{"Everyone has the right to life, \textbf{liberty and security of person}."}
\end{quote}
\textbf{Article 26 (Education):}
\begin{quote}
\textit{"Everyone has the right to \textbf{education}... Education shall be directed to the full development of the human personality and to the strengthening of respect for human rights and fundamental freedoms."}
\end{quote}

\textbf{Cameroon Specifics:}
\begin{itemize}
    \item Cameroon ratified \cle{CEDAW} (Convention on the Elimination of All Forms of Discrimination Against Women) in 1994.
    \item The \cle{Constitution} (Preamble) guarantees equality of rights for all citizens.
    \item \textbf{MINPROFF} (Ministère de la Promotion de la Femme et de la Famille) is the dedicated ministry created in 2004 (Decree 2004/320). It designs national policies for gender mainstreaming.
    \item \textbf{National Gender Policy (2011-2020, revised)} aims to eliminate disparities in education, health, and economic participation.
\end{itemize}
\textbf{Connection:} When you \cle{volunteer} for an \cle{awareness campaign}, you are helping the State implement its international and national legal obligations. You are a \cle{human rights defender}.
\end{savaistu}

\courssub{Writing Preparation: Composition on Volunteering}

This reading prepares you for the \cle{Writing} task (Section 6 of the lesson plan): \textit{"Write a composition about volunteering in gender equality promotion activities."} You will need:
\begin{itemize}
    \item \cle{Vocabulary} from this section (advocate, empowerment, NGO, awareness campaign...).
    \item \cle{Grammar}: \cle{Phrasal verbs} (e.g., \textit{bring about} change, \textit{carry out} a campaign, \textit{reach out to} communities), \cle{Sequence of tenses} (reporting past events), \cle{Impersonal passives} (e.g., \textit{It is reported that..., Volunteers are needed...}).
    \item \cle{Structure}: Introduction (Context/Hook) $\rightarrow$ Body Paragraphs (Actions, Challenges, Results, Connectors) $\rightarrow$ Conclusion (Call to action / Personal commitment).
\end{itemize}

\begin{exemplebox}[Writing Plan Outline]
\textbf{Topic:} "Volunteering for Gender Equality: A Duty for Youth."
\begin{enumerate}
    \item \textbf{Intro:} Define gender equality (Human Right). State the problem (gap in Far North). Thesis: Youth volunteering bridges the gap.
    \item \textbf{Body 1 (Action):} Describe activities (campaigns, dialogue with chiefs, radio). Use \textit{phrasal verbs}: \textit{carry out, reach out to, speak out}.
    \item \textbf{Body 2 (Impact/Results):} Cite evidence (enrolment increase). Use \textit{impersonal passive}: \textit{It has been observed that... / A rise was recorded...}. Use \textit{sequence of tenses} for reported speech.
    \item \textbf{Body 3 (Challenges/Sustainability):} Mention persistence of discrimination. Solution: Training relays/role models. Use \textit{connectors}: \textit{However, Therefore, Consequently}.
    \item \textbf{Conclusion:} Summary. Personal pledge: "As a technical student, I will..."
\end{enumerate}
\end{exemplebox}

\courssub{Summary Checklist}

\begin{aretenir}[Key Takeaways for Section 5]
\begin{itemize}
    \item \textbf{Vocabulary:} Master the 11 key terms (Volunteer, NGO, Gender equality, Empowerment, Discrimination, Awareness campaign, Advocate, Women's rights, Girls' education, Role model, Community service).
    \item \textbf{Reading Skills:} Skim (Main idea) $\rightarrow$ Scan (Details) $\rightarrow$ Analyse Connectors (Logic).
    \item \textbf{Text Content:} "Youth for Equality" in Maroua (Far North). 300 volunteers, 5000 households, 15\% enrolment rise. Strategy: Dialogue + Community Relays.
    \item \textbf{Grammar Links:} Connectors (\textit{Furthermore, However, As a result, Despite, Therefore}) organize the argument.
    \item \textbf{Citizenship:} Volunteering = Active Citizenship = Human Rights Defence (UDHR Art. 3, 26; CEDAW; MINPROFF).
\end{itemize}
\end{aretenir}

\coursec{Grammar focus 3 \& Writing: Phrasal verbs, sequence of tenses, impersonal passives}

In the previous section, we read texts about volunteers who promote gender equality in their communities. We met sentences like \emph{``She joined the NGO to promote girls' education''} and \emph{``It is reported that more girls now go to school.''} These sentences contain three powerful grammar tools: \cle{phrasal verbs} used to express purpose, the \cle{sequence of tenses}, and the \cle{impersonal passive}. In this final section, we will master these three tools and then use them to write a full composition about volunteering in gender equality promotion activities — exactly the kind of writing task you will meet at the Probatoire and the Baccalauréat technique.

\courssub{Phrasal verbs in simple expressions of purpose}

\textbf{What is a phrasal verb?}\par

Look at these two sentences:

\begin{itemize}
\item \emph{The volunteers \textbf{carried} heavy boxes.} (carried = transported)
\item \emph{The volunteers \textbf{carried out} a campaign.} (carried out = conducted, organised)
\end{itemize}

The verb \emph{carry} changes its meaning completely when we add the particle \emph{out}. This is the key intuition: a phrasal verb is a team of words whose meaning is different from the meaning of the verb alone.

\begin{definition}[Phrasal verb]
A \cle{phrasal verb} is a verb combined with a particle (an adverb or a preposition) whose global meaning is different from the meaning of the simple verb. Examples: \emph{carry out} (conduct), \emph{set up} (create), \emph{give out} (distribute).
\end{definition}

\textbf{Key phrasal verbs for citizenship and volunteering}\par

Here are the phrasal verbs you must know for this unit. Learn them with their meanings and one example sentence each.

\begin{popfigure}
\begin{tikzpicture}
\draw[fill=popblueL] (0,0) rectangle (3.4,-0.7);
\draw[fill=popblueL] (3.4,0) rectangle (7.6,-0.7);
\draw[fill=popblueL] (7.6,0) rectangle (15.2,-0.7);
\node at (1.7,-0.35) {\textbf{Phrasal verb}};
\node at (5.5,-0.35) {\textbf{Meaning}};
\node at (11.4,-0.35) {\textbf{Example sentence}};
\draw[fill=popcream] (0,-0.7) rectangle (3.4,-1.5);
\draw[fill=popcream] (3.4,-0.7) rectangle (7.6,-1.5);
\draw[fill=popcream] (7.6,-0.7) rectangle (15.2,-1.5);
\node at (1.7,-1.1) {\emph{carry out}};
\node at (5.5,-1.1) {conduct, do};
\node[text width=7.2cm,align=left] at (11.4,-1.1) {\emph{They carried out a campaign in schools.}};
\draw (0,-1.5) rectangle (3.4,-2.3);
\draw (3.4,-1.5) rectangle (7.6,-2.3);
\draw (7.6,-1.5) rectangle (15.2,-2.3);
\node at (1.7,-1.9) {\emph{set up}};
\node at (5.5,-1.9) {create, establish};
\node[text width=7.2cm,align=left] at (11.4,-1.9) {\emph{We set up a girls' club in our school.}};
\draw[fill=popcream] (0,-2.3) rectangle (3.4,-3.1);
\draw[fill=popcream] (3.4,-2.3) rectangle (7.6,-3.1);
\draw[fill=popcream] (7.6,-2.3) rectangle (15.2,-3.1);
\node at (1.7,-2.7) {\emph{look after}};
\node at (5.5,-2.7) {take care of};
\node[text width=7.2cm,align=left] at (11.4,-2.7) {\emph{She looks after vulnerable children.}};
\draw (0,-3.1) rectangle (3.4,-3.9);
\draw (3.4,-3.1) rectangle (7.6,-3.9);
\draw (7.6,-3.1) rectangle (15.2,-3.9);
\node at (1.7,-3.5) {\emph{give out}};
\node at (5.5,-3.5) {distribute};
\node[text width=7.2cm,align=left] at (11.4,-3.5) {\emph{They gave out leaflets at the market.}};
\draw[fill=popcream] (0,-3.9) rectangle (3.4,-4.7);
\draw[fill=popcream] (3.4,-3.9) rectangle (7.6,-4.7);
\draw[fill=popcream] (7.6,-3.9) rectangle (15.2,-4.7);
\node at (1.7,-4.3) {\emph{bring about}};
\node at (5.5,-4.3) {cause, produce};
\node[text width=7.2cm,align=left] at (11.4,-4.3) {\emph{Education brings about real change.}};
\draw (0,-4.7) rectangle (3.4,-5.5);
\draw (3.4,-4.7) rectangle (7.6,-5.5);
\draw (7.6,-4.7) rectangle (15.2,-5.5);
\node at (1.7,-5.1) {\emph{stand up for}};
\node at (5.5,-5.1) {defend};
\node[text width=7.2cm,align=left] at (11.4,-5.1) {\emph{We must stand up for women's rights.}};
\end{tikzpicture}
\legende{Six essential phrasal verbs for talking about volunteering and citizenship. Note also \emph{take part in} (participate in): \emph{Many students take part in awareness activities.}}
\end{popfigure}

\begin{attention}[Separable or inseparable?]
Some phrasal verbs can be separated by an object (\emph{give out leaflets} / \emph{give leaflets out} / \emph{give them out}), but with \cle{look after}, the particle never separates from the verb before a pronoun: \emph{She looks \textbf{after them}} and never \emph{looks them after}. The same rule applies to \emph{stand up for} and \emph{take part in}: \emph{We stand up \textbf{for them}}, \emph{We take part \textbf{in it}}.
\end{attention}

\textbf{Expressing purpose: to, in order to, so as to}\par

Why do people volunteer? To answer, English uses the infinitive of purpose.

\begin{aretenir}[Expressing purpose]
To say \emph{why} someone does something, use:
\begin{itemize}
\item \textbf{to + infinitive}: \emph{She joined the NGO \textbf{to promote} girls' education.}
\item \textbf{in order to + infinitive} (more formal): \emph{They set up a club \textbf{in order to fight} discrimination.}
\item \textbf{so as to + infinitive} (formal): \emph{He gave out leaflets \textbf{so as to inform} the public.}
\end{itemize}
Negative purpose: \emph{in order \textbf{not} to} / \emph{so as \textbf{not} to}: \emph{She spoke quietly so as not to frighten the children.}
\end{aretenir}

\begin{exemplebox}[Phrasal verbs + purpose together]
\emph{The students \textbf{set up} a gender equality club \textbf{in order to bring about} change in their school. They \textbf{gave out} leaflets \textbf{to inform} parents, and they \textbf{carried out} a survey \textbf{to understand} why some girls drop out of school.}
\end{exemplebox}

\courssub{Sequence of tenses}

\textbf{The intuition}\par

Imagine Amina says on Monday: \emph{``I \textbf{will} volunteer at the centre.''} On Friday, you report her words: \emph{Amina said she \textbf{would} volunteer at the centre.} When the reporting verb is in the past (\emph{said}), the verbs inside the reported sentence ``step back'' into the past. This is the \cle{sequence of tenses} (la concordance des temps).

\begin{propriete}[Sequence of tenses after a past reporting verb]
When the introductory verb is in the past (\emph{said, told, explained, believed}), the tenses in the subordinate clause shift back:
\begin{itemize}
\item present simple $\rightarrow$ past simple: \emph{``I work hard''} $\rightarrow$ \emph{She said she \textbf{worked} hard.}
\item will $\rightarrow$ would: \emph{``I \textbf{will} volunteer''} $\rightarrow$ \emph{She said she \textbf{would} volunteer.}
\item can $\rightarrow$ could: \emph{``We \textbf{can} help''} $\rightarrow$ \emph{He said they \textbf{could} help.}
\item present perfect / past simple $\rightarrow$ past perfect: \emph{``We \textbf{have organised} a campaign''} $\rightarrow$ \emph{They said they \textbf{had organised} a campaign.}
\end{itemize}
\end{propriete}

\begin{popfigure}
\begin{tikzpicture}[>=Stealth]
\node[draw,rounded corners,fill=popgreenL,text width=3.2cm,align=center] (a1) at (0,0) {\emph{will}\\ (direct speech)};
\node[draw,rounded corners,fill=poporangeL,text width=3.2cm,align=center] (b1) at (7,0) {\emph{would}\\ (reported speech)};
\draw[->,very thick,popblue] (a1) -- (b1) node[midway,above]{\small past reporting verb};
\node[draw,rounded corners,fill=popgreenL,text width=3.2cm,align=center] (a2) at (0,-1.8) {\emph{present simple}\\ \emph{``she works''}};
\node[draw,rounded corners,fill=poporangeL,text width=3.2cm,align=center] (b2) at (7,-1.8) {\emph{past simple}\\ \emph{she worked}};
\draw[->,very thick,popblue] (a2) -- (b2);
\node[draw,rounded corners,fill=popgreenL,text width=3.2cm,align=center] (a3) at (0,-3.6) {\emph{past simple / present perfect}\\ \emph{``she organised''}};
\node[draw,rounded corners,fill=poporangeL,text width=3.2cm,align=center] (b3) at (7,-3.6) {\emph{past perfect}\\ \emph{she had organised}};
\draw[->,very thick,popblue] (a3) -- (b3);
\node[draw,rounded corners,fill=popgreenL,text width=3.2cm,align=center] (a4) at (0,-5.4) {\emph{can}\\ \emph{``we can vote''}};
\node[draw,rounded corners,fill=poporangeL,text width=3.2cm,align=center] (b4) at (7,-5.4) {\emph{could}\\ \emph{they could vote}};
\draw[->,very thick,popblue] (a4) -- (b4);
\end{tikzpicture}
\legende{The sequence of tenses: after a reporting verb in the past, every tense ``steps back'' one step into the past.}
\end{popfigure}

\begin{exoresolu}[Applying the sequence of tenses]
Report this sentence: \emph{The coordinator said: ``Our volunteers have set up three clubs and they will train more girls next year.''}
\tcblower
\textbf{Step 1} — Identify the reporting verb: \emph{said} (past). So all tenses shift back.\\
\textbf{Step 2} — \emph{have set up} (present perfect) $\rightarrow$ \emph{had set up} (past perfect); \emph{will train} $\rightarrow$ \emph{would train}.\\
\textbf{Solution}: \emph{The coordinator said that their volunteers \textbf{had set up} three clubs and that they \textbf{would train} more girls the following year.}
\end{exoresolu}

\begin{attention}[When the tense does NOT change]
If the reported sentence expresses a permanent truth or a fact that is still true, the tense can stay the same: \emph{The teacher said that girls \textbf{have} the same rights as boys.} Also, if the reporting verb is in the present (\emph{she says}), there is no shift at all: \emph{She says she \textbf{will} volunteer.}
\end{attention}

\courssub{Impersonal passives}

\textbf{The intuition}\par

Journalists and official reports often avoid saying \emph{who} says or believes something. Instead of \emph{``People say that volunteers change lives,''} they write \emph{``It is said that volunteers change lives.''} This structure sounds more objective and more formal — perfect for reports, news and exam essays.

\begin{definition}[Impersonal passive]
An \cle{impersonal passive} reports what people in general say, believe, know or report, without naming them. The most common verbs are: \emph{say, believe, report, know, think, expect}.
\end{definition}

\begin{methode}[Building an impersonal passive: two possible structures]
Starting from: \emph{People say that volunteers change lives.}
\begin{enumerate}
\item \textbf{Structure 1} — \emph{It + is said + that + clause}: \emph{\textbf{It is said that} volunteers change lives.}
\item \textbf{Structure 2} — \emph{Subject + is said + to-infinitive}: \emph{Volunteers \textbf{are said to} change lives.}
\end{enumerate}
In structure 2, the subject of the \emph{that}-clause becomes the subject of the passive sentence, and the verb becomes a to-infinitive.
\end{methode}

\begin{exemplebox}[Impersonal passives in context]
\emph{\textbf{It is reported that} more girls are enrolled in schools in the Far North thanks to awareness campaigns.}\\
\emph{\textbf{It is believed that} education brings about lasting change.}\\
\emph{Volunteers \textbf{are known to} look after vulnerable people in many communities.}
\end{exemplebox}

\begin{exoresolu}[Active $\rightarrow$ impersonal passive]
Transform: \emph{People believe that the campaign has changed attitudes.}
\tcblower
\textbf{Structure 1}: \emph{\textbf{It is believed that} the campaign has changed attitudes.}\\
\textbf{Structure 2}: \emph{The campaign \textbf{is believed to have} changed attitudes.} (Note: the action happened before the belief, so we use the perfect infinitive \emph{to have changed}.)
\end{exoresolu}

\courssub{Writing: a composition on volunteering in gender equality promotion}

\textbf{Task and plan}\par

\textbf{Exam task}: \emph{Write a composition of 250--300 words on the following topic: ``Volunteering in gender equality promotion activities in my community.''}

\begin{methode}[Building your composition in four steps]
\begin{enumerate}
\item \textbf{Brainstorm}: list ideas — activities (campaigns, clubs, leaflets, door-to-door visits), people (students, NGOs, traditional leaders), results (more girls in school, changed attitudes).
\item \textbf{Plan}: introduction (present the situation and the topic), development in 2--3 paragraphs (activities / people involved / results and difficulties), conclusion (your opinion and hopes).
\item \textbf{Write}: link your paragraphs with connectors: \emph{firstly, moreover, however, as a result, in conclusion}. Reuse the grammar of this lesson: phrasal verbs, purpose with \emph{to/in order to}, sequence of tenses, impersonal passives.
\item \textbf{Check}: word count (250--300), verb tenses, spelling, punctuation.
\end{enumerate}
\end{methode}

\textbf{Model composition}\par

\begin{exemplebox}[Model answer (about 280 words)]
\textbf{Volunteering in gender equality promotion in my community}

In my community, many people still believe that boys are more important than girls. However, things are slowly changing, thanks to the work of volunteers. It is said that volunteering can bring about real change, and I have seen this with my own eyes.

Firstly, young volunteers in my town have set up a gender equality club in order to fight discrimination. Every Saturday, they carry out awareness campaigns in schools and markets. They give out leaflets to inform parents about the importance of sending girls to school. Moreover, some volunteers look after girls who have dropped out of school and help them to return to the classroom.

The results are encouraging. It is reported that more girls are now enrolled in our schools than five years ago. Last year, our club organised a public debate, and many traditional leaders attended it. The organiser said that the leaders had promised to stand up for girls' rights. She added that they would support future campaigns.

However, there are still difficulties. Some families refuse to listen, and volunteers are sometimes discouraged. Nevertheless, I believe that we must not give up. Education is the key to development, and girls have the same right to education as boys.

In conclusion, volunteering in gender equality promotion is a noble activity that transforms lives. I am proud to take part in it, and I encourage every young Cameroonian to join us, so as to build a fairer society for all.
\end{exemplebox}

\begin{attention}[Exam marking criteria]
At the Probatoire and the Baccalauréat technique, your composition is marked on four criteria: \textbf{content} (relevant ideas), \textbf{organisation} (clear paragraphs, connectors), \textbf{vocabulary} (rich, correct words from the unit) and \textbf{grammatical accuracy} (correct tenses, phrasal verbs, passives). A beautiful idea with many grammar mistakes loses marks: always leave five minutes to proofread.
\end{attention}

\begin{savaistu}[Did you know?]
In Cameroon, many NGOs and school clubs promote gender equality through volunteering. Programmes supporting girls' education, especially in the Far North, North and Adamawa regions, have helped thousands of girls stay in school. International days such as the International Day of the Girl Child (11 October) are marked by awareness campaigns organised largely by young volunteers — a perfect topic for your exam essays!
\end{savaistu}

\courssub{Practice exercises}

\begin{exercice}[Exercise 1: Phrasal verbs]
Complete the sentences with the correct phrasal verb: \emph{carry out, set up, look after, give out, bring about, stand up for}.
\begin{enumerate}
\item The students decided to \dots{} a club to defend girls' rights.
\item Volunteers \dots{} leaflets in the market every Friday.
\item We must \dots{} the rights of all children.
\item She \dots{} her younger brothers and sisters after school.
\item The NGO will \dots{} a survey in three villages.
\item Hard work can \dots{} real change in society.
\end{enumerate}
\end{exercice}

\begin{exercice}[Exercise 2: Sequence of tenses]
Report these sentences, starting with \emph{She said that\dots}
\begin{enumerate}
\item ``I will join the volunteer club.''
\item ``We have organised two campaigns.''
\item ``The girls can study at the centre.''
\item ``My brother works for an NGO.''
\end{enumerate}
\end{exercice}

\begin{exercice}[Exercise 3: Impersonal passives]
Transform each sentence into the two impersonal passive structures.
\begin{enumerate}
\item People say that the campaign is a success.
\item People believe that education changes attitudes.
\end{enumerate}
\end{exercice}

\begin{corrige}[Exercise 1]
1. set up \quad 2. give out \quad 3. stand up for \quad 4. looks after \quad 5. carry out \quad 6. bring about.
\end{corrige}

\begin{corrige}[Exercise 2]
1. She said that she \textbf{would join} the volunteer club. 2. She said that they \textbf{had organised} two campaigns. 3. She said that the girls \textbf{could study} at the centre. 4. She said that her brother \textbf{worked} for an NGO.
\end{corrige}

\begin{corrige}[Exercise 3]
1. \emph{It is said that the campaign is a success.} / \emph{The campaign is said to be a success.} 2. \emph{It is believed that education changes attitudes.} / \emph{Education is believed to change attitudes.}
\end{corrige}

\begin{aretenir}[Key points of the section]
\begin{itemize}
\item A phrasal verb = verb + particle, with a new global meaning (\emph{carry out, set up, look after, give out, bring about, stand up for, take part in}).
\item Purpose: \emph{to / in order to / so as to + infinitive}.
\item Sequence of tenses: past reporting verb $\rightarrow$ \emph{will} becomes \emph{would}, \emph{can} becomes \emph{could}, present perfect becomes past perfect.
\item Impersonal passive: \emph{It is said that + clause} or \emph{Subject + is said + to-infinitive}.
\item A good composition = clear plan + connectors + unit grammar + careful proofreading (250--300 words).
\end{itemize}
\end{aretenir}

\newpage

\coursec{Activité d'intégration}

\coursec{Citizenship and human rights}

\courssub{Objectifs de l'activité}
\begin{itemize}
    \item Mobiliser l'ensemble des savoirs de la leçon \cle{Citizenship and human rights} : vocabulaire électoral et de volontariat, questions rapportées, ordre des mots, \emph{phrasal verbs}, concordance des temps, passifs impersonnels.
    \item Produire une tâche complexe authentique combinant production orale (présentation de sondage) et production écrite (rapport d'enquête, composition personnelle).
    \item Développer des compétences transversales : travail collaboratif, enquête de terrain, traitement de données chiffrées, argumentation citoyenne.
    \item S'auto-évaluer et évaluer ses pairs à l'aide d'une grille critériée commune (contenu, organisation, vocabulaire, grammaire).
\end{itemize}

\courssub{Situation}
\textbf{Contexte :} Le Lycée Technique de New Bell (Douala) organise l'élection de ses nouveaux délégués de classe pour l'année scolaire 2026-2027. L'administration encourage les élèves à mener une campagne citoyenne responsable, fondée sur l'écoute des attentes réelles de la communauté scolaire.

\textbf{Votre mission :} En groupes de 4 élèves, vous pilotez le projet \emph{« My school, my voice »}. Vous devez :
\begin{enumerate}
    \item Concevoir un mini-sondage de 5 questions fermées (choix multiples ou échelle de Likert) portant sur les qualités attendues d'un délégué, les priorités d'action (hygiène, discipline, activités culturelles, orientation, égalité filles-garçons) et le mode de scrutin souhaité.
    \item Interroger 20 élèves du lycée (hors votre classe) sur une période de deux jours.
    \item Présenter les résultats sous forme de diagramme en barres (une barre par modalité de réponse par question) et les commenter oralement en 3 minutes devant la classe (\emph{« According to our survey, 55\% of respondents consider honesty as the top quality... »}).
    \item Rédiger collectivement un rapport écrit d'une page (environ 150 mots) rapportant les réponses significatives à l'aide de \textbf{questions rapportées} (\emph{« One student said that... Another asked whether... »}).
    \item Rédiger \textbf{individuellement} une composition de 250 à 300 mots intitulée \emph{« How I would volunteer to promote gender equality in my school »}. Le texte doit impérativement contenir :
    \begin{itemize}
        \item au moins \textbf{4 \emph{phrasal verbs}} (ex. \cle{stand up for}, \cle{team up with}, \cle{set up}, \cle{carry out}, \cle{speak out}, \cle{look into}, \cle{run for});
        \item au moins \textbf{2 passifs impersonnels} (ex. \cle{It is believed that...}, \cle{Students are encouraged to...}, \cle{It has been decided that...});
        \item une \textbf{concordance des temps} correcte tout au long du récit (passé simple / passé continu / plus-que-parfait / futur dans le passé).
    \end{itemize}
\end{enumerate}
Une \textbf{grille d'évaluation commune} (4 critères × 5 points = 20 points) guide la correction et la remédiation.

\courssub{Consignes}
\begin{enumerate}
    \item \textbf{Étape 1 — Conception du sondage (20 min) :} En groupe, rédigez 5 questions claires, neutres et exploitables statistiquement. Validez-les avec l'enseignant avant diffusion. \emph{(Vocabulaire : \cle{poll}, \cle{survey}, \cle{respondent}, \cle{ballot}, \cle{manifesto}, \cle{gender equality})}
    \item \textbf{Étape 2 — Collecte des données (hors classe, 2 jours) :} Chaque membre du groupe interroge 5 élèves différents. Notez les réponses brutes dans un tableau de collecte.
    \item \textbf{Étape 3 — Traitement et présentation orale (30 min en classe) :} Compilez les 20 réponses par question. Construisez un diagramme en barres (papier ou tableur). Préparez un commentaire oral structuré : introduction (objectif, échantillon), résultats principaux (pourcentages), analyse brève, conclusion. Utilisez l'ordre des mots correct (SVOMPT) et le vocabulaire électoral.
    \item \textbf{Étape 4 — Rapport écrit d'enquête (groupe, 30 min) :} Rédigez le rapport en transformant les citations directes en \textbf{questions rapportées} (\emph{wh-questions} et \emph{yes/no questions}). Respectez la concordance des temps (\emph{present $\to$ past, will $\to$ would, can $\to$ could}).
    \item \textbf{Étape 5 — Composition individuelle (devoir maison, 1 semaine) :} Rédigez votre composition selon les contraintes grammaticales imposées. Relisez-vous à l'aide de la checklist fournie.
    \item \textbf{Étape 6 — Évaluation et remédiation (1 séance) :} Échangez les productions, appliquez la grille, discutez des écarts, l'enseignant valide et propose des exercices de remédiation ciblés.
\end{enumerate}

\courssub{Ressources à mobiliser}
\begin{definition}[Questions rapportées — Règles clés]
\textbf{Wh-questions :} \emph{« What are your priorities? » $\to$ He asked what my priorities \textbf{were}.} (Changement de temps, ordre affirmatif S+V, pas d'auxiliaire \emph{do/does/did}).
\textbf{Yes/No questions :} \emph{« Will you vote for her? » $\to$ She asked \textbf{if/whether} I \textbf{would} vote for her.}
\textbf{Concordance des temps :} Present $\to$ Past ; Past $\to$ Past Perfect ; Will $\to$ Would ; Can $\to$ Could ; Must $\to$ Had to.
\end{definition}

\begin{definition}[Passif impersonnel — Deux structures]
\begin{itemize}
    \item \textbf{It + passive verb + that-clause :} \emph{It is said that... / It is believed that... / It has been decided that... / It is expected that...}
    \item \textbf{Subject + passive verb + to-infinitive :} \emph{He is said to be... / Students are encouraged to participate... / The delegates are expected to report...}
\end{itemize}
\end{definition}

\begin{definition}[Phrasal verbs utiles pour le volontariat]
\begin{center}
\begin{tabularx}{\linewidth}{|l|X|X|}\hline
\rowcolor{popblueL} \textbf{Phrasal verb} & \textbf{Sens (FR)} & \textbf{Exemple de contexte} \\ \hline
\cle{stand up for} & défendre (une cause) & I will \emph{stand up for} girls' rights. \\ \hline
\cle{team up with} & faire équipe avec & We \emph{teamed up with} the student council. \\ \hline
\cle{set up} & mettre en place, organiser & We \emph{set up} a mentoring programme. \\ \hline
\cle{carry out} & réaliser, mener à bien & We \emph{carried out} a survey last week. \\ \hline
\cle{speak out} & s'exprimer publiquement & She \emph{spoke out} against discrimination. \\ \hline
\cle{look into} & examiner, enquêter sur & The committee will \emph{look into} the complaints. \\ \hline
\cle{run for} & se présenter, briguer (un poste) & She decided to \emph{run for} class delegate. \\ \hline
\end{tabularx}
\end{center}
\end{definition}

\begin{propriete}[Ordre des mots (SVOMPT) — Rappel]
Sujet + Verbe + Objet + Manière + Place + Temps. \\
Ex. : \emph{The students (S) elected (V) their delegate (O) democratically (M) in the hall (P) yesterday (T).}
\end{propriete}

\courssub{Démarche et corrigé commenté}
\begin{exoresolu}{Projet « My school, my voice » — Production modèle guidée}
\textbf{Étape 1 — Sondage modèle (validé par l'enseignant)}
\begin{center}
\begin{tabularx}{\linewidth}{|l|X|}\hline
\rowcolor{popblueL} \textbf{Q\#} & \textbf{Question (closed-ended)} \\ \hline
1 & Which quality is most important for a class delegate? \newline A) Honesty \quad B) Popularity \quad C) Intelligence \quad D) Boldness \\ \hline
2 & What should be the delegate's first priority? \newline A) Hygiene \quad B) Discipline \quad C) Cultural activities \quad D) Gender equality \\ \hline
3 & How should the vote be organised? \newline A) Secret ballot \quad B) Show of hands \quad C) Online poll \quad D) Consensus \\ \hline
4 & Should gender parity be respected in the student council? \newline A) Yes, absolutely \quad B) Preferably \quad C) Not necessary \quad D) No opinion \\ \hline
5 & Would you volunteer for a gender-equality club? \newline A) Yes, certainly \quad B) Maybe \quad C) Probably not \quad D) Never \\ \hline
\end{tabularx}
\end{center}

\textbf{Étape 2 — Données brutes compilées (20 répondants)}
\begin{center}
\begin{tabular}{|c|c|c|c|c|c|}\hline
\rowcolor{popblueL} \textbf{Question} & \textbf{Option A} & \textbf{Option B} & \textbf{Option C} & \textbf{Option D} & \textbf{Total} \\ \hline
1 (Quality) & 11 & 3 & 4 & 2 & 20 \\ \hline
2 (Priority) & 5 & 6 & 3 & 6 & 20 \\ \hline
3 (Vote method) & 14 & 2 & 3 & 1 & 20 \\ \hline
4 (Gender parity) & 12 & 5 & 2 & 1 & 20 \\ \hline
5 (Volunteering) & 8 & 7 & 3 & 2 & 20 \\ \hline
\end{tabular}
\end{center}

\textbf{Étape 3 — Commentaire oral modèle (3 min)}
\emph{« Good morning everyone. Our group conducted a survey among 20 students of Lycée Technique de New Bell about the upcoming election of class delegates. According to our survey, \textbf{55\%} of respondents consider \textbf{honesty} as the top quality for a delegate, while only 15\% chose popularity. Regarding priorities, the votes are split: 30\% for discipline, 30\% for gender equality, 25\% for hygiene, and 15\% for cultural activities. Interestingly, \textbf{70\%} prefer a \textbf{secret ballot}, which shows a strong attachment to democratic principles. On gender parity, 60\% say it is absolutely necessary and 25\% say preferably. Finally, 40\% would certainly volunteer for a gender-equality club and 35\% maybe. To conclude, students want honest delegates, democratic voting, and are ready to engage for equality. Thank you. »}

\textbf{Étape 4 — Rapport écrit d'enquête (groupe) — Questions rapportées}
\emph{One student \textbf{said that} honesty was the most important quality because delegates manage money. Another \textbf{asked whether} the election would be transparent. A third respondent \textbf{wanted to know} what the delegate would do for gender equality. Several students \textbf{stated that} they \textbf{would} vote only if the ballot was secret. Many \textbf{added that} gender parity \textbf{had to} be respected in the council. One girl \textbf{explained that} she \textbf{could} volunteer if the club was well organised. The group \textbf{concluded that} the majority \textbf{were} ready to \textbf{stand up for} equality.}

\textbf{Analyse grammaticale du rapport :}
\begin{itemize}
    \item \emph{said that honesty was} (Present $\to$ Past Simple) — \checkmark
    \item \emph{asked whether the election would be} (Will $\to$ Would) — \checkmark
    \item \emph{wanted to know what the delegate would do} (Will $\to$ Would, ordre affirmatif) — \checkmark
    \item \emph{stated that they would vote} (Will $\to$ Would) — \checkmark
    \item \emph{added that gender parity had to be respected} (Must $\to$ Had to + Passif) — \checkmark
    \item \emph{explained that she could volunteer} (Can $\to$ Could) — \checkmark
    \item \emph{concluded that the majority were ready} (Are $\to$ Were) — \checkmark
\end{itemize}

\textbf{Étape 5 — Composition individuelle modèle (298 mots)}
\textbf{Title: How I would volunteer to promote gender equality in my school}

\emph{If I \textbf{were} elected as a volunteer coordinator for gender equality, I \textbf{would} immediately \textbf{set up} a mixed committee of boys and girls. \textbf{It is believed that} real change starts when students \textbf{team up with} each other across gender lines. Last year, I \textbf{noticed} that girls \textbf{were often ignored} during class debates. \textbf{It has been decided that} this year, a « Girls' Voice » week \textbf{will be organised} every term.}

\emph{First, I \textbf{would speak out} against stereotypes in morning assemblies. I \textbf{would} also \textbf{look into} the reasons why few girls \textbf{run for} leadership positions. A survey \textbf{carried out} last month \textbf{showed} that 60\% of girls \textbf{lacked} confidence. Therefore, I \textbf{would} \textbf{team up with} the school counsellor to \textbf{carry out} confidence-building workshops.}

\emph{Second, I \textbf{would stand up for} equal access to technical workshops. In our section (Électrotechnique), girls \textbf{are encouraged to} handle tools exactly like boys. \textbf{It is expected that} this practice \textbf{will inspire} younger students. I \textbf{would} also \textbf{set up} a mentoring system where senior girls \textbf{guide} junior ones.}

\emph{Finally, I \textbf{would} write a termly report for the administration. \textbf{It is said that} accountability \textbf{ensures} sustainability. By the end of my mandate, I \textbf{hope} the school \textbf{will have become} a model of parity in Douala.}

\textbf{Vérification des contraintes :}
\begin{itemize}
    \item \textbf{4+ Phrasal verbs :} \cle{set up} (x2), \cle{team up with} (x2), \cle{speak out}, \cle{look into}, \cle{carry out} (x2), \cle{stand up for}, \cle{run for} = \textbf{10 occurrences} — \checkmark
    \item \textbf{2+ Passifs impersonnels :} \cle{It is believed that}, \cle{It has been decided that}, \cle{It is expected that}, \cle{It is said that} = \textbf{4 occurrences} (structure \emph{It + passive + that}) ; \cle{girls are encouraged to}, \cle{delegates are expected to} = \textbf{2 occurrences} (structure \emph{Subject + passive + to-inf}) — \checkmark
    \item \textbf{Concordance des temps :} \emph{were/would (2nd conditional), noticed/were ignored (past/past passive), has been decided/will be organised (present perfect/future passive), would speak/would team/would look/would run/would carry out/would stand/would set up/would guide/would write/will have become (conditionnel/futur antérieur dans le passé)} — \checkmark
    \item \textbf{Mots :} $\approx$ 298 mots — \checkmark
\end{itemize}
\end{exoresolu}

\courssub{Bilan de l'activité}
Cette activité d'intégration a permis de vérifier l'acquisition opérationnelle des savoirs de l'unité \cle{Citizenship and human rights} dans une situation authentique et motivante pour des élèves de Terminale technique.

\begin{aretenir}[Compétences validées]
\begin{itemize}
    \item \textbf{Oral :} Présenter des données chiffrées (pourcentages, tendances) avec un vocabulaire électoral précis (\cle{ballot}, \cle{manifesto}, \cle{constituency}, \cle{turnout}) et une prononciation intelligible.
    \item \textbf{Grammaire en contexte :} Transformer spontanément le discours direct en \textbf{questions rapportées} (choix correct \emph{if/whether}, concordance des temps, ordre affirmatif) ; utiliser l'\textbf{ordre des mots} SVOMPT dans des phrases longues ; employer \textbf{4+ \emph{phrasal verbs}} pertinents ; construire \textbf{2+ passifs impersonnels} variés (\emph{It is... that / Subject + be + V-en + to-inf}) ; maintenir la \textbf{concordance des temps} sur un texte long.
    \item \textbf{Écrit :} Structurer une composition argumentative (introduction, développement par arguments, conclusion) avec connecteurs logiques (\emph{First, Second, Therefore, Finally}).
    \item \textbf{Méthodologie :} Concevoir un outil de collecte, traiter un échantillon (n=20), produire une visualisation (diagramme en barres), travailler en équipe, respecter un cahier des charges contraignant.
\end{itemize}
\end{aretenir}

\begin{attention}[Points de vigilance pour la remédiation]
\begin{itemize}
    \item \textbf{Erreur fréquente 1 :} Oubli du \emph{backshift} dans le rapporté (\emph{*He asked what are my priorities} au lieu de \emph{what my priorities were}). $\to$ Exercice ciblé : transformation de 10 citations directes.
    \item \textbf{Erreur fréquente 2 :} Mauvaise formation du passif impersonnel (\emph{*It is believed students are...} sans \emph{that} ; \emph{*Students are believed to are...}). $\to$ Tableau de substitution \emph{It is said / thought / expected / decided / reported that...} + \emph{Subject + is/are said / thought / expected + to-inf}.
    \item \textbf{Erreur fréquente 3 :} Confusion \emph{phrasal verb} / verbe + préposition (\emph{*stand for} au lieu de \emph{stand up for}). $\to$ Liste à apprendre par cœur + exercices de complétion.
    \item \textbf{Erreur fréquente 4 :} Diagramme en barres mal étiqueté (pas de titre, pas d'échelle, pas d'unité \%). $\to$ Modèle imposé pour la prochaine évaluation.
\end{itemize}
\end{attention}

\begin{methode}[Grille d'évaluation commune (20 pts) — À distribuer aux élèves]
\begin{center}
\begin{tabularx}{\linewidth}{|l|X|c|}\hline
\rowcolor{popblueL} \textbf{Critère} & \textbf{Indicateurs de réussite} & \textbf{Points} \\ \hline
Contenu (5) & Répond à la consigne, idées pertinentes, données exactes, arguments développés, longueur respectée & /5 \\ \hline
Organisation (5) & Structure claire (intro/dev/concl), paragraphes, connecteurs logiques, progression cohérente & /5 \\ \hline
Vocabulaire (5) & Vocabulaire électoral, volontariat, genre ; \emph{phrasal verbs} (4+), variété, registres adaptés & /5 \\ \hline
Grammaire (5) & Questions rapportées correctes, passifs impersonnels (2+), concordance des temps, ordre des mots, peu d'erreurs bloquantes & /5 \\ \hline
\textbf{Total} &  & \textbf{/20} \\ \hline
\end{tabularx}
\end{center}
\textit{Seuil de réussite : 10/20. Remédiation obligatoire si note $<$ 10 ou si un critère $<$ 2/5.}
\end{methode}

\begin{savaistu}[Ancrage camerounais]
Au Cameroun, la loi n°2017/013 du 12 juillet 2017 relative au code électoral et les textes régissant la vie scolaire dans les établissements publics (décret n°2001/041 et arrêtés d'application) prévoient la participation des élèves par l'élection de délégués. Le MINEDUB et le MINESEC encouragent la mise en place de \emph{Clubs Genre} dans les lycées techniques pour lutter contre les stéréotypes dans les filières industrielles (Électrotechnique, Mécanique, BTP) où les filles restent minoritaires (\textasciitilde 15-20\% des effectifs). Le projet « My school, my voice » s'inscrit dans cette dynamique citoyenne nationale.
\end{savaistu}

\newpage

\coursec{Méthodes et exercices résolus}

\coursec{Citizenship and Human Rights}

\courssub{Speaking and Vocabulary: Campaigning and Voting for School Leaders}

\begin{definition}[Key Vocabulary -- Campaigning and Voting]
\begin{itemize}
\item \cle{to campaign} (faire campagne) -- to work in an organised way to achieve a goal, e.g. \emph{to campaign for school president}
\item \cle{candidate} (candidat) -- a person who stands for election
\item \cle{manifesto} (programme électoral) -- a written statement of beliefs and plans
\item \cle{polling station} (bureau de vote) -- the place where people vote
\item \cle{ballot paper} (bulletin de vote) -- the paper used to record a vote
\item \cle{ballot box} (urne) -- the container for ballot papers
\item \cle{to cast a vote} (voter) -- to formally express a choice
\item \cle{turnout} (taux de participation) -- the number of people who vote
\item \cle{to rally support} (mobiliser du soutien) -- to gather people behind a cause
\item \cle{slogan} (slogan) -- a short, striking phrase used in a campaign
\end{itemize}
\end{definition}

\begin{exemplebox}[Dialogue: Campaign Interview]
\textbf{Journalist:} Good morning. You are a candidate for school leader. \textbf{Why did you decide to stand for election?} \\
\textbf{Candidate:} I want to improve our study conditions. \textbf{Will you organise a debate before the vote?} \\
\textbf{Journalist:} Yes, next Tuesday. \textbf{How many students have registered to vote today?} \\
\textbf{Candidate:} About 300 so far.
\end{exemplebox}

\begin{exercice}[Speaking Practice -- Role Play]
Work in pairs. Student A is a journalist, Student B is a candidate. Use the vocabulary above. Ask and answer at least five questions about the campaign, the manifesto, and voting day. Then swap roles.
\end{exercice}

\courssub{Grammar Focus 1: Reported Questions}

\begin{methode}[Reporting Questions Correctly]
To transform a direct question into a reported question (\emph{question au style indirect}), follow these steps:
\begin{enumerate}
\item Choose a reporting verb: \cle{asked}, \cle{wanted to know}, \cle{wondered}, \cle{enquired}.
\item If the question is a \textbf{yes/no question}, introduce it with \cle{if} or \cle{whether}.
\item If it is a \textbf{wh-question}, keep the question word (\emph{what, where, why, how}\ldots).
\item Change the question order into \textbf{statement order}: subject + verb. Never keep the inversion.
\item Backshift the tense: present simple $\rightarrow$ past simple; present continuous $\rightarrow$ past continuous; present perfect / past simple $\rightarrow$ past perfect; \emph{will} $\rightarrow$ \emph{would}; \emph{can} $\rightarrow$ \emph{could}.
\item Adapt time and place expressions: \emph{now} $\rightarrow$ \emph{then}; \emph{today} $\rightarrow$ \emph{that day}; \emph{tomorrow} $\rightarrow$ \emph{the next day}; \emph{here} $\rightarrow$ \emph{there}.
\item Remove the question mark: a reported question ends with a \textbf{full stop}.
\end{enumerate}
\end{methode}

\begin{exoresolu}[Reported Questions -- Bac-Type Exercise]
\textbf{Task.} During the campaign for school leader elections, a journalist interviewed candidates. Report these questions, beginning as shown.
\begin{enumerate}
\item ``Will you organise a debate before the vote?'' The journalist asked the candidate \ldots
\item ``Why did you decide to stand for election?'' She asked him \ldots
\item ``How many students have registered to vote today?'' She wanted to know \ldots
\end{enumerate}
\tcblower
\textbf{Solution.}
\begin{enumerate}
\item Yes/no question with \emph{will} $\rightarrow$ use \cle{if/whether} and backshift \emph{will} to \emph{would}: \\
The journalist asked the candidate \textbf{if (whether) he would organise a debate before the vote}.
\item Wh-question with \emph{why}; past simple \emph{did you decide} $\rightarrow$ past perfect \emph{he had decided}; statement order: \\
She asked him \textbf{why he had decided to stand for election}.
\item Wh-question; present perfect \emph{have registered} $\rightarrow$ past perfect \emph{had registered}; \emph{today} $\rightarrow$ \emph{that day}: \\
She wanted to know \textbf{how many students had registered to vote that day}.
\end{enumerate}
\textbf{Conclusion:} in every reported question, the word order is subject + verb, the tense is backshifted, and there is no question mark.
\end{exoresolu}

\begin{exercice}[Reported Questions -- Practice]
Report the following questions.
\begin{enumerate}
\item ``Can you explain your manifesto?'' The student asked \ldots
\item ``Where is the polling station?'' He wanted to know \ldots
\item ``Did you vote yesterday?'' She asked me \ldots
\item ``What time does the voting start?'' They enquired \ldots
\end{enumerate}
\end{exercice}

\begin{corrige}[Exercice Reported Questions]
\begin{enumerate}
\item The student asked \textbf{if/whether I could explain my manifesto}.
\item He wanted to know \textbf{where the polling station was}.
\item She asked me \textbf{if/whether I had voted the day before}.
\item They enquired \textbf{what time the voting started}.
\end{enumerate}
\end{corrige}

\courssub{Listening and Vocabulary: Training Courses on Democracy}

\begin{definition}[Key Vocabulary -- Democracy Training]
\begin{itemize}
\item \cle{workshop} (atelier) -- a practical training session
\item \cle{facilitator} (formateur, animateur) -- the person leading the training
\item \cle{participant} (participant) -- a person taking part
\item \cle{civic education} (éducation civique) -- teaching about rights and duties of citizens
\item \cle{ballot} (scrutin) -- a secret vote
\item \cle{transparency} (transparence) -- openness, accountability
\item \cle{accountability} (reddition de comptes) -- obligation to explain actions
\item \cle{to empower} (donner du pouvoir à, habiliter) -- to give someone the authority or confidence to act
\item \cle{grassroots} (de base, local) -- at the local community level
\item \cle{consensus} (consensus) -- general agreement
\end{itemize}
\end{definition}

\begin{exemplebox}[Listening Script: Democracy Workshop Announcement]
\textbf{Announcer:} ``Good morning everyone. Welcome to the three-day workshop on \emph{Democratic Participation in Schools}. Our facilitator, Mr. Ngwa, will \textbf{carry out} sessions on voting procedures, transparency, and student leadership. Participants will \textbf{take part in} group discussions and role plays. The aim is to \textbf{empower} you to organise free and fair elections in your schools. Please \textbf{sign up} at the desk before 9 a.m. tomorrow.''
\end{exemplebox}

\begin{exercice}[Listening Comprehension]
Listen to the announcement (or read the script) and answer:
\begin{enumerate}
\item What is the theme of the workshop?
\item Who is the facilitator?
\item Name three activities participants will do.
\item What is the overall aim?
\item What must participants do before 9 a.m. tomorrow?
\end{enumerate}
\end{exercice}

\begin{corrige}[Exercice Listening Comprehension]
\begin{enumerate}
\item Democratic Participation in Schools.
\item Mr. Ngwa.
\item Group discussions, role plays, sessions on voting procedures/transparency/student leadership.
\item To empower participants to organise free and fair elections.
\item Sign up at the desk.
\end{enumerate}
\end{corrige}

\courssub{Grammar Focus 2: Word Order}

\begin{methode}[Getting the Word Order Right in English Sentences]
English word order is much stricter than French word order. To build a correct sentence:
\begin{enumerate}
\item Start with the \textbf{subject}, then the \textbf{verb}, then the \textbf{object/complement}: S + V + O.
\item Place \textbf{adverbs of frequency} (\emph{always, often, never, usually}) \emph{before} the main verb but \emph{after} the verb \emph{to be}: \emph{She often votes. / She is always punctual.}
\item Place \textbf{adjectives before nouns}: \emph{a free election}, never ``an election free'' in basic sentences.
\item Order of complements at the end: \textbf{Manner -- Place -- Time} (MPT): \emph{They voted calmly at school on Monday.}
\item In questions, use inversion with the auxiliary: \emph{Do you vote? / Did they campaign?}
\item In reported questions and indirect speech, \textbf{no inversion}: \emph{He asked where the polling station was.}
\end{enumerate}
\end{methode}

\begin{exoresolu}[Word Order -- Bac-Type Exercise]
\textbf{Task.} Reorder the words to make correct sentences.
\begin{enumerate}
\item always / the students / attend / democracy workshops / on Fridays
\item asked / the trainer / where / was / the polling station
\item a / organised / they / successful / campaign / last week / very
\end{enumerate}
\tcblower
\textbf{Solution.}
\begin{enumerate}
\item Adverb of frequency \emph{always} goes before the main verb \emph{attend}; time expression at the end: \\
\textbf{The students always attend democracy workshops on Fridays.}
\item Reported question: no inversion, subject \emph{the polling station} before verb \emph{was}: \\
\textbf{The trainer asked where the polling station was.}
\item Adjective \emph{successful} before the noun \emph{campaign}; adverb \emph{very} modifies the adjective; time at the end: \\
\textbf{They organised a very successful campaign last week.}
\end{enumerate}
\textbf{Conclusion:} respect S + V + O, put frequency adverbs before the main verb, adjectives before nouns, and never invert in reported questions.
\end{exoresolu}

\begin{exercice}[Word Order -- Practice]
Put the words in the correct order.
\begin{enumerate}
\item never / he / misses / a / training / session
\item the / explained / facilitator / how / works / the / ballot / system
\item carefully / the / read / participants / manifesto / the / yesterday
\item often / are / elections / school / in / held / our
\end{enumerate}
\end{exercice}

\begin{corrige}[Exercice Word Order]
\begin{enumerate}
\item He never misses a training session.
\item The facilitator explained how the ballot system works.
\item The participants read the manifesto carefully yesterday.
\item School elections are often held in our school.
\end{enumerate}
\end{corrige}

\courssub{Reading and Vocabulary: Volunteering in Gender Equality Promotion}

\begin{definition}[Key Vocabulary -- Gender Equality Volunteering]
\begin{itemize}
\item \cle{volunteer} (bénévole) -- a person who works without pay
\item \cle{to volunteer} (faire du bénévolat) -- to offer to do something freely
\item \cle{gender equality} (égalité des sexes) -- equal rights and opportunities for all genders
\item \cle{to raise awareness} (sensibiliser) -- to make people conscious of an issue
\item \cle{to empower} (autonomiser, donner du pouvoir à) -- to give power/confidence
\item \cle{discrimination} (discrimination) -- unfair treatment based on gender, etc.
\item \cle{stereotype} (stéréotype) -- a fixed, oversimplified idea
\item \cle{advocacy} (plaidoyer) -- public support for a cause
\item \cle{outreach} (action de terrain, porte-à-porte) -- reaching out to communities
\item \cle{mentorship} (mentorat) -- guidance by an experienced person
\end{itemize}
\end{definition}

\begin{exemplebox}[Reading Text: Girls' Education Campaign]
\textbf{Last year, the Youth Action Club \textbf{set up} a campaign to promote girls' education in rural areas. Volunteers \textbf{carried out} door-to-door visits to \textbf{raise awareness} about the importance of schooling for girls. They \textbf{took part in} a training course on gender-based violence organised by a local NGO. The campaign \textbf{reached} over 500 families. As a result, 50 girls who had dropped out \textbf{re-enrolled} in school. The project \textbf{was praised} by the Ministry of Secondary Education.}
\end{exemplebox}

\begin{exercice}[Reading Comprehension]
Read the text and answer:
\begin{enumerate}
\item What did the Youth Action Club set up?
\item What two main activities did the volunteers do?
\item What training did they take part in?
\item How many families were reached?
\item What was the concrete result for girls?
\item Who praised the project?
\end{enumerate}
\end{exercice}

\begin{corrige}[Exercice Reading Comprehension]
\begin{enumerate}
\item A campaign to promote girls' education in rural areas.
\item Door-to-door visits to raise awareness; took part in a training course on gender-based violence.
\item A training course on gender-based violence organised by a local NGO.
\item Over 500 families.
\item 50 girls who had dropped out re-enrolled in school.
\item The Ministry of Secondary Education.
\end{enumerate}
\end{corrige}

\courssub{Grammar Focus 3: Phrasal Verbs, Sequence of Tenses, Impersonal Passives}

\begin{methode}[Using Phrasal Verbs to Express Purpose]
Phrasal verbs (\emph{verbes à particule}) are frequent in texts about volunteering. To use them well:
\begin{enumerate}
\item Learn the verb with its particle as one unit: \cle{carry out} (mener à bien), \cle{set up} (créer, mettre en place), \cle{take part in} (participer à), \cle{look after} (s'occuper de), \cle{give up} (abandonner), \cle{stand for} (représenter, défendre), \cle{sign up} (s'inscrire), \cle{reach out} (tendre la main, contacter).
\item To express \textbf{purpose}, use \emph{to + infinitive} or \emph{in order to / so as to + infinitive}: \emph{Volunteers set up clubs \textbf{to promote} gender equality.}
\item With a \textbf{separable} phrasal verb, a short object can go between verb and particle: \emph{They carried the project out.} A pronoun \textbf{must} go in the middle: \emph{They carried it out} (never ``carried out it'').
\item With an \textbf{inseparable} phrasal verb (verb + preposition), the object always follows: \emph{She takes part in the campaign.}
\end{enumerate}
\end{methode}

\begin{exoresolu}[Phrasal Verbs and Purpose -- Bac-Type Exercise]
\textbf{Task.} Complete with the correct phrasal verb from the box and add a purpose clause.
\begin{center}
\begin{tabularx}{\linewidth}{|X|X|X|}
\hline
\rowcolor{popblueL} set up & carry out & give up \\
\hline
\end{tabularx}
\end{center}
\begin{enumerate}
\item The association \ldots\ldots\ldots a girls' club \ldots\ldots\ldots (promote equal access to education).
\item The volunteers \ldots\ldots\ldots a survey \ldots\ldots\ldots (know the needs of the community).
\item Even when the work is difficult, we must not \ldots\ldots\ldots .
\end{enumerate}
\tcblower
\textbf{Solution.}
\begin{enumerate}
\item \emph{Set up} means ``créer''; purpose with \emph{to + infinitive}: \\
The association \textbf{set up} a girls' club \textbf{to promote equal access to education}.
\item \emph{Carry out} means ``mener à bien'': \\
The volunteers \textbf{carried out} a survey \textbf{to know the needs of the community}.
\item \emph{Give up} means ``abandonner'', after \emph{must not} the base form is used: \\
Even when the work is difficult, we must not \textbf{give up}.
\end{enumerate}
\textbf{Conclusion:} the particle changes the meaning of the verb; purpose is expressed with \emph{to + infinitive}.
\end{exoresolu}

\begin{exercice}[Phrasal Verbs -- Practice]
Complete the sentences with a phrasal verb from the box in the correct form.
\begin{center}
\begin{tabularx}{\linewidth}{|X|X|X|X|}
\hline
\rowcolor{popblueL} take part in & look after & stand for & sign up \\
\hline
\end{tabularx}
\end{center}
\begin{enumerate}
\item Many students \ldots\ldots\ldots the democracy workshop last week.
\item Our club \ldots\ldots\ldots transparency and fairness in elections.
\item Please \ldots\ldots\ldots at the library before Friday.
\item Volunteers \ldots\ldots\ldots younger students during the campaign.
\end{enumerate}
\end{exercice}

\begin{corrige}[Exercice Phrasal Verbs]
\begin{enumerate}
\item took part in
\item stands for
\item sign up
\item look after
\end{enumerate}
\end{corrige}

\begin{methode}[Mastering the Sequence of Tenses]
When the main verb is in the past, the verbs in the subordinate clauses usually shift backwards:
\begin{enumerate}
\item Identify the tense of the \textbf{reporting/main verb}. If it is in the past (\emph{said, explained, believed}), apply the backshift.
\item Present simple $\rightarrow$ past simple: \emph{``Women work hard'' $\rightarrow$ He said women worked hard.}
\item Present continuous $\rightarrow$ past continuous: \emph{is training} $\rightarrow$ \emph{was training}.
\item Present perfect and past simple $\rightarrow$ past perfect: \emph{has voted / voted} $\rightarrow$ \emph{had voted}.
\item \emph{will} $\rightarrow$ \emph{would}; \emph{can} $\rightarrow$ \emph{could}; \emph{must} $\rightarrow$ \emph{had to}.
\item \textbf{Exception:} a general truth or a fact still true can stay in the present: \emph{The trainer said that every citizen has the right to vote.}
\end{enumerate}
\end{methode}

\begin{exoresolu}[Sequence of Tenses -- Bac-Type Exercise]
\textbf{Task.} Put the verbs in brackets in the correct tense.
\begin{enumerate}
\item The coordinator announced that the training session (start) \ldots the following Monday.
\item She explained that the volunteers (already / visit) \ldots three villages.
\item He told us that gender equality (be) \ldots a human right.
\end{enumerate}
\tcblower
\textbf{Solution.}
\begin{enumerate}
\item Main verb \emph{announced} (past) $\rightarrow$ \emph{will} becomes \emph{would}: \\
The coordinator announced that the training session \textbf{would start} the following Monday.
\item The action happened before the announcement $\rightarrow$ past perfect: \\
She explained that the volunteers \textbf{had already visited} three villages.
\item General truth, still true $\rightarrow$ present simple is accepted (backshift also possible): \\
He told us that gender equality \textbf{is} (or \emph{was}) a human right.
\end{enumerate}
\textbf{Conclusion:} after a past reporting verb, shift tenses one step back, except for permanent truths.
\end{exoresolu}

\begin{exercice}[Sequence of Tenses -- Practice]
Put the verbs in brackets in the correct tense.
\begin{enumerate}
\item The teacher said that the elections (be) \ldots free and fair.
\item They reported that many girls (drop out) \ldots of school the previous year.
\item She claimed that she (can) \ldots organise the event alone.
\item We knew that the law (protect) \ldots women's rights.
\end{enumerate}
\end{exercice}

\begin{corrige}[Exercice Sequence of Tenses]
\begin{enumerate}
\item were (or are -- general truth)
\item had dropped out
\item could
\item protected (or protects -- general truth)
\end{enumerate}
\end{corrige}

\begin{methode}[Forming Impersonal Passives]
The impersonal passive (\emph{passif impersonnel}) is used in formal reports and news about elections or rights. Two constructions exist:
\begin{enumerate}
\item \textbf{It + passive verb + that-clause}: \emph{It is said that the elections were free.}
\item \textbf{Subject + passive verb + to-infinitive}: \emph{The elections are said to have been free.}
\item Use verbs of saying and thinking: \cle{say}, \cle{believe}, \cle{think}, \cle{report}, \cle{expect}, \cle{claim}, \cle{know}.
\item If the reported action is \textbf{earlier}, use the perfect infinitive \emph{to have + past participle}: \emph{She is believed to have voted.}
\item If the reported action is \textbf{simultaneous or future}, use the simple infinitive \emph{to + base}: \emph{He is expected to win.}
\end{enumerate}
\end{methode}

\begin{exoresolu}[Impersonal Passives -- Bac-Type Exercise]
\textbf{Task.} Rewrite each sentence using an impersonal passive, starting as shown.
\begin{enumerate}
\item People say that the campaign was peaceful. $\rightarrow$ It \ldots
\item Journalists report that many women voted in the election. $\rightarrow$ Many women \ldots
\item Everyone expects the new school leader to defend students' rights. $\rightarrow$ The new school leader \ldots
\end{enumerate}
\tcblower
\textbf{Solution.}
\begin{enumerate}
\item Construction 1 with \emph{It is said that}: \\
\textbf{It is said that the campaign was peaceful.}
\item Construction 2; the voting happened before the report $\rightarrow$ perfect infinitive: \\
\textbf{Many women are reported to have voted in the election.}
\item Construction 2 with \emph{expect} + simple infinitive: \\
\textbf{The new school leader is expected to defend students' rights.}
\end{enumerate}
\textbf{Conclusion:} choose \emph{It is said that\ldots} or \emph{subject + is said + to-infinitive}; use the perfect infinitive for earlier actions.
\end{exoresolu}

\begin{exercice}[Impersonal Passives -- Practice]
Rewrite using an impersonal passive.
\begin{enumerate}
\item People believe that the new law protects minorities. $\rightarrow$ It \ldots
\item They say that the turnout was very high. $\rightarrow$ The turnout \ldots
\item Experts think that the programme will empower women. $\rightarrow$ The programme \ldots
\item Everyone knows that she worked hard. $\rightarrow$ She \ldots
\end{enumerate}
\end{exercice}

\begin{corrige}[Exercice Impersonal Passives]
\begin{enumerate}
\item It is believed that the new law protects minorities.
\item The turnout is said to have been very high.
\item The programme is thought to empower women. (or: to be empowering women)
\item She is known to have worked hard.
\end{enumerate}
\end{corrige}

\courssub{Writing: Composition on Volunteering for Gender Equality}

\begin{methode}[Writing a Composition on Volunteering for Gender Equality]
For the Bac essay (about 150--200 words), follow this plan:
\begin{enumerate}
\item \textbf{Read the topic carefully} and identify the task: describe an activity, give your opinion, or propose solutions.
\item \textbf{Brainstorm} vocabulary from the unit: \emph{volunteer, raise awareness, equal rights, campaign, empower girls, training course}.
\item \textbf{Introduction} (2--3 sentences): present the topic and its importance in your community.
\item \textbf{Body} (2 paragraphs): describe concrete actions (door-to-door campaigns, girls' clubs, surveys) using phrasal verbs and purpose clauses; give results and your opinion.
\item \textbf{Conclusion} (2 sentences): summarise and give a recommendation or a hope for the future.
\item \textbf{Link your ideas} with connectors: \emph{first, moreover, however, as a result, in conclusion}.
\item \textbf{Proofread}: check sequence of tenses, word order, subject-verb agreement, and spelling.
\end{enumerate}
\end{methode}

\begin{exoresolu}[Guided Composition -- Bac-Type Exercise]
\textbf{Task.} Your English club took part in a volunteering activity to promote gender equality in your neighbourhood. Write a composition of about 120 words describing the activity and its results.
\tcblower
\textbf{Solution (model answer).}

\emph{Last month, the English club of our school carried out a volunteering campaign to promote gender equality in our neighbourhood.}

\emph{First, we set up a small team and organised door-to-door visits in order to raise awareness about girls' education. Moreover, we took part in a training session on human rights, where we learned that every child has the right to go to school. During the campaign, we distributed posters and answered the questions of parents.}

\emph{As a result, many families promised to send their daughters to school. It is believed that our action changed mentalities. In conclusion, volunteering is a powerful way to build a fairer society, and I encourage every student to take part in such activities.}

\textbf{Marking points:} clear plan (introduction, body, conclusion); phrasal verbs (\emph{carried out, set up, took part in}); purpose clause (\emph{in order to raise}); impersonal passive (\emph{it is believed that}); connectors; correct tenses.
\end{exoresolu}

\begin{exercice}[Writing Practice -- Bac-Type]
Choose ONE topic and write 150--200 words.
\begin{enumerate}
\item Describe a volunteering experience you had (or imagine one) to promote gender equality in your school. Explain what you did and what changed.
\item ``Volunteering is the best way to learn citizenship.'' Do you agree? Give reasons and examples.
\item Write a formal letter to the Mayor proposing a youth volunteer programme for gender equality in your municipality.
\end{enumerate}
\end{exercice}

\courssub{Integration Activities / Evaluation / Remediation}

\begin{aretenir}[Key Points for the Bac]
\begin{itemize}
\item \textbf{Reported questions:} no inversion, backshift tenses, use \emph{if/whether} for yes/no questions.
\item \textbf{Word order:} S+V+O; frequency adverbs before main verb (after \emph{be}); adjectives before nouns; Manner-Place-Time.
\item \textbf{Phrasal verbs:} separable vs inseparable; pronoun goes in middle of separable verbs; express purpose with \emph{to/in order to}.
\item \textbf{Sequence of tenses:} backshift after past reporting verbs; exception for general truths.
\item \textbf{Impersonal passive:} \emph{It is said that\ldots} OR \emph{Subject + is said + to-infinitive}; perfect infinitive for earlier actions.
\item \textbf{Writing:} structured paragraphs, connectors, unit vocabulary, varied grammar.
\end{itemize}
\end{aretenir}

\begin{attention}[The 5 Mistakes That Cost Marks at the Bac]
\begin{enumerate}
\item \textbf{Keeping the inversion in reported questions.} Never write ``She asked where was the polling station''; write \emph{She asked where the polling station was}. Statement order is compulsory.
\item \textbf{Forgetting the backshift of tenses.} After \emph{said, asked, explained}, \emph{will} must become \emph{would}, \emph{have voted} becomes \emph{had voted}. Only permanent truths may stay in the present.
\item \textbf{Wrong position of frequency adverbs.} Write \emph{Students always attend the course}, not ``Students attend always the course''; but \emph{She is always motivated} (after \emph{to be}).
\item \textbf{Misplacing the pronoun with separable phrasal verbs.} Write \emph{They carried it out}, never ``They carried out it''.
\item \textbf{Faulty impersonal passive.} Do not write ``It is said the elections to be free''. Use either \emph{It is said that the elections were free} or \emph{The elections are said to have been free} -- never mix the two structures.
\end{enumerate}
\end{attention}

\begin{experience}[Integration Activity: Mock Election Project]
\textbf{Objective:} Apply all unit skills in a simulated school election.
\begin{enumerate}
\item \textbf{Week 1 (Speaking/Vocabulary):} Form parties, write manifestos, hold a debate. Record key vocabulary.
\item \textbf{Week 2 (Grammar/Listening):} Listen to a talk on electoral transparency; report questions from the debate using reported speech.
\item \textbf{Week 3 (Reading/Writing):} Read a case study on gender quotas in student councils; write a proposal for gender-balanced representation using phrasal verbs, sequence of tenses, and impersonal passives.
\item \textbf{Week 4 (Evaluation):} Conduct the election; observers write a formal report using impersonal passives. Peer-assess compositions with a rubric covering vocabulary, grammar, organisation.
\end{enumerate}
\textbf{Remediation:} Students with difficulties review the \texttt{methode} boxes and redo the \texttt{exercice} boxes with teacher support.
\end{experience}

\newpage

\coursec{Exercices}

\courssub{Je vérifie mes connaissances}

\begin{exercice}[QCM : Campaigning and voting]
Choose the correct answer (A, B or C).

1. The document in which a candidate presents his or her ideas and promises is called a \ldots
\begin{itemize}
\item[A.] ballot box \quad B. manifesto \quad C. survey
\end{itemize}

2. To \emph{canvass} means \ldots
\begin{itemize}
\item[A.] to count the votes \quad B. to ask people for their support \quad C. to close the polling station
\end{itemize}

3. A \emph{polling station} is \ldots
\begin{itemize}
\item[A.] the place where people vote \quad B. a training centre \quad C. a radio station
\end{itemize}

4. ``She asked me \ldots I would vote for the new head girl.''
\begin{itemize}
\item[A.] that \quad B. if \quad C. what
\end{itemize}
\end{exercice}

\begin{exercice}[Vrai ou faux justifié]
Say whether the following statements are \textbf{True} or \textbf{False}, and justify your answer with one sentence.

1. In a reported yes/no question, we use \emph{if} or \emph{whether}.
2. In a reported question, we keep the question word order (auxiliary before the subject).
3. ``It is believed that volunteering empowers young women'' is an example of an impersonal passive.
4. A phrasal verb is a verb followed by one or two particles whose meaning is often different from the verb alone.
\end{exercice}

\begin{exercice}[Définitions : démocratie et citoyenneté]
Match each word (1--6) with its definition (a--f).

1. democracy \quad 2. volunteer \quad 3. gender equality \quad 4. facilitator \quad 5. election \quad 6. survey

\begin{itemize}
\item[a.] a person who leads a training session or a discussion
\item[b.] a system of government in which citizens choose their leaders by voting
\item[c.] a person who offers to work for free for a good cause
\item[d.] the fact that men and women have the same rights and opportunities
\item[e.] a set of questions asked to a group of people to collect opinions
\item[f.] the process of choosing a person for a position by voting
\end{itemize}
\end{exercice}

\begin{exercice}[Grammaire directe : sequence of tenses]
Complete each sentence with the correct form of the verb in brackets.

1. The facilitator said that the training \ldots (start) at 8 a.m. the next day.
2. She explained that she \ldots (join) the association two years before.
3. They told us that they \ldots (never / see) a ballot box before the workshop.
4. He promised that he \ldots (help) us to organise the campaign.
\end{exercice}

\courssub{Leçon : Campaigning, Democracy and Volunteering}

\begin{definition}[Vocabulaire clé : Campagne et vote scolaire]
\begin{center}
\begin{tabularx}{\linewidth}{|l|X|l|}\hline
\textbf{English term} & \textbf{Definition / Context} & \textbf{Français} \\ \hline
manifesto & a public declaration of aims and policies & profession de foi, manifeste \\
to canvass & to ask people for their votes or support & faire campagne, solliciter des voix \\
polling station & place where people go to vote & bureau de vote \\
ballot box & sealed box where voters put their ballot papers & urne \\
ballot paper & paper on which voters mark their choice & bulletin de vote \\
head girl / head boy & student elected to represent the student body & déléguée / délégué général des élèves \\
to stand for election & to be a candidate in an election & se présenter aux élections \\
turnout & the number of people who vote & taux de participation \\ \hline
\end{tabularx}
\end{center}
\end{definition}

\begin{definition}[Vocabulaire clé : Démocratie et formation]
\begin{center}
\begin{tabularx}{\linewidth}{|l|X|l|}\hline
\textbf{English term} & \textbf{Definition / Context} & \textbf{Français} \\ \hline
democracy & system of government by the whole population & démocratie \\
facilitator & person who helps a group work together effectively & animateur, facilitateur \\
workshop & intensive training session with practical activities & atelier de formation \\
mock election & simulated election for practice & élection simulée \\
civic engagement & participation in community and political life & engagement civique \\
awareness campaign & organised effort to inform people about an issue & campagne de sensibilisation \\ \hline
\end{tabularx}
\end{center}
\end{definition}

\begin{definition}[Vocabulaire clé : Volontariat et égalité de genre]
\begin{center}
\begin{tabularx}{\linewidth}{|l|X|l|}\hline
\textbf{English term} & \textbf{Definition / Context} & \textbf{Français} \\ \hline
volunteer (n/v) & person who works freely for a cause / to offer freely & bénévole / faire du bénévolat \\
gender equality & equal rights and opportunities for all genders & égalité entre les sexes \\
drop out (phr. v.) & leave school or a programme before finishing & abandonner (l'école) \\
carry out (phr. v.) & perform or conduct (a survey, study, etc.) & mener, réaliser \\
set up (phr. v.) & establish or create (an organisation, fund) & créer, mettre en place \\
look after (phr. v.) & take care of & s'occuper de \\
empower & give power, confidence, or authority to & autonomiser, donner du pouvoir à \\
dropout rate & percentage of students leaving school early & taux d'abandon scolaire \\ \hline
\end{tabularx}
\end{center}
\end{definition}

\begin{propriete}[Reported Questions -- Règles essentielles]
\begin{itemize}
\item \textbf{Yes/No questions} : introductory verb + \emph{if} / \emph{whether} + subject + verb (statement order).\\
Ex: ``Do you know the candidates?'' $\rightarrow$ She asked \emph{if I knew} the candidates.
\item \textbf{Wh- questions} : introductory verb + question word + subject + verb (statement order).\\
Ex: ``Where is the polling station?'' $\rightarrow$ He wanted to know \emph{where the polling station was}.
\item \textbf{Tense shift (backshift)} : present $\rightarrow$ past, past $\rightarrow$ past perfect, will $\rightarrow$ would, can $\rightarrow$ could.
\item \textbf{No auxiliary do/does/did} in the reported clause. No question mark.
\end{itemize}
\end{propriete}

\begin{propriete}[Word Order -- Ordre des mots en anglais (SVOMPT)]
\textbf{Subject + Verb + Object + Manner + Place + Time}\\
\begin{itemize}
\item Adverbs of frequency (always, often, usually) go \textbf{before} the main verb but \textbf{after} be/auxiliaries.
\item Adverbs of manner (carefully, attentively) usually go \textbf{after} the verb or object.
\item In passive voice: Subject + be + past participle + (by agent) + Manner/Place/Time.
\end{itemize}
\end{propriete}

\begin{propriete}[Phrasal Verbs -- Verbes à particules (sens idiomatique)]
\begin{center}
\begin{tabularx}{\linewidth}{|l|l|X|l|}\hline
\textbf{Phrasal verb} & \textbf{Type} & \textbf{Meaning} & \textbf{Français} \\ \hline
carry out & separable & conduct, perform & mener, réaliser \\
drop out & intransitive & leave school/activity & abandonner \\
set up & separable & establish, create & créer, mettre en place \\
look after & inseparable (prep.) & take care of & s'occuper de \\
stand for & inseparable (prep.) & represent, be a candidate & représenter, se présenter \\
give up & separable & stop doing, quit & abandonner, renoncer à \\
find out & separable & discover & découvrir \\ \hline
\end{tabularx}
\end{center}
\end{propriete}

\begin{propriete}[Sequence of Tenses -- Concordance des temps]
Lorsque le verbe introducteur est au \textbf{passé} (said, told, explained, promised, announced), les temps se décalent :
\begin{center}
\begin{tabular}{|l|l|l|}\hline
\textbf{Direct speech} & \textbf{Reported speech} & \textbf{Exemple} \\ \hline
Present simple & Past simple & ``I know'' $\rightarrow$ He said he \emph{knew} \\
Present continuous & Past continuous & ``I am waiting'' $\rightarrow$ She said she \emph{was waiting} \\
Present perfect & Past perfect & ``I have finished'' $\rightarrow$ He said he \emph{had finished} \\
Past simple & Past perfect & ``I voted'' $\rightarrow$ She said she \emph{had voted} \\
Will & Would & ``I will help'' $\rightarrow$ He promised he \emph{would help} \\
Can & Could & ``I can come'' $\rightarrow$ She said she \emph{could come} \\ \hline
\end{tabular}
\end{center}
\emph{Note :} Past perfect et modaux would/could/should/might ne changent pas.
\end{propriete}

\begin{propriete}[Impersonal Passive -- Passif impersonnel]
Deux structures pour rapporter des opinions/généralités sans citer la source :
\begin{enumerate}
\item \textbf{It + passive verb + that-clause} : \emph{It is said / believed / thought / considered / expected / reported that...}
\item \textbf{Subject + passive verb + to-infinitive} : \emph{He is said to be... / The project is expected to start...}
\end{enumerate}
\emph{Verbes courants :} say, believe, think, consider, expect, know, report, claim, allege.
\end{propriete}

\courssub{Je m'entraîne}

\begin{exercice}[Reported questions]
Rewrite the following questions in reported speech, starting as indicated.

1. ``Will you vote for the new head girl?'' the journalist asked me.\\
The journalist asked me \ldots
2. ``Where did the training take place?'' she wanted to know.\\
She wanted to know \ldots
3. ``Have you ever taken part in a survey?'' the interviewer asked the students.\\
The interviewer asked the students \ldots
4. ``Why do so few girls volunteer for the campaign?'' the teacher wondered.\\
The teacher wondered \ldots
5. ``When will the results be announced?'' the candidates asked.\\
The candidates asked \ldots
\end{exercice}

\begin{exercice}[Word order]
Put the words in the correct order to make correct sentences.

1. always / the / participants / attentively / listen / facilitator / to / the
2. the / girls / encouraged / to / teacher / stand / for / election / the
3. yesterday / a / very / gave / interesting / the / volunteer / speech
4. carefully / the / filled in / voters / their / ballot papers
5. often / surveys / are / conducted / before / school / elections
\end{exercice}

\begin{exercice}[Texte à trous : une campagne électorale scolaire]
Complete the text with the following words: \emph{ballot box -- manifesto -- to canvass -- polling station -- survey}.

Last month, our school elected its new head boy and head girl. Before the election, each candidate wrote a \ldots (1) presenting his or her projects for the students. Then the candidates went from class to class \ldots (2) for votes. A few days before the election, a \ldots (3) was organised to find out which candidate was the most popular. On election day, all the students went to the \ldots (4), which was set up in the school hall, and put their ballot papers into the \ldots (5). The results were announced in the afternoon.
\end{exercice}

\begin{exercice}[Impersonal passives]
Rewrite each sentence in two ways, as in the example.

Example: People believe that volunteering empowers young women.\\
$\rightarrow$ It is believed that volunteering empowers young women.\\
$\rightarrow$ Volunteering is believed to empower young women.

1. People say that the training course was very useful.
2. Everyone thinks that the new head girl will defend students' rights.
3. Experts consider that gender equality improves school results.
4. People expect that many students will take part in the survey.
\end{exercice}

\courssub{Je me prépare au Bac}

\begin{exercice}[Type Bac -- Reading comprehension]
Read the text carefully, then answer the questions.

\textbf{Girls' Education in the Far North: The Work of ``Bright Future''}

In 2015, a group of young Cameroonian graduates created an association called ``Bright Future'' in Maroua, in the Far North region. Their objective was simple but ambitious: to help more girls stay in school and complete their secondary education. In this region, many girls leave school before the age of fifteen, often because of early marriage, poverty, or long distances between home and school.

The association started with only twelve volunteers, most of them women teachers and nurses. They first carried out a survey in five villages to understand why so many girls dropped out. The results showed that parents often did not know about the importance of girls' education, and that many families could not afford school fees and uniforms.

``Bright Future'' then organised awareness campaigns in markets, churches and mosques. The volunteers also set up a small fund to pay school fees for the poorest families. In addition, they offered free evening classes in mathematics and English to girls who had difficulties at school. A local facilitator was trained in each village to follow the girls' progress and to report any problem to the association.

Ten years later, the results are impressive. More than two thousand girls have been supported, and the dropout rate in the five villages has fallen by almost forty per cent. Several girls who were helped by the association have now become volunteers themselves. ``It is said that educating a girl means educating a whole nation,'' explains Aïcha, one of the founders. ``Our experience proves that this is true.''

The association now plans to extend its activities to ten more villages and to open a training centre where girls will learn computer skills and sewing.

\textbf{Questions}

1. Say whether the following statements are \textbf{True} or \textbf{False}. Justify each answer with a quotation from the text.
\begin{itemize}
\item[a.] ``Bright Future'' was created by the government of Cameroon.
\item[b.] Most of the first volunteers were men.
\item[c.] The survey helped the association understand the causes of school dropout.
\item[d.] The association only paid school fees for poor families.
\item[e.] Some former beneficiaries now work as volunteers.
\end{itemize}

2. Answer these questions in your own words.
\begin{itemize}
\item[a.] Give two reasons why girls leave school early in the Far North region.
\item[b.] What activities did the association organise to help the girls? (Name three.)
\item[c.] What are the association's projects for the future?
\end{itemize}

3. Vocabulary: find in the text words or expressions that mean:
\begin{itemize}
\item[a.] stopped going to school (paragraph 2)
\item[b.] money paid to study at school (paragraph 2)
\item[c.] information campaigns to change people's attitudes (paragraph 3)
\end{itemize}
\end{exercice}

\begin{exercice}[Type Bac -- Grammar in context]
\textbf{A. Reported questions.} During the school elections, a journalist interviewed students. Report the following questions.

1. ``Do you know the candidates' manifestos?'' the journalist asked the students.
2. ``Who will you vote for?'' she asked Manga.
3. ``Why did you decide to stand for election?'' she asked the new head girl.

\textbf{B. Phrasal verbs.} Complete the sentences with the correct phrasal verb from the list: \emph{to carry out -- to drop out -- to set up -- to look after}.

1. The association \ldots a survey in five villages last year.
2. Many girls \ldots of school before the age of fifteen.
3. The volunteers \ldots a fund to help poor families.
4. Aïcha \ldots the evening classes every week.

\textbf{C. Sequence of tenses and impersonal passive.}

1. The facilitator announced that the results \ldots (publish) the following day.
2. Rewrite: ``People say that educating a girl means educating a whole nation.'' Begin with: \emph{It is said that \ldots}
\end{exercice}

\begin{exercice}[Type Bac -- Writing]
Your school organised a training course on democracy for all the students of Terminale. Write a report (about 250 words) for the school magazine in which you:
\begin{itemize}
\item describe the activities that took place during the course (workshops, debates, mock elections, surveys\ldots);
\item explain what you personally learnt from the training;
\item say why such courses matter for young Cameroonians.
\end{itemize}

Your report must have a title, an introduction, two or three body paragraphs and a conclusion. Use at least three phrasal verbs, one reported question and one impersonal passive. Pay attention to word order and to the sequence of tenses.
\end{exercice}

\newpage

\coursec{Corrigés des exercices}

\begin{corrige}[Exercice 1 — QCM : Campaigning and voting]

\textbf{1. B — manifesto.} A \emph{manifesto} is the document in which a candidate presents his or her ideas and promises (profession de foi). A \emph{ballot box} is the urn where ballot papers are placed, and a \emph{survey} is a set of questions used to collect opinions.

\textbf{2. B — to ask people for their support.} \emph{To canvass} means to go from person to person (or class to class) asking for votes or support (faire campagne, solliciter des voix). Counting the votes is done after the poll closes; closing the polling station is the end of voting.

\textbf{3. A — the place where people vote.} A \emph{polling station} (bureau de vote) is the place where voters cast their ballots. A training centre is a place for courses; a radio station broadcasts programmes.

\textbf{4. B — if.} ``She asked me \emph{if} I would vote for the new head girl.'' This is a reported \textbf{yes/no question}: we introduce it with \emph{if} or \emph{whether}, we use statement word order (subject + verb), and \emph{will} shifts to \emph{would}. \emph{That} is used for reported statements, not questions; \emph{what} would require a wh- question in direct speech.
\end{corrige}

\begin{corrige}[Exercice 2 — Vrai ou faux justifié]

\textbf{1. TRUE.} In a reported yes/no question, we use \emph{if} or \emph{whether} followed by statement word order: ``Are you ready?'' $\rightarrow$ He asked \emph{if I was ready}.

\textbf{2. FALSE.} In a reported question we use the \textbf{statement word order} (subject + verb), not the question order: ``Where is the station?'' $\rightarrow$ He asked \emph{where the station was} (not ``where was the station'').

\textbf{3. TRUE.} ``It is believed that volunteering empowers young women'' follows the structure \emph{It + passive verb + that-clause}, which is the first form of the impersonal passive.

\textbf{4. TRUE.} A phrasal verb combines a verb with one or two particles (adverb or preposition), and its meaning is often idiomatic: \emph{to carry out} does not mean ``to carry outside'' but ``to conduct, to perform'' (mener, réaliser).
\end{corrige}

\begin{corrige}[Exercice 3 — Définitions : démocratie et citoyenneté]

\begin{center}
\begin{tabular}{|c|c|l|}\hline
\textbf{Word} & \textbf{Definition} & \textbf{Justification} \\ \hline
1. democracy & b & system of government in which citizens choose their leaders by voting \\
2. volunteer & c & a person who offers to work for free for a good cause \\
3. gender equality & d & the fact that men and women have the same rights and opportunities \\
4. facilitator & a & a person who leads a training session or a discussion \\
5. election & f & the process of choosing a person for a position by voting \\
6. survey & e & a set of questions asked to a group of people to collect opinions \\ \hline
\end{tabular}
\end{center}

\textbf{Points de vigilance :} ne pas confondre \emph{election} (le processus de vote) et \emph{democracy} (le système politique) ; un \emph{survey} (enquête) est différent d'une \emph{election} (on y donne une opinion, on n'y choisit pas un dirigeant).
\end{corrige}

\begin{corrige}[Exercice 4 — Grammaire directe : sequence of tenses]

\textbf{1.} The facilitator said that the training \textbf{would start} at 8 a.m. the next day.\\
\emph{Justification :} le verbe introducteur \emph{said} est au passé ; \emph{will} se transforme en \emph{would} (futur dans le passé).

\textbf{2.} She explained that she \textbf{had joined} the association two years \textbf{before}.\\
\emph{Justification :} le past simple (\emph{joined}) recule au past perfect (\emph{had joined}) ; \emph{ago} devient \emph{before}.

\textbf{3.} They told us that they \textbf{had never seen} a ballot box before the workshop.\\
\emph{Justification :} le present perfect (\emph{have never seen}) recule au past perfect (\emph{had never seen}) ; l'adverbe \emph{never} se place entre \emph{had} et le participe passé.

\textbf{4.} He promised that he \textbf{would help} us to organise the campaign.\\
\emph{Justification :} \emph{will} devient \emph{would} après un verbe introducteur au passé (\emph{promised}).

\textbf{Point de vigilance :} ne jamais garder \emph{will} ou le present perfect après \emph{said / told / explained / promised} au passé : c'est l'erreur la plus fréquente au Bac.
\end{corrige}

\begin{corrige}[Exercice 5 — Reported questions]

\textbf{1.} The journalist asked me \textbf{if I would vote for the new head girl}.\\
\emph{Yes/no question} $\rightarrow$ \emph{if} + ordre de la déclarative ; \emph{will} $\rightarrow$ \emph{would}.

\textbf{2.} She wanted to know \textbf{where the training had taken place}.\\
\emph{Wh- question} $\rightarrow$ on garde \emph{where} ; le past simple (\emph{did take}) recule au past perfect (\emph{had taken}) ; suppression de l'auxiliaire \emph{did} et ordre sujet + verbe.

\textbf{3.} The interviewer asked the students \textbf{if they had ever taken part in a survey}.\\
\emph{Yes/no question} $\rightarrow$ \emph{if} ; le present perfect (\emph{have taken}) recule au past perfect (\emph{had taken}) ; \emph{you} devient \emph{they} (accord de la personne).

\textbf{4.} The teacher wondered \textbf{why so few girls volunteered for the campaign}.\\
\emph{Wh- question} $\rightarrow$ on garde \emph{why} ; le present simple (\emph{do volunteer}) recule au past simple (\emph{volunteered}) ; suppression de l'auxiliaire \emph{do}.

\textbf{5.} The candidates asked \textbf{when the results would be announced}.\\
\emph{Wh- question} $\rightarrow$ on garde \emph{when} ; \emph{will be announced} (futur passif) devient \emph{would be announced} ; ordre sujet (\emph{the results}) + verbe.

\textbf{Points de vigilance :} pas de point d'interrogation à la fin ; pas d'inversion sujet--verbe ; pas d'auxiliaire \emph{do/does/did} ; ne pas oublier le recul des temps.
\end{corrige}

\begin{corrige}[Exercice 6 — Word order]

\textbf{1.} The participants always listen attentively to the facilitator.\\
\emph{SVOMPT :} sujet (\emph{the participants}) + adverbe de fréquence avant le verbe (\emph{always listen}) + adverbe de manière après le verbe (\emph{attentively}) + complément (\emph{to the facilitator}).

\textbf{2.} The teacher encouraged the girls to stand for election.\\
\emph{Structure :} sujet + verbe + objet (\emph{the girls}) + complément infinitif (\emph{to stand for election}).

\textbf{3.} The volunteer gave a very interesting speech yesterday.\\
\emph{Ordre :} sujet + verbe + objet (\emph{a very interesting speech}) + complément de temps en fin de phrase (\emph{yesterday}).

\textbf{4.} The voters filled in their ballot papers carefully.\\
\emph{Ordre :} sujet + verbe à particule (\emph{filled in}) + objet (\emph{their ballot papers}) + adverbe de manière (\emph{carefully}) en fin de phrase.

\textbf{5.} Surveys are often conducted before school elections.\\
\emph{Voix passive :} sujet (\emph{surveys}) + auxiliaire \emph{are} + adverbe de fréquence après l'auxiliaire (\emph{often}) + participe passé (\emph{conducted}) + complément de temps (\emph{before school elections}).

\textbf{Point de vigilance :} les adverbes de fréquence se placent \textbf{avant} le verbe principal mais \textbf{après} l'auxiliaire ou le verbe \emph{be} (phrase 5 : \emph{are often conducted}).
\end{corrige}

\begin{corrige}[Exercice 7 — Texte à trous : une campagne électorale scolaire]

\textbf{(1) manifesto} — chaque candidat rédige une profession de foi présentant ses projets.\\
\textbf{(2) to canvass} — les candidats vont de classe en classe \emph{pour solliciter} des voix (infinitif de but : \emph{to} + base verbale).\\
\textbf{(3) survey} — une enquête est organisée pour connaître le candidat le plus populaire.\\
\textbf{(4) polling station} — le jour du scrutin, les élèves se rendent au bureau de vote installé dans le hall.\\
\textbf{(5) ballot box} — ils déposent leurs bulletins de vote dans l'urne.

\textbf{Texte complété :} Last month, our school elected its new head boy and head girl. Before the election, each candidate wrote a \emph{manifesto} presenting his or her projects for the students. Then the candidates went from class to class \emph{to canvass} for votes. A few days before the election, a \emph{survey} was organised to find out which candidate was the most popular. On election day, all the students went to the \emph{polling station}, which was set up in the school hall, and put their ballot papers into the \emph{ballot box}. The results were announced in the afternoon.

\textbf{Point de vigilance :} pour (2), il faut l'infinitif complet \emph{to canvass} (expression de but) et non la base verbale seule.
\end{corrige}

\begin{corrige}[Exercice 8 — Impersonal passives]

\textbf{1.} People say that the training course was very useful.\\
$\rightarrow$ \textbf{It is said that the training course was very useful.}\\
$\rightarrow$ \textbf{The training course is said to have been very useful.}\\
\emph{Note :} la deuxième structure exige l'infinitif passé \emph{to have been} car le cours est terminé (antériorité).

\textbf{2.} Everyone thinks that the new head girl will defend students' rights.\\
$\rightarrow$ \textbf{It is thought that the new head girl will defend students' rights.}\\
$\rightarrow$ \textbf{The new head girl is thought to be going to defend students' rights} (ou : \emph{is expected to defend students' rights}).\\
\emph{Note :} \emph{will} n'a pas d'infinitif ; on utilise \emph{to be going to} ou le verbe \emph{expect}.

\textbf{3.} Experts consider that gender equality improves school results.\\
$\rightarrow$ \textbf{It is considered that gender equality improves school results.}\\
$\rightarrow$ \textbf{Gender equality is considered to improve school results.}

\textbf{4.} People expect that many students will take part in the survey.\\
$\rightarrow$ \textbf{It is expected that many students will take part in the survey.}\\
$\rightarrow$ \textbf{Many students are expected to take part in the survey.}

\textbf{Points de vigilance :} dans la structure 2, c'est le \textbf{sujet de la subordonnée} qui devient sujet de la principale ; vérifier l'accord de \emph{be} (\emph{is/are}) ; utiliser l'infinitif passé (\emph{to have done}) si l'action rapportée est antérieure.
\end{corrige}

\begin{corrige}[Exercice 9 — Type Bac : Reading comprehension]

\textbf{Barème indicatif (sur 10 points) :} Question 1 : 5 pts (1 pt par item : 0,5 pour True/False + 0,5 pour la citation) ; Question 2 : 3 pts (1 pt par item) ; Question 3 : 2 pts (environ 0,7 pt par item).

\textbf{1. True or False, with quotations}

\begin{itemize}
\item[a.] \textbf{FALSE.} ``A group of young Cameroonian graduates created an association called `Bright'.'' (Ce sont des diplômés, pas le gouvernement.)
\item[b.] \textbf{FALSE.} ``The association started with only twelve volunteers, \emph{most of them women} teachers and nurses.''
\item[c.] \textbf{TRUE.} ``They first carried out a survey in five villages \emph{to understand why so many girls dropped out}.''
\item[d.] \textbf{FALSE.} ``\emph{In addition}, they offered free evening classes in mathematics and English'' — l'association faisait plus que payer les frais scolaires (campagnes de sensibilisation, cours du soir).
\item[e.] \textbf{TRUE.} ``Several girls who were helped by the association \emph{have now become volunteers themselves}.''
\end{itemize}

\textbf{2. Answers in your own words}

\begin{itemize}
\item[a.] Two reasons: \textbf{early marriage} and \textbf{poverty} (on accepte aussi : long distances between home and school).
\item[b.] Three activities: they organised \textbf{awareness campaigns} in markets, churches and mosques; they \textbf{set up a fund} to pay school fees for the poorest families; they offered \textbf{free evening classes} in mathematics and English (on accepte aussi : training a local facilitator in each village).
\item[c.] Future projects: they plan to \textbf{extend their activities to ten more villages} and to \textbf{open a training centre} where girls will learn computer skills and sewing.
\end{itemize}

\textbf{3. Vocabulary}

\begin{itemize}
\item[a.] \textbf{dropped out} (stopped going to school)
\item[b.] \textbf{school fees} (money paid to study at school)
\item[c.] \textbf{awareness campaigns} (information campaigns to change people's attitudes)
\end{itemize}

\textbf{Points de vigilance (attendus du correcteur) :} la justification par une \textbf{citation exacte} du texte est obligatoire pour la question 1 (aucun point complet sans guillemets) ; pour la question 2, le candidat doit \textbf{reformuler} avec ses propres mots, sans recopier intégralement les phrases du texte ; pour la question 3, le mot doit être tiré du \textbf{paragraphe indiqué}.
\end{corrige}

\begin{corrige}[Exercice 10 — Type Bac : Grammar in context]

\textbf{A. Reported questions}

\textbf{1.} The journalist asked the students \textbf{if they knew the candidates' manifestos}.\\
\emph{Yes/no question} $\rightarrow$ \emph{if} ; present simple \emph{know} $\rightarrow$ past simple \emph{knew} ; \emph{you} $\rightarrow$ \emph{they}.

\textbf{2.} She asked Manga \textbf{who he would vote for}.\\
\emph{Wh- question} $\rightarrow$ on garde \emph{who} ; \emph{will} $\rightarrow$ \emph{would} ; ordre sujet + verbe (\emph{he would vote}) ; la préposition \emph{for} reste en fin de proposition.

\textbf{3.} She asked the new head girl \textbf{why she had decided to stand for election}.\\
\emph{Wh- question} $\rightarrow$ on garde \emph{why} ; past simple (\emph{did decide}) $\rightarrow$ past perfect (\emph{had decided}) ; suppression de \emph{did}.

\textbf{B. Phrasal verbs}

\textbf{1.} The association \textbf{carried out} a survey in five villages last year.\\
\emph{Justification :} \emph{to carry out a survey} = mener une enquête ; past simple car \emph{last year}.

\textbf{2.} Many girls \textbf{drop out} of school before the age of fifteen.\\
\emph{Justification :} \emph{to drop out of school} = abandonner l'école ; present simple pour un fait général.

\textbf{3.} The volunteers \textbf{set up} a fund to help poor families.\\
\emph{Justification :} \emph{to set up} = créer, mettre en place ; past simple (\emph{set} est irrégulier : set--set--set).

\textbf{4.} Aïcha \textbf{looks after} the evening classes every week.\\
\emph{Justification :} \emph{to look after} = s'occuper de ; present simple, 3\textsuperscript{e} personne du singulier (\emph{looks}) car \emph{every week}.

\textbf{C. Sequence of tenses and impersonal passive}

\textbf{1.} The facilitator announced that the results \textbf{would be published} the following day.\\
\emph{Justification :} \emph{will} $\rightarrow$ \emph{would} après \emph{announced} ; la voix passive est nécessaire (\emph{be published}) car les résultats subissent l'action ; \emph{tomorrow} devient \emph{the following day}.

\textbf{2.} \textbf{It is said that educating a girl means educating a whole nation.}\\
\emph{Justification :} structure impersonnelle \emph{It + is said + that}-clause ; le gérondif \emph{educating} reste sujet de la subordonnée.

\textbf{Points de vigilance :} en A, aucune inversion sujet--verbe et aucun point d'interrogation ; en B, conjuguer le phrasal verb au temps qui convient au contexte ; en C1, ne pas oublier la voix passive (\emph{would be published}, pas \emph{would publish}).
\end{corrige}

\begin{corrige}[Exercice 11 — Type Bac : Writing]

\textbf{Barème indicatif (sur 20 points) :}
\begin{itemize}
\item Respect de la tâche et du format du report (titre, introduction, 2--3 paragraphes, conclusion, environ 250 mots) : 5 pts
\item Contenu : description des activités + apprentissage personnel + importance pour les jeunes Camerounais : 5 pts
\item Exigences grammaticales : au moins 3 phrasal verbs, 1 reported question, 1 impersonal passive : 4 pts
\item Correction grammaticale (word order, sequence of tenses) : 4 pts
\item Vocabulaire de l'unité et cohérence du texte : 2 pts
\end{itemize}

\textbf{Exemple de production attendue (modèle de correction) :}

\textbf{A Memorable Training Course on Democracy}

Last week, our school organised a three-day training course on democracy for all the students of Terminale. About two hundred students took part in this enriching experience.

On the first day, the facilitators \emph{carried out} a survey to find out what we already knew about citizenship. Then we attended workshops on voting procedures and human rights. The most exciting moment was the mock election: two students \emph{stood for} the position of head of the student council, and the whole class went to the polling station that had been \emph{set up} in the library. One facilitator asked us \emph{if we had ever voted before}, and most of us answered that we had not.

Personally, I learnt a lot from this course. I found out that democracy is not only about voting: it is also about listening to others, respecting different opinions and accepting defeat. I also understood that young people must not \emph{give up} when things are difficult; they must \emph{look after} their rights and their duties as citizens.

In my opinion, such courses matter enormously for young Cameroonians. \emph{It is believed that} educated citizens build stronger nations. If we learn the values of democracy at school, we will become responsible adults who can transform our communities. I hope our school will organise this training again next year.

(Environ 230--250 mots.)

\textbf{Points de vigilance :}
\begin{itemize}
\item Le titre est obligatoire ; un texte sans titre perd des points de format.
\item Vérifier la présence effective des trois outils grammaticaux exigés : ici \emph{carried out, stood for, set up, give up, look after} (phrasal verbs) ; \emph{asked us if we had ever voted before} (reported question avec backshift) ; \emph{It is believed that...} (impersonal passive).
\item Sequence of tenses : après \emph{answered / said / explained} au passé, les temps doivent reculer (\emph{we had not}).
\item Word order : adverbes de fréquence avant le verbe principal, compléments de temps en fin de phrase.
\item Rester dans le thème de l'unité (citizenship, democracy, voting) et utiliser le vocabulaire de la leçon.
\end{itemize}
\end{corrige}

\newpage

\coursec{Fiche bilan}

\begin{popfigure}
\begin{tikzpicture}[
  mindmap,
  concept color=popblue,
  level 1 concept/.append style={level distance=3.5cm, sibling angle=45},
  level 2 concept/.append style={level distance=2.5cm, sibling angle=45},
  every node/.style={font=\small, text=white},
  grow cyclic,
  text width=2.8cm,
  align=center
]
\node[concept, align=center]{Citizenship \&\\ Human Rights}
  child[concept color=poporange] { node[concept, align=center]{Speaking \&\\ Vocabulary\\ \small(Campaigning \& Voting)} }
  child[concept color=popgreen] { node[concept, align=center]{Grammar 1\\ \small(Reported Questions)} }
  child[concept color=poppurple] { node[concept, align=center]{Listening \&\\ Vocabulary\\ \small(Democracy Training)} }
  child[concept color=poppink] { node[concept, align=center]{Grammar 2\\ \small(Word Order)} }
  child[concept color=popgold] { node[concept, align=center]{Reading \&\\ Vocabulary\\ \small(Gender Equality)} }
  child[concept color=popdark] { node[concept, align=center]{Grammar 3 \&\\ Writing\\ \small(Phrasal Verbs, Sequence\\ of Tenses, Impersonal Passive)} };
\end{tikzpicture}
\legende{Mindmap of Lesson 04: Citizenship and Human Rights}
\end{popfigure}

\begin{bilan}

\courssub{Speaking \& Vocabulary: Campaigning and Voting}

\begin{definition}[Key Terms]
\begin{itemize}
  \item \cle{Manifesto} (\emph{programme électoral}) : A public declaration of policy and aims.
  \item \cle{Ballot} (\emph{bulletin de vote}) : The paper used to record a vote.
  \item \cle{Poll} (\emph{scrutin / sondage}) : The process of voting or a survey of opinion.
  \item \cle{Constituency} (\emph{circonscription}) : The body of voters in a specified area.
  \item \cle{Run for office} (\emph{se présenter aux élections}) : To be a candidate.
  \item \cle{Campaign trail} (\emph{la campagne électorale}) : The series of events/visits a candidate makes.
\end{itemize}
\end{definition}

\begin{exemplebox}[Useful Expressions for Exchanges]
\begin{itemize}
  \item \textbf{Asking about campaigns:} "What are the main points of your manifesto?" / "How do you plan to improve school facilities?"
  \item \textbf{Expressing opinion:} "I strongly support [Candidate] because..." / "I am undecided; I need to compare their policies."
  \item \textbf{Taking part in a poll:} "Have you filled in the survey?" / "The poll shows a tight race between the two candidates."
\end{itemize}
\end{exemplebox}

\courssub{Grammar Focus 1: Reported Questions}

\begin{propriete}[Rules for Reported Questions]
\begin{tabularx}{\linewidth}{|l|X|X|}\hline
\textbf{Type} & \textbf{Direct Question} & \textbf{Reported Question} \\ \hline
\cle{Yes/No Questions} & "Do you vote?" & He asked \textbf{if / whether} I voted. \\ \hline
\cle{Wh- Questions} & "Who won?" & She asked \textbf{who} had won. \\ \hline
\end{tabularx}
\begin{itemize}
  \item \textbf{Word Order:} Subject + Verb (NOT auxiliary do/does/did). \emph{Ex: He asked where the ballot box \textbf{was} (not: was the ballot box).}
  \item \textbf{Tense Backshift:} Present $\to$ Past; Past $\to$ Past Perfect; Will $\to$ Would; Can $\to$ Could.
  \item \textbf{No Question Mark:} Reported questions end with a full stop.
\end{itemize}
\end{propriete}

\courssub{Listening \& Vocabulary: Training Courses on Democracy}

\begin{definition}[Democracy Training Vocabulary]
\begin{itemize}
  \item \cle{Civic education} (\emph{éducation civique}) : Study of rights and duties of citizens.
  \item \cle{Workshop} (\emph{atelier}) : Interactive training session.
  \item \cle{Stakeholder} (\emph{partie prenante}) : Person/group with an interest in a decision.
  \item \cle{Consensus} (\emph{consensus}) : General agreement.
  \item \cle{Accountability} (\emph{reddition de comptes}) : Obligation to explain actions.
  \item \cle{Grassroots} (\emph{à la base / populaire}) : Ordinary people, not leaders.
\end{itemize}
\end{definition}

\courssub{Grammar Focus 2: Word Order}

\begin{propriete}[Standard English Word Order]
\textbf{SVOMPT} Rule: \cle{Subject} + \cle{Verb} + \cle{Object} + \cle{Manner} + \cle{Place} + \cle{Time}.
\begin{itemize}
  \item \textbf{Adverbs of Frequency:} Before main verb (\emph{always votes}), after \emph{be} (\emph{is always late}).
  \item \textbf{Indirect Object vs Direct Object:} \emph{Give [me] [the ballot]} or \emph{Give [the ballot] [to me]}.
  \item \textbf{Questions:} Auxiliary + Subject + Verb (ASV). \emph{Did you vote?}
  \item \textbf{Reported Speech:} Subject + Verb (SV). \emph{He asked if I had voted.}
\end{itemize}
\end{propriete}

\courssub{Reading \& Vocabulary: Volunteering in Gender Equality}

\begin{definition}[Gender Equality Vocabulary]
\begin{itemize}
  \item \cle{Gender bias} (\emph{préjugé sexiste}) : Prejudice based on gender.
  \item \cle{Empowerment} (\emph{autonomisation / renforcement des capacités}) : Giving power/authority.
  \item \cle{Advocacy} (\emph{plaidoyer}) : Public support for a cause.
  \item \cle{Stereotype} (\emph{stéréotype}) : Fixed, oversimplified idea.
  \item \cle{Glass ceiling} (\emph{plafond de verre}) : Invisible barrier to promotion.
  \item \cle{Inclusive} (\emph{inclusif}) : Not excluding any section of society.
\end{itemize}
\end{definition}

\courssub{Grammar Focus 3 \& Writing}

\begin{propriete}[Phrasal Verbs (Purpose/Action)]
\begin{tabularx}{\linewidth}{|l|l|X|}\hline
\textbf{Phrasal Verb} & \textbf{Meaning (FR)} & \textbf{Example} \\ \hline
\cle{Stand for} & Représenter / Se présenter & She decided to \textbf{stand for} president. \\ \hline
\cle{Speak out} & S'exprimer hautement & We must \textbf{speak out} against discrimination. \\ \hline
\cle{Take part in} & Participer à & Many students \textbf{took part in} the survey. \\ \hline
\cle{Give up} & Abandonner & Don't \textbf{give up} on your rights. \\ \hline
\cle{Look into} & Examiner / Enquêter & The committee will \textbf{look into} the complaints. \\ \hline
\end{tabularx}
\end{propriete}

\begin{propriete}[Sequence of Tenses (Concordance des temps)]
If the reporting verb is \textbf{Past} (\emph{said, asked, reported}), the reported verb shifts back:
\begin{align*}
\text{Present Simple} &\to \text{Past Simple} \\
\text{Present Perfect} &\to \text{Past Perfect} \\
\text{Past Simple} &\to \text{Past Perfect} \\
\text{Will} &\to \text{Would} \\
\text{Can} &\to \text{Could}
\end{align*}
\emph{Exception:} Universal truths or current facts may stay present. \emph{Ex: The teacher said the Earth \textbf{revolves} around the Sun.}
\end{propriete}

\begin{propriete}[Impersonal Passive (Passif impersonnel)]
Used to report beliefs, opinions, or findings without naming the agent.
\begin{itemize}
  \item \textbf{Structure 1:} \emph{It} + \textbf{passive verb} + \textbf{that-clause}.
  \emph{Ex: It \textbf{is said that} women earn less. / It \textbf{was reported that} the vote was fair.}
  \item \textbf{Structure 2:} \emph{Subject} + \textbf{passive verb} + \textbf{to-infinitive}.
  \emph{Ex: Women \textbf{are said to earn} less. / The results \textbf{are expected to be} announced tomorrow.}
  \item Common verbs: \cle{say, think, believe, know, expect, report, allege, estimate}.
\end{itemize}
\end{propriete}

\begin{methode}[Writing a Composition on Volunteering]
\begin{enumerate}
  \item \textbf{Analyse the prompt:} Identify the topic (gender equality), format (article, letter, speech), and target audience.
  \item \textbf{Brainstorm vocabulary:} List 10-15 key words (empowerment, advocacy, bias, inclusive, mentor, campaign).
  \item \textbf{Structure:}
    \begin{itemize}
      \item \textbf{Introduction:} Hook (quote/statistic), define the issue, thesis statement.
      \item \textbf{Body Paragraphs (2-3):} One main idea per paragraph (e.g., Education, Economic access, Political representation). Use linking words (\emph{Furthermore, However, Consequently}). Include \textbf{impersonal passive} and \textbf{reported speech}.
      \item \textbf{Conclusion:} Summarize, call to action, final thought.
    \end{itemize}
  \item \textbf{Draft \& Check:} Verify Sequence of Tenses, Word Order (SVOMPT), Phrasal Verbs usage, Spelling.
\end{enumerate}
\end{methode}

\end{bilan}

\begin{aretenir}[Checklist avant le Bac]
\begin{itemize}
  \item $\square$ I can define and use 10+ words related to \textbf{campaigning and voting} (manifesto, ballot, constituency, poll).
  \item $\square$ I can transform \textbf{Direct Questions} into \textbf{Reported Questions} correctly (Word order SV, Backshift, If/Whether).
  \item $\square$ I understand vocabulary for \textbf{democracy training} (civic education, stakeholder, accountability, consensus).
  \item $\square$ I apply the \textbf{SVOMPT} rule and correct adverb placement in sentences.
  \item $\square$ I can define and use 10+ words related to \textbf{gender equality volunteering} (empowerment, advocacy, bias, glass ceiling).
  \item $\square$ I correctly use \textbf{Phrasal Verbs} (stand for, speak out, take part in, look into) in context.
  \item $\square$ I master the \textbf{Sequence of Tenses} rules after a past reporting verb.
  \item $\square$ I can form the \textbf{Impersonal Passive} both ways (\emph{It is said that...} / \emph{He is said to...}).
  \item $\square$ I can write a structured composition (Intro/Body/Conclusion) on volunteering using target grammar.
  \item $\square$ I can participate in a dialogue simulating a school election campaign or a survey interview.
  \item $\square$ I can listen to a text on democracy training and answer comprehension questions.
  \item $\square$ I can read a text on gender equality and extract specific information / vocabulary.
\end{itemize}
\end{aretenir}

\findecours

\end{document}
